% La identidad, la tolerancia y la convivencia (continuación) : debate y redacción — Espagnol 1ère A — Cours signé OnBuch+
% Programme officiel MINESEC. Compiler avec : tectonic EA12-identidad-tolerancia-convivencia-continuacion.tex
\newcommand{\DOCMATIERE}{Espagnol}
\newcommand{\DOCNIVEAU}{1ère A}
\newcommand{\DOCMODULE}{Módulo 3 : Citoyenneté}
\newcommand{\DOCLECON}{EA12}
\newcommand{\DOCTITRE}{La identidad, la tolerancia y la convivencia (continuación) : debate y redacción}
% OnBuch+ — préambule « cours signés » ENRICHI (compilé avec tectonic / XeLaTeX)
% Système visuel « Pop » d'OnBuch+ : fond crème (jamais blanc), bordures encre
% épaisses, ombres « dures » décalées, titres Archivo Black en capitales.
% Version MATHÉMATIQUES : graphes (pgfplots/tikz), boîtes pédagogiques
% (définition, propriété, méthode, attention, le savais-tu, exercice résolu,
% exercice, TP, bilan), page de garde et signature OnBuch+ ancrée en bas.
%
% Macros attendues dans main.tex AVANT \input{preamble} :
%   \DOCMATIERE  \DOCNIVEAU  \DOCMODULE  \DOCLECON (numéro)  \DOCTITRE
\documentclass[11pt]{article}

\usepackage{fontspec}
\usepackage[spanish,french]{babel} % français = langue principale ; l'espagnol via \esp{..} / espagnol
\usepackage[a4paper,margin=2cm,top=3.4cm,bottom=2.8cm,headheight=1.4cm,footskip=1.3cm]{geometry}
\usepackage{amsmath,amssymb,amsfonts}
\usepackage{mathtools} % \xmapsto, \coloneqq... souvent utilisés par les modèles
\usepackage{cancel} % simplifications barrées (bilans de forces)
% \abs{z} (module, valeur absolue) et \norm{u} (norme), utilisés par les modèles
\providecommand{\abs}{}\let\abs\relax\DeclarePairedDelimiter{\abs}{\lvert}{\rvert}
\providecommand{\norm}{}\let\norm\relax\DeclarePairedDelimiter{\norm}{\lVert}{\rVert}
\usepackage[table]{xcolor} % [table] : fournit aussi \rowcolors (alternance de lignes)
\usepackage{pagecolor}
\usepackage{tikz}
\usetikzlibrary{arrows.meta,positioning,calc,shapes.geometric,shapes.misc,shapes.arrows,decorations.pathmorphing,decorations.pathreplacing,decorations.markings,patterns,shadows,mindmap,trees,fit,backgrounds,matrix,intersections,angles,quotes}
\usepackage{pgfplots}
\usepackage[version=4]{mhchem} % équations nucléaires \ce{^{235}_{92}U}
\usepackage[european,siunitx]{circuitikz} % schémas électriques
\pgfplotsset{compat=1.18}
\usepgfplotslibrary{fillbetween,groupplots} % aires entre courbes (calcul intégral)
\usepackage{enumitem}
\usepackage{fancyhdr}
% [most] charge toutes les bibliothèques tcolorbox (listings, minted, raster,
% documentation...) et gonfle inutilement la mémoire TeX sur un document long
% (~300 boîtes) : on ne charge que ce qui est réellement utilisé.
\usepackage[skins,breakable]{tcolorbox}
\usepackage{graphicx}
\usepackage{colortbl}
\usepackage{booktabs}
\usepackage{tabularx}
\usepackage{diagbox} % case d'en-tête diagonale des tableaux à double entrée
% Colonnes extensibles centrée / gauche / droite (C, L, R), souvent utilisées par les modèles
\newcolumntype{C}{>{\centering\arraybackslash}X}
\newcolumntype{L}{>{\raggedright\arraybackslash}X}
\newcolumntype{R}{>{\raggedleft\arraybackslash}X}
\usepackage{array}
\usepackage{multicol}
\usepackage{siunitx}
% parse-numbers=false : accepte aussi \SI{2,5 \times 10^{-3}}{...}
\sisetup{locale=FR,output-decimal-marker={,},group-minimum-digits=4,parse-numbers=false}
\DeclareSIUnit\ohm{\ensuremath{\symup{\Omega}}} % U+2126 absent de la police
\DeclareSIUnit\division{div} % réglages d'oscilloscope (ms/div, V/div)
\DeclareSIUnit\tr{tr} % tours (vitesse de rotation, ex. \SI{1500}{\tr\per\minute})
\DeclareSIUnit\molar{M} % concentration molaire (ex. \SI{3}{\molar}, \SI{1}{\micro\molar})
\DeclareSIUnit\an{an} % année (le modèle confond parfois avec \year, un compteur TeX) — ex. \SI{5}{\kilo\meter\per\an}
\DeclareSIUnit\FCFA{FCFA} % franc CFA (ex. \SI{500}{\FCFA\per\kilo\gram})
\DeclareSIUnit\degreeBrix{\textdegree Brix} % teneur en sucre d'un sirop/jus (réfractométrie)
\DeclareSIUnit\month{mois} % le modèle utilise parfois \month, un compteur TeX, comme unité de durée
\usepackage{qrcode}
% \degree : commande fréquemment utilisée par les modèles pour un angle ou une
% température en dehors de \SI{}{} (non fournie par les paquets chargés).
\providecommand{\degree}{^{\circ}}
% \qed (fin de démonstration), utilisé par les modèles sans amsthm
\providecommand{\qed}{\hfill$\square$}
% \barre : double barre des valeurs interdites dans un tableau de variations (tkz-tab)
\providecommand{\barre}{\ensuremath{\|}}
% environnement proof (amsthm non chargé) utilisé par les modèles
\ifdefined\proof\else\newenvironment{proof}[1][Démonstration]{\par\noindent\textit{#1.}\ }{\qed\par}\fi
\usepackage{lastpage}
\usepackage{hyperref}
\hypersetup{hidelinks}

% ============================================================
% Palette « Pop » OnBuch+ — mêmes hex que la classe P (plus_ui.dart)
% ============================================================
\definecolor{poporange}{HTML}{F59321}
\definecolor{popdark}{HTML}{E07A0C}
\definecolor{popcream}{HTML}{FFF6E8}
\definecolor{popink}{HTML}{1C1714}
\definecolor{poppurple}{HTML}{7A5AE0}
\definecolor{popgreen}{HTML}{2E9E6A}
\definecolor{poppink}{HTML}{E0457B}
\definecolor{popblue}{HTML}{2F7DE1}
\definecolor{popgold}{HTML}{C98A12}
\definecolor{popmuted}{HTML}{8A7F76}
\definecolor{popline}{HTML}{EADFD2}
\definecolor{popgray}{HTML}{8A7F76}
\colorlet{popgrey}{popgray}
% Alias tolérants (le modèle nomme parfois les couleurs différemment)
\colorlet{popwhite}{white}
\colorlet{popblack}{popink}
\colorlet{popred}{poppink}
\colorlet{popyellow}{popgold}
% Teintes claires (fonds de tableaux/encadrés) — le modèle nomme parfois
% les couleurs avec un suffixe L (ex. popblueL) sans les avoir définies.
\colorlet{poporangeL}{poporange!20}
\colorlet{popdarkL}{popdark!20}
\colorlet{poppurpleL}{poppurple!20}
\colorlet{popgreenL}{popgreen!20}
\colorlet{poppinkL}{poppink!20}
\colorlet{popblueL}{popblue!20}
\colorlet{popgoldL}{popgold!20}
\colorlet{popredL}{popred!20}
\colorlet{popgrayL}{popgray!20}
% Teintes claires pour fonds de tableaux / zones
\colorlet{popblueL}{popblue!10!white}
\colorlet{poporangeL}{poporange!14!white}
\colorlet{popgreenL}{popgreen!12!white}
\colorlet{poppurpleL}{poppurple!12!white}
\colorlet{poppinkL}{poppink!12!white}
\colorlet{popgoldL}{popgold!14!white}

\pagecolor{popcream}

% ============================================================
% Typographie
% ============================================================
\defaultfontfeatures{Ligatures=TeX}
\setmainfont{PlusJakartaSans-Medium.ttf}[Path=../../../fonts/,BoldFont=PlusJakartaSans-Bold.ttf]
% Maths en sans-serif (Fira Math) avec chiffres et lettres droites de Plus
% Jakarta, pour que les formules se fondent dans le texte.
\usepackage[math-style=ISO,bold-style=ISO]{unicode-math}
\setmathfont{FiraMath-Regular.otf}[Path=../../../fonts/]
\setmathfont{PlusJakartaSans-Medium.ttf}[Path=../../../fonts/,range={up/num,up/{Latin,latin},\mathup}]
\usepackage{icomma} % virgule décimale sans espace en mode math
% Caractères mathématiques Unicode absents des polices de texte (Plus Jakarta,
% Archivo Black) : rendus par leur équivalent mathématique, en texte comme en
% formule — sinon ils s'affichent en carré vide (ex. « Arithmétique dans ℤ »).
\usepackage{newunicodechar}
\newunicodechar{ℝ}{\ensuremath{\mathbb{R}}}
\newunicodechar{ℤ}{\ensuremath{\mathbb{Z}}}
\newunicodechar{ℂ}{\ensuremath{\mathbb{C}}}
\newunicodechar{ℕ}{\ensuremath{\mathbb{N}}}
\newunicodechar{ℚ}{\ensuremath{\mathbb{Q}}}
\newunicodechar{𝔻}{\ensuremath{\mathbb{D}}}
\newunicodechar{α}{\ensuremath{\alpha}}
\newunicodechar{β}{\ensuremath{\beta}}
\newunicodechar{γ}{\ensuremath{\gamma}}
\newunicodechar{δ}{\ensuremath{\delta}}
\newunicodechar{ε}{\ensuremath{\varepsilon}}
\newunicodechar{θ}{\ensuremath{\theta}}
\newunicodechar{λ}{\ensuremath{\lambda}}
\newunicodechar{μ}{\ensuremath{\mu}}
\newunicodechar{σ}{\ensuremath{\sigma}}
\newunicodechar{τ}{\ensuremath{\tau}}
\newunicodechar{φ}{\ensuremath{\varphi}}
\newunicodechar{ψ}{\ensuremath{\psi}}
\newunicodechar{ω}{\ensuremath{\omega}}
\newunicodechar{Σ}{\ensuremath{\Sigma}}
% majuscules produites par \MakeUppercase dans les titres (α-aminés -> Α)
\newunicodechar{Α}{\ensuremath{\alpha}}
\newunicodechar{Β}{\ensuremath{\beta}}
\newunicodechar{Γ}{\ensuremath{\Gamma}}
\newunicodechar{Δ}{\ensuremath{\Delta}}
\newunicodechar{Ε}{\ensuremath{\varepsilon}}
\newunicodechar{Θ}{\ensuremath{\Theta}}
\newunicodechar{Λ}{\ensuremath{\Lambda}}
\newunicodechar{Μ}{\ensuremath{\mu}}
\newunicodechar{Τ}{\ensuremath{\tau}}
\newunicodechar{Φ}{\ensuremath{\Phi}}
\newunicodechar{Ψ}{\ensuremath{\Psi}}
\newunicodechar{Ω}{\ensuremath{\Omega}}
\newunicodechar{↦}{\ensuremath{\mapsto}}
\newunicodechar{∈}{\ensuremath{\in}}
\newunicodechar{∉}{\ensuremath{\notin}}
\newunicodechar{∧}{\ensuremath{\wedge}}
\newunicodechar{⇔}{\ensuremath{\Leftrightarrow}}
\newunicodechar{⟺}{\ensuremath{\Longleftrightarrow}}
\newunicodechar{⇒}{\ensuremath{\Rightarrow}}
\newunicodechar{⟹}{\ensuremath{\Longrightarrow}}
\newunicodechar{∘}{\ensuremath{\circ}}
\newunicodechar{∪}{\ensuremath{\cup}}
\newunicodechar{∩}{\ensuremath{\cap}}
\newunicodechar{≡}{\ensuremath{\equiv}}
\newunicodechar{⊂}{\ensuremath{\subset}}
\newunicodechar{⊥}{\ensuremath{\perp}}
\newunicodechar{∃}{\ensuremath{\exists}}
\newunicodechar{∀}{\ensuremath{\forall}}
\newunicodechar{∛}{\ensuremath{\sqrt[3]{\,}}}
\newunicodechar{∜}{\ensuremath{\sqrt[4]{\,}}}
\newunicodechar{⃗}{\ensuremath{{}^{\to}}}
\newunicodechar{ⁿ}{\ifmmode^{n}\else\textsuperscript{\ensuremath{n}}\fi}
\newunicodechar{ˣ}{\ifmmode^{x}\else\textsuperscript{\ensuremath{x}}\fi}
\newunicodechar{ᵇ}{\ifmmode^{b}\else\textsuperscript{\ensuremath{b}}\fi}
\newunicodechar{ʳ}{\ifmmode^{r}\else\textsuperscript{\ensuremath{r}}\fi}
\newunicodechar{ᵘ}{\ifmmode^{u}\else\textsuperscript{\ensuremath{u}}\fi}
\newunicodechar{ʸ}{\ifmmode^{y}\else\textsuperscript{\ensuremath{y}}\fi}
\newunicodechar{ᵃ}{\ifmmode^{a}\else\textsuperscript{\ensuremath{a}}\fi}
\newunicodechar{ᵉ}{\ifmmode^{e}\else\textsuperscript{\ensuremath{e}}\fi}
\newunicodechar{ᵏ}{\ifmmode^{k}\else\textsuperscript{\ensuremath{k}}\fi}
\newunicodechar{ᵖ}{\ifmmode^{p}\else\textsuperscript{\ensuremath{p}}\fi}
\newunicodechar{ᴺ}{\ifmmode^{N}\else\textsuperscript{\ensuremath{N}}\fi}
\newunicodechar{ᶜ}{\ifmmode^{c}\else\textsuperscript{\ensuremath{c}}\fi}
\newunicodechar{ʰ}{\ifmmode^{h}\else\textsuperscript{\ensuremath{h}}\fi}
\newunicodechar{ᵠ}{\ifmmode^{q}\else\textsuperscript{\ensuremath{q}}\fi}
\newunicodechar{ₙ}{\ifmmode_{n}\else\textsubscript{\ensuremath{n}}\fi}
\newunicodechar{ₐ}{\ifmmode_{a}\else\textsubscript{\ensuremath{a}}\fi}
\newunicodechar{ᵢ}{\ifmmode_{i}\else\textsubscript{\ensuremath{i}}\fi}
\newunicodechar{ᵣ}{\ifmmode_{r}\else\textsubscript{\ensuremath{r}}\fi}
\newunicodechar{ₖ}{\ifmmode_{k}\else\textsubscript{\ensuremath{k}}\fi}
\newunicodechar{ᵦ}{\ifmmode_{\beta}\else\textsubscript{\ensuremath{\beta}}\fi}
\newunicodechar{ₓ}{\ifmmode_{x}\else\textsubscript{\ensuremath{x}}\fi}
\newfontfamily\popdisplay{ArchivoBlack-Regular.ttf}[Path=../../../fonts/]
\newfontfamily\poplogo{SpaceGrotesk-Bold.ttf}[Path=../../../fonts/]
\newfontfamily\popsemi{PlusJakartaSans-SemiBold.ttf}[Path=../../../fonts/]
\newfontfamily\popxbold{PlusJakartaSans-ExtraBold.ttf}[Path=../../../fonts/]
\newfontfamily\popxboldls{PlusJakartaSans-ExtraBold.ttf}[Path=../../../fonts/,LetterSpace=3.0]

\setlist[enumerate]{leftmargin=1.7em,itemsep=5pt,topsep=4pt}
\setlist[itemize]{leftmargin=1.7em,itemsep=4pt,topsep=4pt,label={\color{poporange}\textbullet}}
\setlength{\parindent}{0pt}
\setlength{\parskip}{5pt}
\renewcommand{\arraystretch}{1.25}

% pgfplots : style Pop par défaut
\pgfplotsset{
  every axis/.append style={axis line style={popink,line width=1pt},tick style={popink},
    grid style={popline,line width=0.5pt},label style={font=\small},tick label style={font=\footnotesize}},
  popcurve/.style={poporange,line width=1.8pt,no markers,smooth},
}
% Même style disponible dans un \draw TikZ simple (hors pgfplots)
\tikzset{popcurve/.style={poporange,line width=1.8pt}}
\tikzset{popgrid/.style={popline,very thin}}
\colorlet{popcurve}{poporange} % le nom du style est parfois utilisé comme couleur

% ============================================================
% Pill / squiggle
% ============================================================
\newcommand{\poppill}[3]{%
  \tikz[baseline=-0.55ex]\node[draw=popink,line width=1.2pt,rounded corners=2.6pt,%
    fill=#2,inner xsep=6.5pt,inner ysep=3pt]{%
    \popxboldls\fontsize{7}{7}\selectfont\color{#3}\MakeUppercase{#1}};%
}
\newcommand{\popsquiggle}[1]{%
  \tikz[baseline=-0.4ex]{
    \draw[#1,line width=2.6pt,line cap=round]
      plot [smooth] coordinates {(0,0.09) (0.34,0.01) (0.68,0.21) (1.02,0.01) (1.36,0.10)};
  }%
}
% Mot-clé surligné (marqueur orange sous le texte)
\newcommand{\cle}[1]{{\popxbold\color{popdark}#1}}

% ============================================================
% En-tête / pied
% ============================================================
\pagestyle{fancy}
\fancyhf{}
\renewcommand{\headrulewidth}{0pt}
\renewcommand{\footrulewidth}{0pt}
\lhead{\raisebox{-0.15ex}{\poplogo\fontsize{13}{13}\selectfont\color{popink}OnBuch\color{poporange}+}}
\rhead{\poppill{\DOCMATIERE\ \textbullet\ \DOCNIVEAU\ \textbullet\ Leçon \DOCLECON}{white}{popink}}
\lfoot{\popxbold\fontsize{7.6}{7.6}\selectfont\color{popmuted}\MakeUppercase{Cours signé OnBuch+}\ \textbullet\ {\popsemi\DOCTITRE}}
\rfoot{\tikz[baseline=-0.6ex]\node[rounded corners=8.5pt,fill=popink,minimum height=17pt,minimum width=17pt,inner xsep=5pt,inner sep=0]{\popxbold\fontsize{8}{8}\selectfont\color{popcream}\thepage\,/\,\pageref*{LastPage}};}

% ============================================================
% Titres
% ============================================================
\newcommand{\chaptitre}[2]{%
  \begin{tcolorbox}[enhanced,colback=poporange,colframe=popink,boxrule=2.6pt,arc=11pt,
      drop shadow=popink,left=15pt,right=15pt,top=12pt,bottom=11pt,before skip=3pt,after skip=18pt]
    \poppill{#2}{popink}{popcream}\par\vspace{9pt}
    {\raggedright\hyphenpenalty=10000\exhyphenpenalty=10000\popdisplay\fontsize{22}{24}\selectfont\color{popcream}\MakeUppercase{#1}\par}
  \end{tcolorbox}
}
% Section numérotée : pastille orange avec le numéro + titre Archivo
\newcounter{popsec}
\newcommand{\coursec}[1]{%
  \stepcounter{popsec}\setcounter{popsub}{0}%
  \par\vspace{16pt}\noindent
  \begin{minipage}{\linewidth}
  \tikz[baseline=-0.7ex]\node[draw=popink,line width=1.6pt,fill=poporange,rounded corners=4pt,minimum size=22pt,inner sep=0]{\popdisplay\fontsize{12}{12}\selectfont\color{popcream}\thepopsec};%
  \hspace{8pt}{\popdisplay\fontsize{15}{17}\selectfont\color{popink}\MakeUppercase{#1}}\par
  \vspace{4pt}\noindent{\color{poporange}\rule{36pt}{3.6pt}}
  \end{minipage}\par\nopagebreak\vspace{8pt}
}
\newcounter{popsub}
\newcommand{\courssub}[1]{%
  \stepcounter{popsub}%
  \par\vspace{9pt}\noindent{\popxbold\fontsize{11.5}{13}\selectfont\color{popink}{\color{poporange}\thepopsec.\thepopsub}\hspace{6pt}#1}\par\nopagebreak\vspace{4pt}
}

% ============================================================
% Boîtes pédagogiques — même recette : fond blanc, bordure épaisse colorée,
% ombre dure encre, badge plein. Titre optionnel [#1].
% ============================================================
\tcbset{popbase/.style={breakable,enhanced,colback=white,boxrule=2.3pt,arc=9pt,
  shadow={3pt}{-3pt}{0pt}{popink},left=14pt,right=14pt,top=11pt,bottom=11pt,before skip=11pt,after skip=15pt,
  fontupper=\popsemi\fontsize{10.6}{14.5}\selectfont\color{popink},
  fontlower=\popsemi\fontsize{10.6}{14.5}\selectfont\color{popink}}}
% Coche dessinée (le glyphe ✓ n'existe pas dans les polices du document)
\newcommand{\popcheck}[1]{\tikz[baseline=-0.5ex]\draw[#1,line width=1.6pt,line cap=round,line join=round](-0.13,0.0)--(-0.04,-0.09)--(0.13,0.1);}
\providecommand{\coche}{\popcheck{popgreen}}
\providecommand{\resultat}[1]{\par\smallskip\noindent\fcolorbox{popgreen}{popgreenL}{\parbox{\dimexpr\linewidth-2\fboxsep-2\fboxrule\relax}{\textbf{Résultat.} #1}}\par\smallskip}
\newcommand{\popboxhead}[3]{\poppill{#1}{#2}{white}\ifx\relax#3\relax\else\hspace{6pt}{\popxbold\fontsize{10.6}{12}\selectfont\color{popink}#3}\fi\par\vspace{7pt}}

\newtcolorbox{aretenir}[1][]{popbase,colframe=popgreen,before upper={\popboxhead{À retenir}{popgreen}{#1}}}
% Texte en espagnol (document de lecture, dialogue, poème, extrait) : cadre
% d'étude. Un titre optionnel (source, auteur) s'affiche dans l'en-tête.
\newtcolorbox{textebox}[1][]{popbase,colframe=popblue,before upper={\popboxhead{Texto}{popblue}{#1}\begin{otherlanguage}{spanish}},after upper={\end{otherlanguage}}}
% Marques ✓ / ✗ (forme correcte / fautive) souvent utilisées par les modèles
\providecommand{\cross}{\textcolor{poppink}{\ensuremath{\times}}}
\providecommand{\xmark}{\cross}\providecommand{\crossmark}{\cross}\providecommand{\cmark}{\ensuremath{\checkmark}}
% Ligne de réponse / justification à compléter (inventée par les modèles)
\providecommand{\justification}[1]{\par\noindent\textit{Justification~:} \dotfill\par}
% \textipa (package tipa non chargé) : la police ne couvre pas l'API -> simple passage
\providecommand{\textipa}[1]{#1}
% Soulignement (ulem non chargé) : \uline, \ul
\providecommand{\uline}[1]{\underline{#1}}\providecommand{\ul}[1]{\underline{#1}}
% Ordinaux français écrits en macro par les modèles : 1\ière, 3\ième
\newcommand{\ière}{\textsuperscript{re}}\newcommand{\ième}{\textsuperscript{e}}
% Mot ou phrase en espagnol à l'intérieur d'un paragraphe français
\newcommand{\esp}[1]{\ifmmode\text{\textit{#1}}\else\begingroup\selectlanguage{spanish}\itshape #1\endgroup\fi}
\newtcolorbox{exemplebox}[1][]{popbase,colframe=poporange,before upper={\popboxhead{Exemple}{poporange}{#1}}}
\newtcolorbox{definition}[1][]{popbase,colframe=popblue,before upper={\popboxhead{Définition}{popblue}{#1}}}
\newtcolorbox{propriete}[1][]{popbase,colframe=poppurple,before upper={\popboxhead{Propriété}{poppurple}{#1}}}
\newtcolorbox{methode}[1][]{popbase,colframe=popgold,before upper={\popboxhead{Méthode}{popgold}{#1}}}
\newtcolorbox{attention}[1][]{popbase,colframe=poppink,before upper={\popboxhead{Attention}{poppink}{#1}}}
\newtcolorbox{experience}[1][]{popbase,colframe=popblue,colback=popblueL,before upper={\popboxhead{Expérience}{popblue}{#1}}}
% « Le savais-tu ? » : carte violette pleine (contexte, culture, Cameroun)
\newtcolorbox{savaistu}[1][]{popbase,colback=poppurple,colframe=popink,
  fontupper=\popsemi\fontsize{10.4}{14.2}\selectfont\color{white},
  before upper={\poppill{Le savais-tu\,?}{white}{poppurple}\ifx\relax#1\relax\else\hspace{6pt}{\popxbold\color{white}#1}\fi\par\vspace{7pt}}}
% Exercice résolu : énoncé en haut, \tcblower puis la solution sur fond crème
\newcounter{popexo}\newcounter{popexr}
\newtcolorbox{exoresolu}[1][]{popbase,colframe=poporange,colbacklower=popcream,
  segmentation style={popink,line width=1.2pt,dashed},
  before upper={\stepcounter{popexr}\popboxhead{Exercice résolu \thepopexr}{poporange}{#1}},
  before lower={\poppill{Solution}{popink}{popcream}\par\vspace{6pt}}}
% Exercice (non résolu en place) — la correction est dans la partie Corrigés
\newtcolorbox{exercice}[1][]{popbase,colframe=popink,
  before upper={\stepcounter{popexo}\popboxhead{Exercice \thepopexo}{popink}{#1}}}
% Corrigé
\newtcolorbox{corrige}[1][]{popbase,colframe=popgreen,colback=popgreenL,
  before upper={\popboxhead{Corrigé}{popgreen}{#1}}}
% Objectifs / prérequis / bilan
\newtcolorbox{objectifs}{popbase,colframe=popink,colback=poporangeL,
  before upper={\popboxhead{Objectifs de la leçon}{popink}{}}}
\newtcolorbox{prerequis}{popbase,colframe=popink,colback=popblueL,
  before upper={\popboxhead{Prérequis}{popblue}{}}}
\newtcolorbox{bilan}{popbase,colframe=popink,colback=white,boxrule=2.8pt,
  before upper={\popboxhead{Fiche bilan}{poporange}{L'essentiel à connaître pour l'examen}}}
% Figure encadrée avec légende
\newtcolorbox{popfigure}[1][]{enhanced,breakable=false,colback=white,colframe=popline,boxrule=1.4pt,arc=8pt,
  left=10pt,right=10pt,top=10pt,bottom=8pt,before skip=10pt,after skip=12pt,halign=center,
  lower separated=false,halign lower=center,
  fontlower=\popsemi\fontsize{9}{11.5}\selectfont\color{popmuted},
  #1}
\newcommand{\legende}[1]{\par\vspace{6pt}{\popsemi\fontsize{9}{11.5}\selectfont\color{popmuted}#1}}

% ============================================================
% Signature OnBuch+
% ============================================================
\newcommand{\poplogoinline}[1]{{\poplogo\fontsize{#1}{#1}\selectfont\color{popink}OnBuch\color{poporange}+}}
\newcommand{\signatureonbuch}{%
  \begin{tcolorbox}[enhanced,colback=popink,colframe=popink,boxrule=2.2pt,arc=10pt,
      drop shadow=poporange,left=14pt,right=14pt,top=10pt,bottom=10pt,before skip=0pt,after skip=0pt]
    \tikz[baseline=-0.7ex]\node[circle,fill=popgreen,draw=popcream,line width=1.3pt,minimum size=22pt,inner sep=0]{\popcheck{white}};%
    \hspace{9pt}\begin{minipage}[c]{0.62\linewidth}
      {\popxbold\fontsize{12}{13}\selectfont\color{popcream}Signé\ \textbullet\ {\poplogo OnBuch\color{poporange}+}}\par\vspace{2pt}
      {\raggedright\popsemi\fontsize{8}{10.5}\selectfont\color{popline}\MakeUppercase{Équipe pédagogique OnBuch+}\\\MakeUppercase{Conforme au programme officiel MINESEC}\par}
    \end{minipage}\hfill
    {\poplogo\fontsize{22}{22}\selectfont\color{popcream}OnBuch\color{poporange}+}
  \end{tcolorbox}%
}

% ============================================================
% Page de garde
% #1 titre · #2 matière · #3 niveau · #4 module · #5 n° leçon · #6 durée ·
% #7 description / objectifs (texte ou itemize)
% ============================================================
\newcommand{\pagegarde}[7]{%
  \thispagestyle{empty}
  \newgeometry{a4paper,margin=1.7cm,top=1.6cm,bottom=1.4cm}%
  \begin{tikzpicture}[remember picture,overlay]
    \fill[poporange!18] ([xshift=-3.2cm,yshift=-3.6cm]current page.north east) circle (4.6cm);
    \fill[poppurple!14] ([xshift=2.4cm,yshift=7.6cm]current page.south west) circle (3.8cm);
  \end{tikzpicture}%
  {\poplogo\fontsize{30}{30}\selectfont\color{popink}OnBuch\color{poporange}+}\hfill
  \raisebox{0.6ex}{\poppill{Cours signé \textbullet\ Premium}{popink}{popcream}}\par
  \vspace{1pt}\popsquiggle{poppurple}\par
  \vspace{8pt}
  \begin{tcolorbox}[enhanced,colback=poporange,colframe=popink,boxrule=2.8pt,arc=13pt,
      drop shadow=popink,left=18pt,right=18pt,top=12pt,bottom=13pt,after skip=10pt]
    \poppill{#2 \textbullet\ #3}{popink}{popcream}\hspace{5pt}\poppill{#4}{white}{popink}\par\vspace{8pt}
    {\popdisplay\fontsize{14}{14}\selectfont\color{popink}LEÇON #5}\par\vspace{4pt}
    {\raggedright\hyphenpenalty=10000\exhyphenpenalty=10000\popdisplay\fontsize{25}{27}\selectfont\color{popcream}\MakeUppercase{#1}\par}
    \vspace{8pt}
    {\popxbold\fontsize{9.5}{9.5}\selectfont\color{popink}\MakeUppercase{Durée conseillée : #6}}
  \end{tcolorbox}
  \begin{tcolorbox}[enhanced,colback=white,colframe=popink,boxrule=2.2pt,arc=10pt,
      drop shadow=popink,left=16pt,right=16pt,top=10pt,bottom=10pt,after skip=8pt]
    \poppill{Dans cette leçon}{poppurple}{white}\par\vspace{5pt}
    \popsemi\fontsize{9.3}{12.6}\selectfont\color{popink}#7
  \end{tcolorbox}
  \noindent\begin{minipage}[t]{0.487\linewidth}
    \begin{tcolorbox}[enhanced,colback=popgreen,colframe=popink,boxrule=2.2pt,arc=10pt,
        drop shadow=popink,left=11pt,right=11pt,top=10pt,bottom=10pt,height=3.1cm,valign=center]
      \tcbox[colback=white,colframe=popink,boxrule=1.3pt,arc=5pt,boxsep=0pt,nobeforeafter,
        left=3pt,right=3pt,top=3pt,bottom=3pt,valign=center]{\qrcode[height=1.9cm]{https://play.google.com/store/apps/details?id=com.luvvix.onbuch&hl=fr}}%
      \hspace{8pt}\begin{minipage}[c]{0.5\linewidth}
        {\popdisplay\fontsize{11}{12.5}\selectfont\color{white}\MakeUppercase{Télécharge OnBuch}}\par\vspace{4pt}
        {\raggedright\popsemi\fontsize{8.4}{11}\selectfont\color{white}Gratuit \textbullet\ Google Play\\Tuteur IA, annales, concours, résultats\par}
      \end{minipage}
    \end{tcolorbox}
  \end{minipage}\hfill
  \begin{minipage}[t]{0.487\linewidth}
    \begin{tcolorbox}[enhanced,colback=poppurple,colframe=popink,boxrule=2.2pt,arc=10pt,
        drop shadow=popink,left=11pt,right=11pt,top=10pt,bottom=10pt,height=3.1cm,valign=center]
      \tikz[baseline=-0.7ex]\node[circle,fill=white,draw=popink,line width=1.3pt,minimum size=1.6cm,inner sep=0]{\popdisplay\fontsize{24}{24}\selectfont\color{poppurple}+};%
      \hspace{8pt}\begin{minipage}[c]{0.6\linewidth}
        {\popdisplay\fontsize{11}{12.5}\selectfont\color{white}\MakeUppercase{Passe à OnBuch+}}\par\vspace{4pt}
        {\raggedright\popsemi\fontsize{8.4}{11}\selectfont\color{white}Tous les cours signés, Tuteur IA illimité, fiches PDF — onglet OnBuch+ de l'app\par}
      \end{minipage}
    \end{tcolorbox}
  \end{minipage}
  \vspace{8pt}
  \popcontenu
  \vfill
  \signatureonbuch
  \restoregeometry
  \newpage
}

% Rangée « Ce cours contient » (page de garde)
\newcommand{\poptile}[3]{% #1 couleur #2 gros texte #3 légende
  \begin{tcolorbox}[enhanced,colback=#1,colframe=popink,boxrule=2pt,arc=9pt,shadow={2.5pt}{-2.5pt}{0pt}{popink},
      left=6pt,right=6pt,top=9pt,bottom=9pt,halign=center,valign=center,height=1.75cm,width=\linewidth,nobeforeafter]
    {\popdisplay\fontsize{12.5}{13}\selectfont\color{white}\MakeUppercase{#2}}\par\vspace{3pt}
    {\centering\popsemi\fontsize{7.8}{9.5}\selectfont\color{white}#3\par}
  \end{tcolorbox}}
\newcommand{\popcontenu}{%
  \par\noindent{\popxboldls\fontsize{7.5}{7.5}\selectfont\color{popmuted}CE COURS SIGNÉ CONTIENT}\par\vspace{7pt}
  \noindent\begin{minipage}[t]{0.235\linewidth}\poptile{popblue}{Cours}{complet, illustré, pas à pas}\end{minipage}\hfill
  \begin{minipage}[t]{0.235\linewidth}\poptile{poporange}{Méthodes}{et exercices résolus}\end{minipage}\hfill
  \begin{minipage}[t]{0.235\linewidth}\poptile{poppink}{Exercices}{type examen + corrigés}\end{minipage}\hfill
  \begin{minipage}[t]{0.235\linewidth}\poptile{popgreen}{Fiche bilan}{l'essentiel à retenir}\end{minipage}\par}

% Sommaire
\newtcolorbox{sommaire}{enhanced,colback=white,colframe=popink,boxrule=2.2pt,arc=10pt,
  drop shadow=popink,left=18pt,right=18pt,top=13pt,bottom=13pt,before skip=6pt,after skip=20pt,
  before upper={\poppill{Sommaire}{poppurple}{white}\par\vspace{9pt}\popsemi\fontsize{10.6}{14.5}\selectfont\color{popink}}}

% Signature de fin de cours — poussée tout en bas de la dernière page
\newcommand{\findecours}{%
  \par\vfill\nopagebreak
  \begin{center}\popsquiggle{poporange}\end{center}\vspace{4pt}
  \signatureonbuch
}
\AtBeginDocument{\def\llbracket{\mathopen{[\mkern-3.5mu[}}\def\rrbracket{\mathclose{]\mkern-3.5mu]}}} % intervalles d'entiers (glyphe ⟦ absent de la police)
% Glyphes absents de Fira Math : redéfinis à partir de glyphes disponibles
\AtBeginDocument{%
  \def\setminus{\mathbin{\backslash}}\let\smallsetminus\setminus
  \def\vdots{\mathord{\vbox{\baselineskip3.2pt\lineskiplimit0pt\kern4pt\hbox{.}\hbox{.}\hbox{.}}}}%
  \def\ddots{\mathinner{\mkern1mu\raise6.5pt\hbox{.}\mkern2mu\raise3.6pt\hbox{.}\mkern2mu\raise0.7pt\hbox{.}\mkern1mu}}%
  \def\triangle{\mathord{\scalebox{0.9}{$\symup{\Delta}$}}}%
  \def\nabla{\mathord{\raisebox{\depth}{\scalebox{1}[-1]{$\symup{\Delta}$}}}}%
  \def\checkmark{\popcheck{popdark}}%
  \def\star{\mathbin{\ast}}\def\bigstar{\mathord{\ast}}%
  \def\diamond{\mathbin{\scalebox{0.6}{$\Diamond$}}}%
  \def\bigcup{\mathop{\mathchoice{\raisebox{-0.3em}{\scalebox{1.7}{$\cup$}}}{\scalebox{1.25}{$\cup$}}{\cup}{\cup}}\displaylimits}%
  \def\bigcap{\mathop{\mathchoice{\raisebox{-0.3em}{\scalebox{1.7}{$\cap$}}}{\scalebox{1.25}{$\cap$}}{\cap}{\cap}}\displaylimits}%
  \def\lesseqqgtr{\mathrel{\lessgtr}}\def\bigoplus{\mathop{\oplus}\displaylimits}%
}
\providecommand{\ssi}{si et seulement si} % abréviation fréquente
\AtBeginDocument{\providecommand{\wideparen}{\overparen}} % arc de cercle

\begin{document}

\pagegarde{La identidad, la tolerancia y la convivencia (continuación) : debate y redacción}{Espagnol}{1ère A}{Módulo 3 : Citoyenneté}{EA12}{2 h}{Cette leçon mixte poursuit l'étude de l'identité et du vivre-ensemble en développant deux compétences clés : participer à un débat oral et rédiger un court texte argumentatif en espagnol. Elle consolide les connecteurs logiques et la traduction de phrases sur la tolérance et la discrimination.
\vspace{4pt}
\begin{multicols}{2}\small
\begin{itemize}
\item À la fin de cette leçon, je sais comprendre un dialogue-débat sur la tolérance et la discrimination
\item À la fin de cette leçon, je sais utiliser le lexique du débat, de l'argumentation et du vivre-ensemble
\item À la fin de cette leçon, je sais exprimer mon opinion, mon accord et mon désaccord en espagnol
\item À la fin de cette leçon, je sais structurer un texte argumentatif avec en primer lugar, además, sin embargo, por último
\item À la fin de cette leçon, je sais rédiger un court texte argumenté d'environ 100 mots sur le vivre-ensemble
\item À la fin de cette leçon, je sais traduire des phrases du français vers l'espagnol sur l'identité et la tolérance (thème)
\end{itemize}\end{multicols}}

\begin{sommaire}
\begin{multicols}{2}
\begin{enumerate}
\item Texte support : un débat au lycée
\item Lexique thématique : débat et vivre-ensemble
\item Grammaire : les connecteurs logiques de l'argumentation
\item Traduction : phrases sur l'identité et la tolérance
\item Méthode du débat oral
\item Production écrite guidée : texte argumenté (100 mots)
\item Activité d'intégration
\item Méthodes et exercices résolus
\item Exercices
\item Corrigés des exercices
\item Fiche bilan
\end{enumerate}
\end{multicols}
\end{sommaire}

\begin{objectifs}
\begin{itemize}
\item À la fin de cette leçon, je sais comprendre un dialogue-débat sur la tolérance et la discrimination
\item À la fin de cette leçon, je sais utiliser le lexique du débat, de l'argumentation et du vivre-ensemble
\item À la fin de cette leçon, je sais exprimer mon opinion, mon accord et mon désaccord en espagnol
\item À la fin de cette leçon, je sais structurer un texte argumentatif avec en primer lugar, además, sin embargo, por último
\item À la fin de cette leçon, je sais rédiger un court texte argumenté d'environ 100 mots sur le vivre-ensemble
\item À la fin de cette leçon, je sais traduire des phrases du français vers l'espagnol sur l'identité et la tolérance (thème)
\item À la fin de cette leçon, je sais traduire des phrases de l'espagnol vers le français sur ces mêmes thèmes (version)
\item À la fin de cette leçon, je sais participer à un débat réglé sur la convivencia en respectant la parole d'autrui
\end{itemize}
\end{objectifs}

\begin{prerequis}
\begin{itemize}
\item Le présent de l'indicatif des verbes réguliers et irréguliers fréquents (Seconde)
\item Le vocabulaire de l'identité et de la tolérance vu à la leçon EA11
\item L'expression de l'opinion simple : pienso que, creo que, me parece que (Seconde)
\item Les pronoms personnels et la construction de la phrase simple en espagnol
\item La notion de connecteurs temporels simples (primero, luego, después)
\end{itemize}
\end{prerequis}

\newpage

\coursec{Texte support : un débat au lycée}

Au Cameroun, plus de 250 langues coexistent et pourtant les élèves d'un même lycée partagent la même cour de récréation. À Yaoundé, à Douala ou à Bafoussam, un élève bamiléké, une élève beti et un élève du Nord se retrouvent chaque matin dans la même salle de classe, mangent au même réfectoire et jouent au football dans la même équipe. En Espagne ou au Mexique, les débats sur l'immigration et la diversité culturelle occupent quotidiennement les médias : comment accueillir les nouveaux arrivants ? Comment respecter les différences ? Comment vivre ensemble ?

Comment exprimer son opinion, nuancer, réfuter et conclure dans un débat en espagnol ? Comment rédiger un texte argumenté sur le vivre-ensemble ? C'est tout l'enjeu de cette leçon : apprendre à débattre et à écrire sur \esp{la identidad, la tolerancia y la convivencia}.

Mais avant de débattre nous-mêmes, observons un modèle : le dialogue entre deux élèves, Lucía et Etienne, qui discutent de la diversité culturelle dans leur lycée de Yaoundé.

\courssub{Le texte : un débat entre Lucía et Etienne}

Lis d'abord le dialogue en silence, puis le professeur le lira à voix haute. Pendant l'écoute, ton objectif est simple : repérer \cle{qui parle}, \cle{de quoi on parle} et \cle{quelle est la position de chaque personnage}.

\begin{textebox}[Un debate en el liceo de Yaoundé]
— Lucía : Hola Etienne, ¿has leído el artículo del periódico escolar sobre la diversidad cultural en nuestro liceo ?

— Etienne : Sí, lo he leído. Es un tema muy interesante, pero también complicado.

— Lucía : En primer lugar, creo que la diversidad cultural es una riqueza. Además, en nuestro liceo conviven alumnos de muchas regiones : del Norte, del Oeste, del Litoral... Cada uno trae sus costumbres, su música y su lengua. Para mí, eso es fantástico.

— Etienne : Sin embargo, a veces hay conflictos por los prejuicios. Ayer, en el patio, dos chicos se pelearon porque uno se burló del acento del otro. En mi opinión, la convivencia no es tan fácil como dices.

— Lucía : Es verdad, Etienne, no digo que sea fácil. Pero pienso que los prejuicios nacen de la ignorancia. Si compartimos nuestras raíces, si hablamos de nuestras tradiciones, aprendemos a respetarnos.

— Etienne : Tienes razón en parte. Pero también hay discriminación : algunos alumnos no quieren sentarse con compañeros de otras regiones.

— Lucía : Precisamente por eso hay que hablar del problema. La tolerancia no significa ignorar las diferencias, sino aceptarlas. Mira nuestro equipo de fútbol : somos de diez regiones diferentes y ganamos juntos.

— Etienne : Bueno, eso es cierto. El deporte une a la gente.

— Lucía : Por último, pienso que el diálogo y el respeto son la clave de la convivencia. Si nos escuchamos, el liceo será un lugar mejor para todos.

— Etienne : De acuerdo, Lucía. Al final, compartimos la misma escuela, el mismo patio y el mismo futuro.
\end{textebox}

\begin{definition}[Glossaire du texte]
\begin{itemize}
\item \esp{la convivencia} : le vivre-ensemble, la vie en communauté
\item \esp{la tolerancia} : la tolérance
\item \esp{la discriminación} : la discrimination
\item \esp{respetar} : respecter (verbe régulier en \esp{-ar})
\item \esp{la costumbre} : la coutume, l'habitude
\item \esp{el prejuicio} : le préjugé
\item \esp{compartir} : partager
\item \esp{la raíz} (pluriel \esp{las raíces}) : la racine ; au figuré, les origines
\item \esp{burlarse de alguien} : se moquer de quelqu'un
\item \esp{la clave} : la clé, l'élément essentiel
\item \esp{pelearse} : se disputer, se battre
\item \esp{el patio} : la cour de récréation
\item \esp{el periódico escolar} : le journal scolaire
\item \esp{hay que} + infinitif : il faut (tournure impersonnelle)
\end{itemize}
\end{definition}

\begin{savaistu}[Le 16 novembre : Día Internacional de la Tolerancia]
En Espagne, comme dans de nombreux pays hispanophones, le 16 novembre est le \esp{Día Internacional de la Tolerancia}, proclamé par l'UNESCO en 1995. Ce jour-là, les écoles organisent des débats, des expositions et des ateliers sur le respect des différences. Au Cameroun, pays de plus de 250 langues, et en Guinée équatoriale, seul pays hispanophone d'Afrique subsaharienne, cette journée a une résonance toute particulière : la diversité n'y est pas une théorie, c'est la vie quotidienne.
\end{savaistu}

\courssub{Compréhension globale : qui parle ? de quoi ?}

Après une première lecture, réponds mentalement à ces trois questions fondamentales :

\begin{enumerate}
\item \textbf{¿Quiénes hablan ?} Qui sont les locuteurs ? Où se trouvent-ils ?
\item \textbf{¿De qué hablan ?} Quel est le sujet du débat ?
\item \textbf{¿Qué posiciones defienden ?} Quelle est la position de Lucía ? Et celle d'Etienne ?
\end{enumerate}

\begin{corrige}[Compréhension globale]
\begin{enumerate}
\item Deux élèves du lycée de Yaoundé dialoguent : \textbf{Lucía} et \textbf{Etienne}. Ils parlent sans doute pendant la récréation ou après les cours, à propos d'un article du journal scolaire.
\item Ils parlent de \textbf{la diversidad cultural} dans leur lycée et, plus largement, du \textbf{vivre-ensemble} (\esp{la convivencia}) entre élèves de régions différentes.
\item \textbf{Lucía} défend une position optimiste : la diversité est une richesse (\esp{una riqueza}) et le dialogue est la solution. \textbf{Etienne} est plus nuancé : il reconnaît l'intérêt du sujet mais souligne les difficultés réelles (conflits, préjugés, discrimination). À la fin, Etienne se rapproche de la position de Lucía : \esp{De acuerdo, Lucía}.
\end{enumerate}
\end{corrige}

\courssub{Compréhension détaillée : les arguments des deux personnages}

Un débat, c'est un échange d'\cle{arguments}. Cherchons maintenant dans le texte les arguments précis de chaque locuteur.

\begin{enumerate}
\item Quels sont les \textbf{trois arguments} de Lucía en faveur de la diversité culturelle ?
\item Quels sont les \textbf{deux problèmes} signalés par Etienne ?
\item Quel \textbf{exemple concret} Lucía utilise-t-elle pour convaincre Etienne ?
\item Comment le débat se \textbf{termine}-t-il ? Qui a convaincu qui ?
\end{enumerate}

\begin{corrige}[Compréhension détaillée]
\begin{enumerate}
\item Lucía avance trois arguments : (a) la diversité culturelle est une \textbf{riqueza} (une richesse) ; (b) chaque élève apporte ses coutumes, sa musique et sa langue ; (c) les préjugés naissent de l'ignorance, donc le partage des racines et le dialogue permettent d'apprendre à se respecter.
\item Etienne signale deux problèmes : (a) les \textbf{conflits} causés par les préjugés (la dispute dans la cour à cause d'une moquerie sur l'accent) ; (b) la \textbf{discriminación} : certains élèves refusent de s'asseoir avec des camarades d'autres régions.
\item Lucía utilise l'exemple de l'\textbf{équipe de football} du lycée : des joueurs de dix régions différentes qui gagnent ensemble. C'est un argument par l'exemple, très efficace dans un débat.
\item Le débat se termine par un \textbf{accord} : Etienne reconnaît que le sport unit les gens, puis conclut \esp{De acuerdo, Lucía. Al final, compartimos la misma escuela, el mismo patio y el mismo futuro.} Lucía a donc convaincu Etienne, mais grâce aux objections d'Etienne, sa position est devenue plus nuancée : c'est la marque d'un vrai débat.
\end{enumerate}
\end{corrige}

\courssub{Recherche dans le texte : marqueurs d'opinion et connecteurs}

Pour débattre en espagnol, il faut deux outils linguistiques : les \cle{marqueurs d'opinion} (pour dire « je pense que... ») et les \cle{connecteurs logiques} (pour organiser ses arguments). Souligne-les dans le texte, puis vérifie avec le tableau ci-dessous.

\begin{center}
\begin{tabularx}{\linewidth}{|X|X|X|}
\hline
\rowcolor{popblueL} Español (dans le texte) & Français & Fonction \\
\hline
\esp{creo que} & je crois que & marqueur d'opinion \\
\hline
\esp{para mí} & pour moi & marqueur d'opinion \\
\hline
\esp{en mi opinión} & à mon avis & marqueur d'opinion \\
\hline
\esp{pienso que} & je pense que & marqueur d'opinion \\
\hline
\esp{es verdad} & c'est vrai & concession \\
\hline
\esp{tienes razón (en parte)} & tu as raison (en partie) & accord (nuancé) \\
\hline
\esp{de acuerdo} & d'accord & accord final \\
\hline
\esp{en primer lugar} & en premier lieu, d'abord & connecteur : 1er argument \\
\hline
\esp{además} & de plus, en outre & connecteur : ajout \\
\hline
\esp{sin embargo} & cependant, pourtant & connecteur : opposition \\
\hline
\esp{por último} & enfin, pour finir & connecteur : conclusion \\
\hline
\end{tabularx}
\end{center}

\begin{propriete}[Les connecteurs logiques de l'argumentation]
Les connecteurs logiques sont des mots-outils invariables qui organisent le discours. Dans un débat ou un texte argumenté, on suit presque toujours le même schéma :
\begin{itemize}
\item \esp{En primer lugar...} : on annonce le premier argument.
\item \esp{Además...} : on ajoute un argument du même côté.
\item \esp{Sin embargo...} : on introduit une objection, une nuance ou la thèse adverse.
\item \esp{Por último...} : on donne le dernier argument ou la conclusion.
\end{itemize}
\textbf{Attention à la ponctuation} : en espagnol comme en français, ces connecteurs sont suivis d'une virgule : \esp{Sin embargo, a veces hay conflictos.} « Cependant, il y a parfois des conflits. »

\textbf{Attention à l'orthographe} : \esp{además} (de plus) prend un accent sur le dernier \esp{a} ; ne le confonds pas avec \esp{ademas} sans accent, qui n'existe pas. \esp{Sin embargo} s'écrit en deux mots.
\end{propriete}

\begin{exemplebox}[Les connecteurs en action]
\esp{En primer lugar, la diversidad es una riqueza.} « D'abord, la diversité est une richesse. »

\esp{Además, aprendemos mucho de los demás.} « De plus, nous apprenons beaucoup des autres. »

\esp{Sin embargo, la convivencia exige esfuerzo.} « Cependant, le vivre-ensemble exige un effort. »

\esp{Por último, el respeto es la clave.} « Enfin, le respect est la clé. »
\end{exemplebox}

\begin{attention}[Convivencia ou coexistencia ?]
Ne confonds pas \esp{la convivencia} et \esp{la coexistencia} !

\esp{La coexistencia} désigne le simple fait d'exister côte à côte, sans échange réel : deux groupes peuvent \emph{coexister} en s'ignorant.

\esp{La convivencia}, elle, implique la vie commune active : partager, dialoguer, respecter. C'est le mot juste pour parler du \cle{vivre-ensemble} à l'école ou dans la société.

Dans un texte sur la tolérance, écris donc \esp{convivencia}, pas \esp{coexistencia}. Autre piège : \esp{convivir} (vivre ensemble) ne se traduit pas par « convier » ! Enfin, \esp{la tolerancia} s'écrit avec un seul \esp{l}, contrairement au français « tolérance » qui prend deux \emph{l}... non, un seul aussi : mais méfie-toi du verbe \esp{tolerar}, qui ne s'emploie pas pour « tolérer une douleur » au sens médical ; il signifie « accepter, supporter une personne ou une attitude ».
\end{attention}

\courssub{Carte mentale du lexique du vivre-ensemble}

\begin{popfigure}
\begin{tikzpicture}[mindmap, grow cyclic, every node/.style={concept, text=white}, root concept/.append style={concept color=popblue, font=\Large\bfseries}, level 1/.append style={level distance=3.4cm, sibling angle=72}, concept color=popblue]
\node[root concept]{CONVIVENCIA}
  child[concept color=popgreen]{ node {tolerancia} }
  child[concept color=poporange]{ node {respeto} }
  child[concept color=poppurple]{ node {diversidad} }
  child[concept color=poppink]{ node {diálogo} }
  child[concept color=popgold]{ node {identidad} };
\end{tikzpicture}
\legende{Carte mentale : les cinq piliers du vivre-ensemble. La \esp{convivencia} repose sur la \esp{tolerancia} (accepter), le \esp{respeto} (respecter), la \esp{diversidad} (valoriser les différences), le \esp{diálogo} (parler et écouter) et l'\esp{identidad} (connaître ses propres racines).}
\end{popfigure}

\courssub{Méthode : comment lire un dialogue-débat}

\begin{methode}[Lire et analyser un dialogue-débat en quatre étapes]
\begin{enumerate}
\item \textbf{Identifier les locuteurs.} Qui parle ? Quel est leur rapport (amis, adversaires, camarades) ? Où et quand se déroule l'échange ? Ici : deux camarades de lycée à Yaoundé, après la lecture d'un article.
\item \textbf{Souligner les marqueurs d'opinion.} Repère \esp{creo que}, \esp{pienso que}, \esp{en mi opinión}, \esp{para mí} : ils signalent la prise de position personnelle.
\item \textbf{Relever les connecteurs logiques.} \esp{En primer lugar}, \esp{además}, \esp{sin embargo}, \esp{por último} : ils dessinent le plan du débat (arguments, objection, conclusion).
\item \textbf{Dégager la thèse de chacun.} Résume en une phrase la position de chaque locuteur, puis observe comment le débat évolue : accord total, accord partiel ou désaccord final.
\end{enumerate}
\end{methode}

\courssub{Exercice résolu : reformuler le débat}

\begin{exoresolu}[Resumir el debate en dos frases]
Énoncé. Résume le débat entre Lucía et Etienne en \textbf{deux phrases} en espagnol, en utilisant au moins un marqueur d'opinion et un connecteur logique.
\tcblower
Solution détaillée. Un bon résumé doit contenir : le sujet (la diversité culturelle), la thèse de Lucía, l'objection d'Etienne et la conclusion commune. Modèle de réponse :

\esp{Lucía piensa que la diversidad cultural es una riqueza porque cada alumno trae sus costumbres y su lengua. Sin embargo, Etienne señala los prejuicios y la discriminación, pero al final los dos están de acuerdo : el diálogo y el respeto son la clave de la convivencia.}

« Lucía pense que la diversité culturelle est une richesse parce que chaque élève apporte ses coutumes et sa langue. Cependant, Etienne souligne les préjugés et la discrimination, mais finalement tous les deux sont d'accord : le dialogue et le respect sont la clé du vivre-ensemble. »

Remarque la méthode : une phrase pour la thèse de Lucía, une phrase (avec \esp{sin embargo} et \esp{al final}) pour l'objection d'Etienne et la conclusion commune. Deux phrases, tout le débat. Note aussi l'emploi de \esp{estar de acuerdo} (être d'accord) : on dit \esp{los dos están de acuerdo}, avec \esp{estar} et non \esp{ser}.
\end{exoresolu}

\begin{exercice}[À ton tour]
\begin{enumerate}
\item Lis le dialogue à voix haute avec un camarade : l'un joue Lucía, l'autre Etienne. Respecte l'intonation du débat (affirmation, objection, concession).
\item Retrouve dans le texte trois mots de la famille du vivre-ensemble et donne leur traduction française.
\item Classe ces expressions en deux colonnes, « opinion » et « connecteur » : \esp{además}, \esp{creo que}, \esp{sin embargo}, \esp{en mi opinión}, \esp{por último}, \esp{para mí}.
\item Reformule la position d'Etienne en une phrase commençant par \esp{Etienne piensa que...}.
\end{enumerate}
\end{exercice}

\begin{corrige}[À ton tour]
\begin{enumerate}
\item Activité orale : veiller à l'intonation montante dans les questions (\esp{¿has leído...?}) et à l'accent d'insistance sur les marqueurs d'opinion (\esp{creo que}, \esp{en mi opinión}).
\item Par exemple : \esp{la convivencia} (le vivre-ensemble), \esp{la tolerancia} (la tolérance), \esp{el respeto} (le respect), \esp{la diversidad} (la diversité).
\item Opinion : \esp{creo que}, \esp{en mi opinión}, \esp{para mí}. Connecteurs : \esp{además}, \esp{sin embargo}, \esp{por último}.
\item Modèle : \esp{Etienne piensa que la convivencia no es fácil porque hay prejuicios y discriminación entre los alumnos.} « Etienne pense que le vivre-ensemble n'est pas facile parce qu'il y a des préjugés et de la discrimination entre les élèves. »
\end{enumerate}
\end{corrige}

Tu sais maintenant lire et analyser un dialogue-débat : identifier les locuteurs, repérer les marqueurs d'opinion, suivre les connecteurs et dégager les thèses. Dans la prochaine partie de la leçon, nous enrichirons le lexique du débat et du vivre-ensemble pour que tu puisses, à ton tour, prendre la parole.

\coursec{Lexique thématique : débat et vivre-ensemble}

Dans la section précédente, nous avons lu le dialogue \esp{Un debate en el aula}, où des élèves d'un lycée de Yaoundé discutaient de la tolérance et de la discrimination. Ce texte regorgeait de mots nouveaux : \esp{respeto}, \esp{prejuicio}, \esp{estar de acuerdo}, \esp{opinar}... Pour participer vous-même à un débat en espagnol et rédiger un texte argumenté, il faut d'abord maîtriser ce vocabulaire. C'est l'objectif de cette section : organiser, comprendre et employer les mots du \cle{débat} et du \cle{vivre-ensemble}.

\courssub{Observer d'abord : les mots en contexte}

Avant toute liste de vocabulaire, observons un court texte qui rassemble naturellement les mots essentiels de cette leçon. Lisez-le deux fois : une première fois pour le sens général, une deuxième fois pour repérer les mots des quatre champs lexicaux.

\begin{textebox}[La convivencia en nuestro barrio]
En nuestro barrio de Yaoundé viven personas de diferentes regiones y culturas. En mi opinión, esta diversidad es una riqueza, pero exige también un esfuerzo de todos. La convivencia no es automática: hay que construirla cada día.

Primero, el respeto mutuo es fundamental. Respetamos las costumbres de los demás, aunque sean diferentes de las nuestras. Mi vecina, por ejemplo, habla ewondo con sus hijos, y mi familia habla duala: nadie tiene derecho a burlarse de la lengua materna del otro. La igualdad significa que todos tenemos los mismos derechos, sin importar nuestra región de origen.

Sin embargo, a veces aparecen prejuicios. Algunas personas discriminan a los recién llegados y los excluyen de las actividades del barrio. Esta intolerancia es peligrosa: el racismo y la xenofobia destruyen la convivencia. A mi juicio, la mejor arma contra la discriminación es el diálogo. Cuando hablamos con los demás, descubrimos que nuestros prejuicios no tienen razón de ser.

En la escuela, aprendemos a debatir: opinamos, argumentamos, defendemos nuestras ideas, pero también escuchamos y aceptamos refutar o ser refutados. Estar de acuerdo no es obligatorio; respetar al otro, sí. Como dice nuestro profesor: « Podemos no estar de acuerdo con una idea, pero debemos respetar a la persona ». Esa es la verdadera tolerancia: no una simple paciencia, sino un respeto activo de las diferencias.
\end{textebox}

\textbf{Glossaire}

\begin{center}
\begin{tabularx}{\linewidth}{|X|X|}\hline
\rowcolor{popblueL} \textbf{Español} & \textbf{Français} \\ \hline
el barrio & le quartier \\ \hline
el esfuerzo & l'effort \\ \hline
burlarse de & se moquer de \\ \hline
los recién llegados & les nouveaux arrivants \\ \hline
no tener razón de ser & ne pas avoir de raison d'être \\ \hline
refutar / ser refutado & réfuter / être réfuté \\ \hline
sino & mais (après négation) \\ \hline
\end{tabularx}
\end{center}

\textbf{Questions de compréhension}
\begin{enumerate}
\item ¿Qué es fundamental para la convivencia, según el texto?
\item ¿Qué peligros amenazan la convivencia en el barrio?
\item Según el profesor, ¿qué debemos respetar siempre?
\end{enumerate}

\begin{definition}[La convivencia]
\esp{La convivencia} signifie \esp{vivir juntos en armonía respetando las diferencias} : vivre ensemble en harmonie en respectant les différences. C'est le mot-clé de cette leçon : il désigne à la fois le fait de vivre ensemble et l'art de le faire pacifiquement. Prononciation : « kon-bi-bén-sia ».
\end{definition}

\courssub{Les quatre champs lexicaux}

Un \cle{champ lexical} est un ensemble de mots qui se rapportent à une même idée. Retenir les mots par familles (et non isolément) est la méthode la plus efficace : chaque mot en rappelle un autre. Voici les quatre familles de cette leçon.

\begin{center}
\begin{tabularx}{\linewidth}{|X|X|}\hline
\rowcolor{popblueL} \textbf{Español} & \textbf{Français} \\ \hline
\multicolumn{2}{|l|}{\cellcolor{popgreenL}\textbf{1. La tolerancia}} \\ \hline
tolerar & tolérer \\ \hline
respetar & respecter \\ \hline
aceptar & accepter \\ \hline
la aceptación & l'acceptation \\ \hline
el respeto mutuo & le respect mutuel \\ \hline
la igualdad & l'égalité \\ \hline
\multicolumn{2}{|l|}{\cellcolor{popgreenL}\textbf{2. La discriminación}} \\ \hline
discriminar & discriminer \\ \hline
el racismo & le racisme \\ \hline
la xenofobia & la xénophobie \\ \hline
el prejuicio & le préjugé \\ \hline
excluir & exclure \\ \hline
la exclusión & l'exclusion \\ \hline
la intolerancia & l'intolérance \\ \hline
\multicolumn{2}{|l|}{\cellcolor{popgreenL}\textbf{3. El debate}} \\ \hline
opinar & donner son avis \\ \hline
argumentar & argumenter \\ \hline
defender una idea & défendre une idée \\ \hline
refutar & réfuter \\ \hline
estar de acuerdo / no estar de acuerdo & être d'accord / ne pas être d'accord \\ \hline
tener razón & avoir raison \\ \hline
\multicolumn{2}{|l|}{\cellcolor{popgreenL}\textbf{4. La identidad}} \\ \hline
la identidad & l'identité \\ \hline
las raíces & les racines \\ \hline
la cultura & la culture \\ \hline
la lengua materna & la langue maternelle \\ \hline
las costumbres & les coutumes \\ \hline
la pertenencia & l'appartenance \\ \hline
\end{tabularx}
\end{center}

\begin{exemplebox}[Les mots en phrases]
\begin{itemize}
\item \esp{Respetamos las costumbres de los demás.} « Nous respectons les coutumes des autres. »
\item \esp{La igualdad es un derecho de todos los ciudadanos.} « L'égalité est un droit de tous les citoyens. »
\item \esp{No debemos excluir a nadie por su lengua materna.} « Nous ne devons exclure personne à cause de sa langue maternelle. »
\item \esp{Los prejuicios nacen de la ignorancia.} « Les préjugés naissent de l'ignorance. »
\item \esp{Mis raíces camerunesas son parte de mi identidad.} « Mes racines camerounaises font partie de mon identité. »
\item \esp{En el debate, cada alumno defiende su idea con argumentos.} « Dans le débat, chaque élève défend son idée avec des arguments. »
\end{itemize}
\end{exemplebox}

\begin{popfigure}
\begin{center}
\begin{tikzpicture}
\node[draw=popgreen, fill=popgreenL, rounded corners, thick, align=center, minimum width=4.2cm, minimum height=1.1cm, font=\bfseries] (tol) at (0,0) {TOLERANCIA};
\node[draw=popgreen, fill=white, rounded corners, align=center] (res) at (-2.6,1.6) {respeto};
\node[draw=popgreen, fill=white, rounded corners, align=center] (igu) at (0,2.1) {igualdad};
\node[draw=popgreen, fill=white, rounded corners, align=center] (dia) at (2.6,1.6) {diálogo};
\node[draw=poporange, fill=poporangeL, rounded corners, thick, align=center, minimum width=4.2cm, minimum height=1.1cm, font=\bfseries] (dis) at (0,-3.4) {DISCRIMINACIÓN};
\node[draw=poporange, fill=white, rounded corners, align=center] (rac) at (-2.6,-5) {racismo};
\node[draw=poporange, fill=white, rounded corners, align=center] (pre) at (0,-5.5) {prejuicio};
\node[draw=poporange, fill=white, rounded corners, align=center] (exc) at (2.6,-5) {exclusión};
\draw[-{Stealth}, thick, popgreen] (tol) -- (res);
\draw[-{Stealth}, thick, popgreen] (tol) -- (igu);
\draw[-{Stealth}, thick, popgreen] (tol) -- (dia);
\draw[-{Stealth}, thick, poporange] (dis) -- (rac);
\draw[-{Stealth}, thick, poporange] (dis) -- (pre);
\draw[-{Stealth}, thick, poporange] (dis) -- (exc);
\draw[{Stealth}-{Stealth}, very thick, poppurple, dashed] (tol) -- (dis) node[midway, right, align=left, font=\small\itshape, text=popdark]{dos polos\\opuestos};
\end{tikzpicture}
\end{center}
\legende{Deux pôles opposés : la tolérance construit la convivencia, la discrimination la détruit.}
\end{popfigure}

\begin{attention}[Familles de mots : verbe, nom, adjectif]
En espagnol comme en français, les mots d'une même famille changent de terminaison selon leur fonction. Apprenez-les ensemble : \esp{tolerar} (verbe) → \esp{la tolerancia} (nom) → \esp{tolerante} (adjectif) ; \esp{excluir} (verbe) → \esp{la exclusión} (nom) → \esp{excluido} (participe/adjectif). Attention à \esp{excluir} : son nom est \esp{exclusión} avec un \textbf{s}, pas \emph{exclución}.
\end{attention}

\courssub{Les expressions du débat : opinion, accord, désaccord}

Pour débattre, le vocabulaire ne suffit pas : il faut des \cle{expressions toutes faites} (des « formules ») qui permettent de prendre la parole poliment. Retenez-les par cœur : ce sont des blocs indissociables.

\textbf{Exprimer son opinion}

\begin{center}
\begin{tabularx}{\linewidth}{|X|X|X|}\hline
\rowcolor{popblueL} \textbf{Español} & \textbf{Français} & \textbf{Exemple} \\ \hline
en mi opinión & à mon avis & \esp{En mi opinión, la diversidad es una riqueza.} \\ \hline
a mi juicio & à mon avis / selon moi & \esp{A mi juicio, el diálogo es necesario.} \\ \hline
desde mi punto de vista & de mon point de vue & \esp{Desde mi punto de vista, todos somos iguales.} \\ \hline
me parece que & il me semble que & \esp{Me parece que tienes razón.} \\ \hline
\end{tabularx}
\end{center}

\textbf{Exprimer l'accord}

\begin{center}
\begin{tabularx}{\linewidth}{|X|X|}\hline
\rowcolor{popgreenL} \textbf{Español} & \textbf{Français} \\ \hline
estoy de acuerdo (contigo) & je suis d'accord (avec toi) \\ \hline
es verdad & c'est vrai \\ \hline
tienes razón & tu as raison \\ \hline
efectivamente & effectivement \\ \hline
\end{tabularx}
\end{center}

\textbf{Exprimer le désaccord poli}

\begin{center}
\begin{tabularx}{\linewidth}{|X|X|}\hline
\rowcolor{poporangeL} \textbf{Español} & \textbf{Français} \\ \hline
no estoy de acuerdo & je ne suis pas d'accord \\ \hline
sin embargo & cependant, pourtant \\ \hline
puede ser, pero... & c'est possible, mais... \\ \hline
no lo creo & je ne le crois pas \\ \hline
\end{tabularx}
\end{center}

\begin{exemplebox}[Un mini-échange typique]
\esp{— En mi opinión, las lenguas locales deben enseñarse en la escuela.} « À mon avis, les langues locales doivent être enseignées à l'école. »

\esp{— Estoy de acuerdo contigo: son parte de nuestra identidad.} « Je suis d'accord avec toi: elles font partie de notre identité. »

\esp{— Puede ser, pero sin embargo, el español une a todos los cameruneses hispanohablantes.} « C'est possible, mais cependant, l'espagnol unit tous les Camerounais hispanophones. »
\end{exemplebox}

\begin{propriete}[La construction de « estar de acuerdo »]
\esp{Estar de acuerdo} se construit avec la préposition \esp{con}: \esp{estar de acuerdo \textbf{con} alguien} (être d'accord avec quelqu'un) ou \esp{estar de acuerdo \textbf{con} una idea} (être d'accord avec une idée). Jamais avec \esp{de}.

\esp{Estoy de acuerdo con María.} « Je suis d'accord avec María. »

\esp{No estamos de acuerdo con esa opinión.} « Nous ne sommes pas d'accord avec cette opinion. »
\end{propriete}

\begin{attention}[Pièges fréquents des francophones]
\begin{itemize}
\item Ne dites jamais \emph{estar de acuerdo de alguien}: c'est un calque du français « être d'accord \textbf{de} ». La bonne préposition est \esp{con}.
\item \esp{Tener razón} (avoir raison) se construit avec \esp{tener}, pas avec \esp{ser}: \esp{Tú tienes razón}, jamais \emph{tú eres razón}.
\item \esp{Actualmente} est un faux ami: il signifie « actuellement, en ce moment », et non « effectivement ». Pour « effectivement », dites \esp{efectivamente} ou \esp{en efecto}.
\item \esp{Sin embargo} (cependant) s'écrit en deux mots et sert à nuancer un désaccord poliment: \esp{Entiendo tu idea; sin embargo, no lo creo.} « Je comprends ton idée; cependant, je ne le crois pas. »
\end{itemize}
\end{attention}

\begin{savaistu}[D'où vient le mot « tolerancia »?]
Le mot \esp{tolerancia} vient du latin \emph{tolerare}, qui signifiait « supporter, endurer ». En espagnol moderne, le sens a évolué: la tolérance n'est plus une simple patience envers ce qui nous dérange, mais un \cle{respect actif} des différences de l'autre. Tolérer, aujourd'hui, c'est reconnaître que l'autre a le droit d'être différent. C'est exactement l'esprit de la \esp{convivencia}: non pas supporter l'autre malgré ses différences, mais vivre avec lui grâce à elles.
\end{savaistu}

\begin{exoresolu}[Compléter avec le mot juste]
Énoncé. Complétez chaque phrase avec le mot ou l'expression qui convient: \esp{respeto mutuo}, \esp{prejuicios}, \esp{tienes razón}, \esp{en mi opinión}, \esp{lengua materna}, \esp{excluir}.

\begin{enumerate}
\item \esp{............, la xenofobia es un peligro para la convivencia.}
\item \esp{No debemos ............ a los alumnos que hablan otra lengua.}
\item \esp{Mi ............ es el ewondo, pero también hablo español.}
\item \esp{— El diálogo resuelve los conflictos. — Sí, ............ .}
\item \esp{El ............ entre los vecinos hace la vida más fácil.}
\item \esp{Los ............ desaparecen cuando conocemos a los demás.}
\end{enumerate}
\tcblower
Solution détaillée.
\begin{enumerate}
\item \esp{En mi opinión, la xenofobia es un peligro para la convivencia.} — La virgule et l'opinion personnelle appellent une formule d'opinion.
\item \esp{No debemos excluir a los alumnos...} — Il faut un verbe à l'infinitif après \esp{debemos}; le sens (« mettre à part ») impose \esp{excluir}. Attention à la préposition \esp{a} devant la personne (a personnelle).
\item \esp{Mi lengua materna es el ewondo...} — On parle d'une langue apprise dès l'enfance.
\item \esp{Sí, tienes razón.} — Réponse d'accord à une affirmation; \esp{tener razón} avec \esp{tener}.
\item \esp{El respeto mutuo entre los vecinos...} — Nom composé du champ de la tolérance.
\item \esp{Los prejuicios desaparecen...} — Sujet pluriel au sens de « idées reçues ».
\end{enumerate}
\end{exoresolu}

\begin{exercice}[À vous de jouer]
\begin{enumerate}
\item Traduisez en espagnol: a) « Je suis d'accord avec toi. » b) « Nous respectons les coutumes des autres. » c) « À mon avis, l'égalité est importante. » d) « C'est possible, mais je ne le crois pas. »
\item Classez ces mots dans les deux pôles TOLERANCIA / DISCRIMINACIÓN: \esp{igualdad, racismo, diálogo, exclusión, aceptación, xenofobia, respeto, intolerancia}.
\item Rédigez deux répliques de débat: une d'accord et une de désaccord poli, à partir de l'affirmation \esp{Las costumbres de cada región deben conservarse}.
\end{enumerate}
\end{exercice}

\begin{corrige}[Exercice « À vous de jouer »]
\begin{enumerate}
\item a) \esp{Estoy de acuerdo contigo.} b) \esp{Respetamos las costumbres de los demás.} c) \esp{En mi opinión, la igualdad es importante.} d) \esp{Puede ser, pero no lo creo.}
\item TOLERANCIA: \esp{igualdad, diálogo, aceptación, respeto}. DISCRIMINACIÓN: \esp{racismo, exclusión, xenofobia, intolerancia}.
\item Exemples: \esp{Estoy de acuerdo: las costumbres son parte de nuestra identidad.} / \esp{Puede ser, pero sin embargo, algunas costumbres deben cambiar con el tiempo.}
\end{enumerate}
\end{corrige}

\begin{aretenir}[L'essentiel du lexique]
\begin{itemize}
\item Quatre champs lexicaux: la \esp{tolerancia} (respetar, aceptar, igualdad), la \esp{discriminación} (racismo, prejuicio, excluir), el \esp{debate} (opinar, argumentar, refutar), la \esp{identidad} (raíces, lengua materna, costumbres).
\item Trois familles de formules: opinion (\esp{en mi opinión}, \esp{me parece que}), accord (\esp{estoy de acuerdo}, \esp{tienes razón}), désaccord poli (\esp{no estoy de acuerdo}, \esp{puede ser, pero...}).
\item Construction obligatoire: \esp{estar de acuerdo \textbf{con} alguien}; \esp{tener razón} avec le verbe \esp{tener}.
\item La \esp{convivencia}: vivre ensemble en harmonie en respectant les différences.
\end{itemize}
\end{aretenir}

Ce lexique maîtrisé, nous pourrons passer à la section suivante: les connecteurs logiques (\esp{en primer lugar}, \esp{además}, \esp{sin embargo}, \esp{por último}), qui permettront d'organiser ces mots et ces formules dans un véritable texte argumenté.

\coursec{Grammaire : les connecteurs logiques de l'argumentation}

Dans la section précédente, nous avons découvert le lexique du débat et du vivre-ensemble : \esp{la tolerancia}, \esp{la discriminación}, \esp{el respeto}, \esp{la convivencia}. Mais posséder du vocabulaire ne suffit pas pour convaincre. Dans le débat du lycée que nous avons lu, les élèves ne juxtaposent pas leurs idées : ils les \cle{organisent}, les \cle{relient} et les \cle{opposent} grâce à de petits mots-outils appelés \cle{connecteurs logiques} (\esp{conectores lógicos}). C'est eux qui transforment une simple liste d'opinions en un véritable raisonnement. Cette section est consacrée à ces outils indispensables de l'argumentation, à l'oral comme à l'écrit.

\courssub{Observation : repérer les connecteurs dans un texte}

Lisons d'abord ce court extrait d'une prise de parole lors du débat sur la tolérance au lycée de Yaoundé.

\begin{textebox}[Intervención de Amina en el debate]
En mi opinión, la tolerancia es la base de la convivencia en nuestro liceo. En primer lugar, somos estudiantes de regiones y culturas diferentes, y debemos respetarnos. Además, la tolerancia favorece la paz y la amistad entre nosotros. Sin embargo, todavía hay personas intolerantes que rechazan a los demás por su lengua o su religión. Por eso, es necesario educar en el respeto desde la escuela. Por último, quiero decir que un debate como este nos ayuda a comprendernos mejor. En conclusión, tolerar no es aceptarlo todo, sino escuchar y dialogar.
\end{textebox}

\textbf{Glossaire :} \esp{la base} = la base ; \esp{respetarnos} = nous respecter (mutuellement) ; \esp{rechazar} = rejeter ; \esp{todavía} = encore ; \esp{comprendernos} = nous comprendre (mutuellement) ; \esp{sino} = mais (après négation).

\textbf{Questions de compréhension :}
\begin{enumerate}
\item Quelle est l'opinion (\esp{opinión}) d'Amina sur la tolérance ?
\item Combien d'arguments donne-t-elle ? Souligne les mots qui les introduisent.
\item Quel connecteur annonce une difficulté, une objection ?
\item Quels connecteurs annoncent la fin du raisonnement ?
\end{enumerate}

Observons maintenant la phrase modèle suivante, qui résume la mécanique du texte :

\begin{exemplebox}[La mécanique des connecteurs]
\esp{En primer lugar, la tolerancia es necesaria. Además, favorece la paz. Sin embargo, hay personas intolerantes. Por último, debemos educar en el respeto.}

« D'abord, la tolérance est nécessaire. De plus, elle favorise la paix. Cependant, il y a des personnes intolérantes. Enfin, nous devons éduquer au respect. »
\end{exemplebox}

On remarque que chaque connecteur joue un rôle précis : \esp{en primer lugar} ouvre la liste des arguments, \esp{además} ajoute, \esp{sin embargo} oppose, \esp{por último} conclut. Chacun est suivi d'une \cle{virgule}. C'est exactement le squelette d'un texte argumenté.

\courssub{Les cinq fonctions des connecteurs logiques}

\textbf{Fonction 1 — Ordonner les arguments}

Ces connecteurs annoncent la place de l'argument dans le plan : premier, deuxième, dernier.

\begin{definition}[Connecteurs d'ordre]
\esp{en primer lugar} = en premier lieu, d'abord (registre soutenu) ; \esp{en segundo lugar} = en deuxième lieu ; \esp{luego} = ensuite ; \esp{por último} = enfin (dernier argument) ; \esp{finalmente} = finalement (après une attente).
\end{definition}

\begin{exemplebox}[Exemples traduits]
\esp{En primer lugar, la escuela debe enseñar el respeto.} « En premier lieu, l'école doit enseigner le respect. »

\esp{En segundo lugar, los medios de comunicación tienen una gran responsabilidad.} « En deuxième lieu, les médias ont une grande responsabilité. »

\esp{Luego, cada ciudadano debe rechazar la discriminación.} « Ensuite, chaque citoyen doit rejeter la discrimination. »

\esp{Por último, el diálogo es siempre mejor que la violencia.} « Enfin, le dialogue est toujours meilleur que la violence. »
\end{exemplebox}

\textbf{Fonction 2 — Ajouter une idée}

\begin{definition}[Connecteurs d'addition]
\esp{además} = de plus, en outre ; \esp{también} = aussi ; \esp{asimismo} = de même (soutenu) ; \esp{incluso} = même, jusqu'à.
\end{definition}

\begin{exemplebox}[Exemples traduits]
\esp{La tolerancia une a las personas. Además, crea un ambiente de paz.} « La tolérance unit les personnes. De plus, elle crée un climat de paix. »

\esp{El respeto también se aprende en la familia.} « Le respect s'apprend aussi dans la famille. »

\esp{En nuestro liceo hay incluso estudiantes de otros países.} « Dans notre lycée, il y a même des élèves d'autres pays. »
\end{exemplebox}

\textbf{Fonction 3 — Opposer, nuancer, réfuter}

C'est la fonction la plus importante dans un débat : elle permet de présenter l'avis contraire avant de le réfuter.

\begin{definition}[Connecteurs d'opposition]
\esp{sin embargo} = cependant, pourtant ; \esp{pero} = mais ; \esp{no obstante} = néanmoins (soutenu) ; \esp{en cambio} = en revanche, par contre ; \esp{por el contrario} = au contraire.
\end{definition}

\begin{exemplebox}[Exemples traduits]
\esp{Todos somos iguales; sin embargo, algunos sufren discriminación.} « Nous sommes tous égaux ; cependant, certains subissent la discrimination. »

\esp{La convivencia es difícil, pero no imposible.} « Le vivre-ensemble est difficile, mais pas impossible. »

\esp{Algunos rechazan el diálogo; no obstante, la mayoría lo apoya.} « Certains rejettent le dialogue ; néanmoins, la majorité le soutient. »

\esp{Pedro prefiere discutir; en cambio, María prefiere escuchar.} « Pedro préfère discuter ; en revanche, María préfère écouter. »
\end{exemplebox}

\textbf{Fonction 4 — Conclure}

\begin{definition}[Connecteurs de conclusion]
\esp{por último} = enfin (dernier point du plan) ; \esp{en conclusión} = en conclusion ; \esp{en resumen} = en résumé ; \esp{por eso} = c'est pourquoi ; \esp{por lo tanto} = par conséquent, donc.
\end{definition}

\begin{exemplebox}[Exemples traduits]
\esp{En conclusión, la tolerancia es un deber de todos.} « En conclusion, la tolérance est un devoir de tous. »

\esp{La discriminación destruye la sociedad; por lo tanto, debemos combatirla.} « La discrimination détruit la société ; par conséquent, nous devons la combattre. »

\esp{No queremos conflictos; por eso dialogamos.} « Nous ne voulons pas de conflits ; c'est pourquoi nous dialoguons. »
\end{exemplebox}

\textbf{Fonction 5 — Donner un exemple (complément Bac)}

\begin{definition}[Connecteurs d'exemple et d'explication]
\esp{por ejemplo} = par exemple ; \esp{es decir} = c'est-à-dire.
\end{definition}

\begin{exemplebox}[Exemples traduits]
\esp{Hay muchas formas de discriminación, por ejemplo, el racismo.} « Il y a beaucoup de formes de discrimination, par exemple le racisme. »

\esp{La convivencia, es decir, vivir juntos en paz, es nuestro objetivo.} « Le vivre-ensemble, c'est-à-dire vivre ensemble en paix, est notre objectif. »
\end{exemplebox}

\courssub{La règle : forme, place et ponctuation}

\begin{propriete}[Règle des connecteurs logiques]
\begin{itemize}
\item Les \cle{connecteurs d'ordre} (\esp{en primer lugar}, \esp{en segundo lugar}, \esp{luego}, \esp{por último}) structurent le plan du texte : ils annoncent la place de chaque argument.
\item Les \cle{connecteurs d'addition} (\esp{además}, \esp{también}, \esp{incluso}) renforcent l'argumentation en apportant un élément nouveau.
\item Les \cle{connecteurs d'opposition} (\esp{sin embargo}, \esp{no obstante}) introduisent la nuance ou la réfutation d'une idée contraire.
\item \textbf{Place :} le connecteur se met généralement \cle{en tête de phrase}, suivi d'une \cle{virgule} : \esp{Además, la tolerancia favorece la paz.}
\item \textbf{Ponctuation spéciale :} \esp{sin embargo} et \esp{no obstante} peuvent être précédés d'un \cle{point-virgule} pour relier deux phrases indépendantes : \esp{Somos diferentes; sin embargo, somos iguales en derechos.}
\item \esp{pero} se place au milieu de la phrase, après une virgule : \esp{Es difícil, pero necesario.}
\end{itemize}
\end{propriete}

\begin{attention}[\esp{sin embargo} : un bloc figé]
\esp{Sin embargo} signifie « cependant ». C'est un \cle{bloc figé} : il ne se traduit pas mot à mot et ne se décompose jamais pour l'apprenant. N'écris jamais \esp{sin embargo} en pensant à « sans embargo » : le nom \esp{embargo} existe en espagnol (sens juridique : « saisie, embargo commercial »), mais dans cette locution, le sens est figé. De même, ne confonds pas \esp{sin embargo} (cependant) avec \esp{sin} + nom : \esp{sin respeto} = sans respect.
\end{attention}

\courssub{Comparaison français / espagnol : les faux amis des connecteurs}

Le français et l'espagnol partagent beaucoup de connecteurs, mais certains sont des pièges pour les francophones.

\begin{center}
\begin{tabularx}{\linewidth}{|X|X|X|}
\hline
\rowcolor{popblueL} Français & Español & Remarque \\
\hline
d'abord (familier) / en premier lieu & en primer lugar & registre soutenu ; évite \esp{primeramente} à l'écrit \\
\hline
de plus, en outre & además & ne pas confondre avec \esp{además de} + nom (= en plus de) \\
\hline
cependant, pourtant & sin embargo / no obstante & \esp{no obstante} est plus soutenu \\
\hline
en revanche, par contre & en cambio & contraste entre deux éléments \\
\hline
enfin (dernier argument) & por último & \esp{finalmente} = finalement (après une attente) \\
\hline
par conséquent, donc & por lo tanto / por eso & \esp{por eso} est plus courant à l'oral \\
\hline
c'est-à-dire & es decir & introduit une explication \\
\hline
en fait & en realidad / de hecho & jamais \esp{en facto} ! \\
\hline
actuellement & actualmente & \esp{actualmente} = actuellement (vrai ami) \\
\hline
\end{tabularx}
\end{center}

\begin{attention}[Faux amis à retenir absolument]
\begin{itemize}
\item « En fait » se traduit \esp{en realidad} ou \esp{de hecho}, jamais \esp{en facto} (ce mot n'existe pas : c'est un calque du français).
\item \esp{Actualmente} signifie bien « actuellement » : \esp{Actualmente vivo en Yaoundé.} « Actuellement, je vis à Yaoundé. » Mais attention au sens inverse : si tu veux insister sur la réalité par opposition à l'apparence (« en fait, en réalité »), dis \esp{en realidad}.
\item « Enfin » n'est pas toujours \esp{finalmente} : pour le dernier argument d'un plan, utilise \esp{por último}. \esp{Finalmente} exprime le résultat après une attente : \esp{Finalmente llegó el autobús.} « Finalement, le bus est arrivé. »
\item « D'abord » familier ne se traduit pas par \esp{en primer lugar} dans une conversation : à l'oral spontané, on dira plutôt \esp{primero} ou \esp{lo primero}. \esp{En primer lugar} appartient au registre soutenu du débat et de la rédaction.
\end{itemize}
\end{attention}

\courssub{Schéma : la structure d'un texte argumenté}

\begin{popfigure}
\begin{tikzpicture}[
  node distance=0.55cm,
  box/.style={draw=popblue, fill=popblueL, rounded corners=3pt, align=center, text width=4.6cm, inner sep=5pt, font=\small},
  conn/.style={draw=poporange, fill=poporangeL, rounded corners=3pt, align=center, text width=3.4cm, inner sep=4pt, font=\small\itshape},
  fleche/.style={-{Stealth[length=2.5mm]}, thick, popdark}
]
\node[box, fill=poppurpleL, draw=poppurple, align=center] (intro){\textbf{Introducción}\\ opinión : \esp{En mi opinión...}};
\node[conn, below=of intro] (c1) {\esp{En primer lugar}};
\node[box, below=of c1, align=center] (arg1){\textbf{Argumento 1}\\ la tolerancia es necesaria};
\node[conn, below=of arg1] (c2) {\esp{Además}};
\node[box, below=of c2, align=center] (arg2){\textbf{Argumento 2}\\ favorece la paz};
\node[conn, below=of arg2] (c3) {\esp{Sin embargo}};
\node[box, below=of c3, align=center] (obj){\textbf{Objeción}\\ hay personas intolerantes};
\node[conn, below=of obj] (c4) {\esp{Por último / En conclusión}};
\node[box, fill=popgreenL, draw=popgreen, below=of c4, align=center] (conc){\textbf{Conclusión}\\ debemos educar en el respeto};
\draw[fleche] (intro) -- (c1);
\draw[fleche] (c1) -- (arg1);
\draw[fleche] (arg1) -- (c2);
\draw[fleche] (c2) -- (arg2);
\draw[fleche] (arg2) -- (c3);
\draw[fleche] (c3) -- (obj);
\draw[fleche] (obj) -- (c4);
\draw[fleche] (c4) -- (conc);
\end{tikzpicture}
\legende{Organigramme : la structure d'un texte argumenté et ses connecteurs}
\end{popfigure}

\courssub{Méthode : choisir le bon connecteur}

\begin{methode}[Trois questions pour choisir son connecteur]
\begin{enumerate}
\item \textbf{Quelle est la fonction de ma phrase ?} Avant d'écrire, demande-toi : est-ce que j'\cle{ajoute} une idée semblable, est-ce que j'\cle{oppose} une idée contraire, ou est-ce que je \cle{conclus} ?
\item \textbf{Si j'ajoute} : utilise \esp{además}, \esp{también} ou \esp{incluso}. Si je numérote mes arguments : \esp{en primer lugar}, \esp{en segundo lugar}, \esp{luego}.
\item \textbf{Si j'oppose} : utilise \esp{sin embargo} (en tête de phrase ou après un point-virgule), \esp{pero} (au milieu) ou \esp{en cambio} (contraste entre deux personnes ou situations).
\item \textbf{Si je conclus} : utilise \esp{por último} pour le dernier argument, \esp{en conclusión} ou \esp{en resumen} pour la conclusion générale, \esp{por eso} ou \esp{por lo tanto} pour la conséquence logique.
\item \textbf{Vérifie la ponctuation} : connecteur en tête de phrase = virgule obligatoire derrière lui.
\end{enumerate}
\end{methode}

\courssub{Exercice résolu et entraînement}

\begin{exoresolu}[Complète avec le connecteur qui convient]
Énoncé. Complète ce paragraphe avec les connecteurs suivants : \esp{además}, \esp{sin embargo}, \esp{en primer lugar}, \esp{por último}, \esp{por ejemplo}.

\esp{\ldots, la convivencia es posible si nos respetamos. \ldots, el diálogo resuelve los conflictos sin violencia. \ldots, todavía existe la discriminación en algunas escuelas; \ldots, hay estudiantes que rechazan a los compañeros de otras regiones. \ldots, debemos educar en la tolerancia desde la infancia.}

\tcblower

Solution détaillée. On applique la méthode des trois questions.

\begin{itemize}
\item Première phrase : elle ouvre le raisonnement, c'est le premier argument $\rightarrow$ \esp{En primer lugar}.
\item Deuxième phrase : elle ajoute une idée du même type (le dialogue aide aussi) $\rightarrow$ \esp{Además}.
\item Troisième phrase : elle introduit une réalité contraire (la discrimination existe encore) $\rightarrow$ \esp{sin embargo}.
\item Quatrième phrase : elle illustre l'idée précédente par un cas concret $\rightarrow$ \esp{por ejemplo}.
\item Cinquième phrase : c'est le dernier argument, qui annonce la conclusion $\rightarrow$ \esp{Por último}.
\end{itemize}

Texte complété : \esp{En primer lugar, la convivencia es posible si nos respetamos. Además, el diálogo resuelve los conflictos sin violencia. Sin embargo, todavía existe la discriminación en algunas escuelas; por ejemplo, hay estudiantes que rechazan a los compañeros de otras regiones. Por último, debemos educar en la tolerancia desde la infancia.}

« En premier lieu, le vivre-ensemble est possible si nous nous respectons. De plus, le dialogue résout les conflits sans violence. Cependant, la discrimination existe encore dans certaines écoles ; par exemple, il y a des élèves qui rejettent les camarades d'autres régions. Enfin, nous devons éduquer à la tolérance dès l'enfance. »
\end{exoresolu}

\begin{exercice}[À toi de jouer]
\begin{enumerate}
\item Traduis en espagnol : « D'abord, le respect est indispensable. De plus, il rend la vie en communauté plus facile. Cependant, certaines personnes restent intolérantes. En conclusion, nous devons dialoguer. »
\item Relie les deux phrases avec \esp{sin embargo} et un point-virgule : \esp{Somos diferentes. / Somos iguales en derechos.}
\item Corrige l'erreur : \esp{En facto, la tolerancia es difícil.} et \esp{Finalmente, mi último argumento es el respeto.}
\end{enumerate}
\end{exercice}

\begin{corrige}[Exercice « À toi de jouer »]
\begin{enumerate}
\item \esp{En primer lugar, el respeto es indispensable. Además, hace más fácil la vida en comunidad. Sin embargo, algunas personas siguen siendo intolerantes. En conclusión, debemos dialogar.}
\item \esp{Somos diferentes; sin embargo, somos iguales en derechos.}
\item \esp{En realidad, la tolerancia es difícil.} (« en fait » = \esp{en realidad}, jamais \esp{en facto}) ; \esp{Por último, mi argumento final es el respeto.} (dernier argument d'un plan = \esp{por último}, pas \esp{finalmente}).
\end{enumerate}
\end{corrige}

\begin{aretenir}[L'essentiel des connecteurs logiques]
\begin{itemize}
\item \textbf{Ordonner} : \esp{en primer lugar}, \esp{en segundo lugar}, \esp{luego}, \esp{por último}.
\item \textbf{Ajouter} : \esp{además}, \esp{también}, \esp{incluso}.
\item \textbf{Opposer} : \esp{sin embargo}, \esp{pero}, \esp{no obstante}, \esp{en cambio}.
\item \textbf{Conclure} : \esp{en conclusión}, \esp{en resumen}, \esp{por eso}, \esp{por lo tanto}.
\item \textbf{Exemplifier} : \esp{por ejemplo}, \esp{es decir}.
\item Connecteur en tête de phrase $\rightarrow$ virgule. \esp{Sin embargo} accepte le point-virgule devant lui.
\item Pièges : « en fait » = \esp{en realidad} ; dernier argument = \esp{por último} ; \esp{sin embargo} est un bloc figé qui signifie « cependant ».
\end{itemize}
\end{aretenir}

Ces connecteurs sont désormais tes outils. Dans la section suivante, tu les réutiliseras pour traduire des phrases sur l'identité et la tolérance, puis pour construire ton propre texte argumenté sur le vivre-ensemble.

\coursec{La identidad, la tolerancia y la convivencia (continuación) : débat et rédaction}

\courssub{Texte support : « Voces del aula »}

\begin{textebox}[Voces del aula — Débat au lycée bilingue de Yaoundé]
— \textbf{En primer lugar}, quiero decir que la convivencia no es un sueño, es una tarea diaria. \textbf{Además}, \textbf{pienso que} el respeto mutuo es la base de todo. Si no respetamos a \textbf{los demás}, no hay paz posible.

— Estoy de acuerdo, \textbf{sin embargo} \textbf{no creo que} sea fácil cambiar las mentalidades. \textbf{Hay que} educar desde la infancia. \textbf{Debemos} enseñar a los niños que \textbf{cada cultura tiene sus propias} riquezas. Los prejuicios nacen de la ignorancia.

— Exacto. \textbf{Por último}, quiero añadir que la discriminación duele. Yo he visto compañeros marginados por su acento o su origen. \textbf{No opino que} eso \textbf{sea} normal. \textbf{Hay que} luchar \textbf{contra} la exclusión, no solo con palabras, sino con hechos.

— \textbf{En resumen}, la tolerancia no es pasividad. Es actuar. \textbf{Debemos} construir puentes, no muros. \textbf{¿Están de acuerdo?}
\end{textebox}

\emph{Traduction libre :}
« En premier lieu, je veux dire que le vivre-ensemble n'est pas un rêve, c'est une tâche quotidienne. De plus, je pense que le respect mutuel est la base de tout. Si nous ne respectons pas les autres, pas de paix possible. — Je suis d'accord, cependant je ne crois pas que ce soit facile de changer les mentalités. Il faut éduquer dès l'enfance. Nous devons apprendre aux enfants que chaque culture a ses propres richesses. Les préjugés naissent de l'ignorance. — Exact. Pour finir, je veux ajouter que la discrimination fait mal. J'ai vu des camarades marginalisés à cause de leur accent ou de leur origine. Je ne pense pas que cela soit normal. Il faut lutter contre l'exclusion, pas seulement avec des mots, mais avec des actes. — En résumé, la tolérance n'est pas de la passivité. C'est agir. Nous devons construire des ponts, pas des murs. Êtes-vous d'accord ? »

\begin{definition}[Glossaire du texte : vocabulaire clé du vivre-ensemble]
\begin{itemize}
    \item \esp{La convivencia} : le vivre-ensemble, la cohabitation.
    \item \esp{Una tarea diaria} : une tâche quotidienne.
    \item \esp{El respeto mutuo} : le respect mutuel.
    \item \esp{Los demás} : les autres (l'humanité, le prochain).
    \item \esp{Cambiar las mentalidades} : changer les mentalités.
    \item \esp{Desde la infancia} : dès l'enfance.
    \item \esp{Propias riquezas / propias costumbres} : propres richesses / propres coutumes (adj. \emph{propio} accordé).
    \item \esp{Marginado / marginada} : marginalisé(e), mis(e) à l'écart.
    \item \esp{El acento} : l'accent (façon de parler).
    \item \esp{La exclusión} : l'exclusion.
    \item \esp{La ignorancia} : l'ignorance.
    \item \esp{Construir puentes, no muros} : construire des ponts, pas des murs (métaphore du dialogue).
    \item \esp{La pasividad} : la passivité.
    \item \esp{Estar de acuerdo} : être d'accord (ne pas oublier \esp{de}).
\end{itemize}
\end{definition}

\courssub{Lexique thématique : débat et vivre-ensemble}

\begin{propriete}[Mots et expressions utiles pour argumenter]
\begin{center}
\begin{tabularx}{\linewidth}{|X|X|X|}
\hline
\rowcolor{popblueL} \textbf{Français} & \textbf{Espagnol} & \textbf{Catégorie / Note} \\ \hline
L'identité & \esp{La identidad} & Nom fém. \\ \hline
La tolérance & \esp{La tolerancia} & Nom fém. \\ \hline
La discrimination & \esp{La discriminación} & Nom fém. \\ \hline
Le préjugé & \esp{El prejuicio} & Nom masc. \\ \hline
Le stéréotype & \esp{El estereotipo} & Nom masc. \\ \hline
L'exclusion & \esp{La exclusión} & Nom fém. \\ \hline
Le respect & \esp{El respeto} & Nom masc. \\ \hline
La diversité & \esp{La diversidad} & Nom fém. \\ \hline
L'égalité & \esp{La igualdad} & Nom fém. \\ \hline
La justice / L'injustice & \esp{La justicia / La injusticia} & Nom fém. \\ \hline
Respecter & \esp{Respetar} & Verbe 1er groupe \\ \hline
Accepter & \esp{Aceptar} & Verbe 1er groupe \\ \hline
Dialoguer & \esp{Dialogar} & Verbe 1er groupe \\ \hline
Lutter contre & \esp{Luchar contra} & \esp{contra} = contre \\ \hline
Éduquer & \esp{Educar} & Verbe 1er groupe \\ \hline
Changer & \esp{Cambiar} & Verbe 1er groupe \\ \hline
Enrichir & \esp{Enriquecer} & Verbe 2e groupe (c$\rightarrow$zc) \\ \hline
Diviser & \esp{Dividir} & Verbe 3e groupe \\ \hline
Empêcher & \esp{Impedir} & Verbe 3e groupe (e$\rightarrow$i) \\ \hline
Être d'accord & \esp{Estar de acuerdo con} & \textbf{Jamais} \esp{*estar acuerdo} \\ \hline
Penser que / Croire que & \esp{Pensar que / Creer que / Opinar que} & + Indicatif (aff.) / Subjonctif (nég.) \\ \hline
Devoir (moral) & \esp{Deber} + inf. & \esp{Debemos} = Nous devons \\ \hline
Il faut (général) & \esp{Hay que} + inf. & Impersonnel \\ \hline
Propre (à soi) & \esp{Propio / propia / propios / propias} & Accord genre/nombre \\ \hline
Les autres (prochain) & \esp{Los demás} & Invariable, avec article \\ \hline
Cependant & \esp{Sin embargo} & Connecteur adversatif \\ \hline
De plus / En outre & \esp{Además} & Connecteur additif \\ \hline
En premier lieu & \esp{En primer lugar} & Connecteur d'ordre \\ \hline
Pour finir / Enfin & \esp{Por último} & Connecteur d'ordre \\ \hline
En résumé / En conclusion & \esp{En resumen / En conclusión} & Connecteur de synthèse \\ \hline
\end{tabularx}
\end{center}
\end{propriete}

\courssub{Grammaire : les connecteurs logiques de l'argumentation}

\begin{propriete}[Les connecteurs essentiels pour structurer un discours]
Pour obtenir une note élevée à l'expression écrite et orale (Bac), vous devez articuler vos idées avec des connecteurs variés, bien placés et suivis de la ponctuation correcte (virgule après le connecteur en début de phrase).
\begin{center}
\begin{tabularx}{\linewidth}{|l|X|X|}
\hline
\rowcolor{popblueL} \textbf{Connecteur} & \textbf{Fonction} & \textbf{Exemple d'emploi} \\ \hline
\esp{En primer lugar} & Introduire le 1er argument & \esp{En primer lugar, la tolerancia es un derecho.} \\ \hline
\esp{Además} / \esp{También} / \esp{Por añadidura} & Ajouter un argument (addition) & \esp{Además, enriquece a la sociedad.} \\ \hline
\esp{Sin embargo} / \esp{No obstante} / \esp{Pero} & Opposer / Nuancer (adversatif) & \esp{Sin embargo, no es fácil de lograr.} \\ \hline
\esp{Por último} / \esp{Finalmente} & Introduire le dernier argument (ordre) & \esp{Por último, hay que actuar.} \\ \hline
\esp{En resumen} / \esp{En conclusión} / \esp{Por tanto} & Conclure / Résumer (conséquence) & \esp{En resumen, debemos construir puentes.} \\ \hline
\esp{Por un lado... por otro lado} & Structurer en deux pôles (balançoire) & \esp{Por un lado hay avances, por otro lado persisten retos.} \\ \hline
\end{tabularx}
\end{center}
\end{propriete}

\begin{attention}[Ponctuation des connecteurs]
En espagnol, les connecteurs en début de phrase sont \textbf{toujours suivis d'une virgule} : \esp{En primer lugar, ...} ; \esp{Sin embargo, ...} ; \esp{Por último, ...}. Oublier la virgule est une faute d'orthographe sanctionnée.
\end{attention}


\courssub{Grammaire : opinion, obligation et adjectif \emph{propio}}

\courssub{Observation : indicatif ou subjonctif après un verbe d'opinion ?}

Comparez ces deux phrases issues du texte \emph{Voces del aula} :

\begin{textebox}[Extraits du débat]
— \textbf{Pienso que el respeto mutuo es la base de todo.} \\
— \textbf{No creo que} sea fácil cambiar las mentalidades. \\
— \textbf{No opino que} eso \textbf{sea} normal.
\end{textebox}

\emph{Traduction : « Je pense que le respect mutuel est la base de tout. / Je ne crois pas que ce soit facile de changer les mentalités. / Je ne pense pas que cela soit normal. »}

\begin{propriete}[Règle : \emph{Creer / Pensar / Opinar / Considerar / Suponer} + \emph{que} + Indicatif ou Subjonctif]
\begin{itemize}
    \item \textbf{Forme affirmative} (\esp{creo que, pienso que, opino que, considero que, supongo que}) $\rightarrow$ \textbf{Indicatif} (le fait est présenté comme certain, réel, objectif).
    \item \textbf{Forme négative} (\esp{no creo que, no pienso que, no opino que, no considero que, no supongo que}) $\rightarrow$ \textbf{Subjonctif} (doute, subjectivité, négation de la certitude).
\end{itemize}
\textbf{Exemples traduits} :
\esp{Creo que la igualdad es un derecho fundamental.} « Je crois que l'égalité est un droit fondamental. » (Indicatif \emph{es}) \\
\esp{No creo que la igualdad sea fácil de alcanzar.} « Je ne crois pas que l'égalité soit facile à atteindre. » (Subjonctif \emph{sea}) \\
\esp{Pienso que el diálogo resuelve los conflictos.} « Je pense que le dialogue résout les conflits. » (Indicatif \emph{resuelve}) \\
\esp{No pienso que la violencia se justifique.} « Je ne pense pas que la violence se justifie. » (Subjonctif \emph{se justifique}) \\
\esp{Considero que la diversidad enriquece.} « Je considère que la diversité enrichit. » (Indicatif \emph{enriquece}) \\
\esp{No considero que la discriminación sea normal.} « Je ne considère pas que la discrimination soit normale. » (Subjonctif \emph{sea})
\end{propriete}

\begin{attention}[Piège « Bac » : la négation change le mode]
En français, « je ne pense pas que » garde souvent l'indicatif (ou le conditionnel). En espagnol, \textbf{la négation impose systématiquement le subjonctif présent} pour ces verbes d'opinion. C'est un point classique du Baccalauréat.
\begin{itemize}
    \item \esp{Creo que es justo.} (Indicatif \emph{es})
    \item \esp{No creo que \textbf{sea} justo.} (Subjonctif \emph{sea}, \textbf{pas} \emph{es})
\end{itemize}
\emph{Autres verbes concernés :} \esp{suponer que, imaginar que, parecer que} (affirmatif = indicatif ; négatif = subjonctif).
\end{attention}

\courssub{Thème 1 : Exprimer l'obligation morale — « Nous devons »}

\begin{exemplebox}[Thème 1]
\textbf{Français :} « Nous devons respecter les différences. » \\
\textbf{Espagnol :} \esp{Debemos respetar las diferencias.}
\end{exemplebox}

\begin{propriete}[\esp{Deber} + infinitif = obligation morale, devoir, conseil fort]
\esp{Deber} (conjugué au présent de l'indicatif) + infinitif exprime un devoir moral, une obligation interne, une recommandation forte.
\begin{center}
\begin{tabular}{|l|l|l|}
\hline
\textbf{Personne} & \textbf{\esp{Deber} (Présent Indicatif)} & \textbf{Exemple} \\ \hline
yo & debo & \esp{Debo escuchar a los demás.} \\ \hline
tú & debes & \esp{Debes ser tolerante.} \\ \hline
él/ella/usted & debe & \esp{Él debe cambiar de actitud.} \\ \hline
nosotros/as & \textbf{debemos} & \esp{\textbf{Debemos respetar} las diferencias.} \\ \hline
vosotros/as & debéis & \esp{Debéis dialogar más.} \\ \hline
ellos/ellas/ustedes & deben & \esp{Deben luchar contra la injusticia.} \\ \hline
\end{tabular}
\end{center}
\textbf{Traduction des exemples} : « Je dois écouter les autres / Tu dois être tolérant / Il doit changer d'attitude / \textbf{Nous devons respecter les différences} / Vous devez dialoguer plus / Ils doivent lutter contre l'injustice. »
\end{propriete}

\begin{attention}[Erreur fréquente : \esp{Deber} vs \esp{Tener que}]
\begin{itemize}
    \item \textbf{\esp{Deber} + inf.} = Obligation \textbf{morale, interne, devoir} (« Je dois » par conscience). Ex. : \esp{Debo ayudar a mi prójimo.}
    \item \textbf{\esp{Tener que} + inf.} = Obligation \textbf{extérieure, circonstancielle, contrainte} (« Avoir à », « être obligé de »). Ex. : \esp{Tengo que ir al médico mañana.} (J'ai rendez-vous).
\end{itemize}
Pour « nous devons respecter les différences » (devoir moral, éthique), on utilise \textbf{\esp{debemos}}.
\end{attention}

\courssub{Thème 2 : Exprimer l'obligation impersonnelle — « Il faut »}

\begin{exemplebox}[Thème 2]
\textbf{Français :} « Il faut lutter contre la discrimination. » \\
\textbf{Espagnol :} \esp{Hay que luchar contra la discriminación.}
\end{exemplebox}

\begin{propriete}[\esp{Hay que} + infinitif = « il faut », obligation générale et impersonnelle]
\esp{Hay que} (forme figée du verbe \esp{haber} à l'infinitif présent, 3\textsuperscript{e} personne du singulier) s'emploie \textbf{uniquement à la 3\textsuperscript{e} personne du singulier}. Il n'y a \textbf{pas de sujet}. C'est la structure standard pour traduire « il faut » au niveau Première A.
\begin{center}
\begin{tabular}{|l|l|}
\hline
\textbf{Structure} & \textbf{Exemple} \\ \hline
\esp{Hay que} + infinitif & \esp{Hay que respetar a todos.} \\ \hline
 & \esp{Hay que luchar contra la discriminación.} \\ \hline
 & \esp{Hay que educar en la tolerancia.} \\ \hline
 & \esp{Hay que promover la igualdad.} \\ \hline
\end{tabular}
\end{center}
\emph{Traduction :} « Il faut respecter tout le monde / Il faut lutter contre la discrimination / Il faut éduquer à la tolérance / Il faut promouvoir l'égalité. »
\end{propriete}

\begin{attention}[Pièges à éviter avec « il faut »]
\begin{enumerate}
    \item \textbf{Ne pas écrire} \esp{*Es necesario que} + subjonctif (structure correcte mais niveau B1/B2, hors programme pour l'obligation simple en 1ère A).
    \item \textbf{Ne pas écrire} \esp{*Hay que} + subjonctif (jamais : \esp{hay que} appelle \textbf{toujours l'infinitif}).
    \item \textbf{Ne pas confondre} \esp{hay que} (il faut) et \esp{hay} (il y a) :
    \begin{itemize}
        \item \esp{Hay muchos prejuicios} = « Il y a beaucoup de préjugés » (verbe \esp{hay} = il y a).
        \item \esp{Hay que combatir los prejuicios} = « Il faut combattre les préjugés » (\esp{hay que} = il faut).
    \end{itemize}
\end{enumerate}
\end{attention}

\courssub{Thème 3 : L'adjectif \emph{propio} — « propre »}

\begin{exemplebox}[Thème 3]
\textbf{Français :} « Chaque culture a ses propres coutumes. » \\
\textbf{Espagnol :} \esp{Cada cultura tiene sus propias costumbres.}
\end{exemplebox}

\begin{propriete}[Accord de \emph{propio / propia / propios / propias}]
\emph{Propio} est un adjectif qualificatif qui s'accorde en \textbf{genre et en nombre} avec le nom qu'il qualifie. Il se place \textbf{généralement après le nom} (mais peut aller devant pour insister sur l'appartenance exclusive).
\begin{center}
\begin{tabular}{|l|l|l|l|}
\hline
\textbf{Masculin singulier} & \textbf{Féminin singulier} & \textbf{Masculin pluriel} & \textbf{Féminin pluriel} \\ \hline
\esp{propio} & \esp{propia} & \esp{propios} & \esp{propias} \\ \hline
\end{tabular}
\end{center}
\textbf{Exemples} :
\esp{Cada persona tiene su \textbf{propia} identidad.} (féminin singulier : \emph{identidad}) \\
\esp{Todos los pueblos tienen sus \textbf{propias} tradiciones.} (féminin pluriel : \emph{tradiciones}) \\
\esp{Mi \textbf{propio} hermano no me entiende.} (masculin singulier, placé avant pour insister : « Mon \emph{propre} frère… ») \\
\esp{Cada estudiante tiene su \textbf{propio} cuaderno.} (masculin singulier : \emph{cuaderno})
\end{propriete}

\courssub{Traduction : phrases sur l'identité et la tolérance}

\courssub{Version 1 : De l'espagnol au français}

\begin{exemplebox}[Version 1]
\textbf{Espagnol :} \esp{La convivencia es posible si respetamos a los demás.} \\
\textbf{Français :} « Le vivre-ensemble est possible si nous respectons les autres. »
\end{exemplebox}

\begin{propriete}[Analyse grammaticale]
\begin{itemize}
    \item \esp{La convivencia} = le vivre-ensemble / la cohabitation (nom féminin, sujet).
    \item \esp{es posible} = est possible (\esp{ser} + adjectif, indicatif présent, vérité générale).
    \item \esp{si respetamos} = si nous respectons (\esp{si} + indicatif présent pour une condition réelle/possible ; \esp{respetamos} = 1\textsuperscript{re} personne du pluriel du présent de l'indicatif de \esp{respetar}).
    \item \esp{a los demás} = les autres (préposition \esp{a} + article \esp{los} + pronom \esp{demás} ; \esp{respetar a alguien}).
\end{itemize}
\end{propriete}

\begin{attention}[Piège majeur : \emph{los demás} vs \emph{los otros}]
\begin{itemize}
    \item \textbf{\esp{Los demás}} = « les autres » (dans le sens de « le reste des gens », « les autres personnes » de manière générale, l'humanité). \textbf{Invariable en genre} (toujours \emph{los demás} même pour des femmes), prend toujours l'article défini \emph{los}.
    \item \textbf{\esp{Los otros / las otras}} = « les autres » (dans le sens de « les autres \emph{choses/personnes} distinctes de celles déjà citées », avec accord de genre). Ex. : \esp{Unos alumnos estudian, los otros juegan.} « Certains élèves étudient, les autres jouent. »
\end{itemize}
\textbf{Règle pratique :} Pour « respecter les autres » (le prochain, l'altérité), toujours \esp{respetar a \textbf{los demás}}.
\end{attention}

\courssub{Version 2 : Connecteur adversatif et sujet collectif}

\begin{exemplebox}[Version 2]
\textbf{Espagnol :} \esp{Sin embargo, los prejuicios dividen a la sociedad.} \\
\textbf{Français :} « Cependant, les préjugés divisent la société. »
\end{exemplebox}

\begin{propriete}[Analyse]
\begin{itemize}
    \item \esp{Sin embargo} = cependant, nevertheless (connecteur adversatif, suivi d'une virgule).
    \item \esp{Los prejuicios} = les préjugés (sujet, masculin pluriel).
    \item \esp{Dividen} = divisent (verbe \esp{dividir}, indicatif présent, 3\textsuperscript{e} personne du pluriel).
    \item \esp{A la sociedad} = la société (COI introduit par \esp{a} car \esp{dividir} quelqu'un = \esp{dividir a alguien}).
\end{itemize}
\end{propriete}

\begin{savaistu}[\esp{Prejuicio} : étymologie]
\esp{Prejuicio} vient du latin \emph{praeiudicium} (\emph{prae-} « avant » + \emph{iudicium} « jugement »). Un préjugé est donc un « jugement porté \textbf{avant} de connaître la réalité ». En français, \emph{préjugé} a la même étymologie. Lutter contre les préjugés, c'est refuser de juger sans connaître.
\end{savaistu}

\courssub{Récapitulatif : correspondances clés français / espagnol}

\begin{center}
\begin{tabularx}{\linewidth}{|X|X|X|}
\hline
\rowcolor{popblueL} \textbf{Français} & \textbf{Espagnol} & \textbf{Remarque} \\ \hline
Il faut + infinitif & \esp{Hay que} + infinitif & Impersonnel, obligation générale \\ \hline
Nous devons + infinitif & \esp{Debemos} + infinitif & Obligation morale, 1\textsuperscript{re} pers. pl. \\ \hline
Les autres (prochain, humanité) & \esp{Los demás} & Invariable, avec article défini \\ \hline
Les autres (les distincts) & \esp{Los otros / las otras} & Accord en genre et nombre \\ \hline
Être d'accord avec & \esp{Estar de acuerdo con} & \textbf{Jamais} \esp{*estar acuerdo} (sans \esp{de}) \\ \hline
Lutter contre & \esp{Luchar contra} & \esp{Contra} = contre (opposition) \\ \hline
Les différences & \esp{Las diferencias} & Féminin pluriel \\ \hline
Propre (à soi) & \esp{Propio / propia / propios / propias} & Accord avec le nom qualifié \\ \hline
Penser que (affirmatif) & \esp{Pienso/Creo que} + \textbf{Indicatif} & Certitude \\ \hline
Ne pas penser que & \esp{No pienso/creo que} + \textbf{Subjonctif} & Doute / Négation \\ \hline
\end{tabularx}
\end{center}

\begin{popfigure}
\begin{tikzpicture}[
    mindmap,
    concept color=popblue,
    level 1/.style={sibling angle=90, level distance=4.5cm},
    level 2/.style={level distance=3.5cm},
    every node/.style={concept, execute at begin node=\small, align=center},
    grow cyclic
]
\node[align=center]{Identité,\\Tolérance,\\Convivencia}
    child[concept color=poporange] {node {Obligation}
        child {node[align=center]{\esp{Hay que} + inf.\\« Il faut »}}
        child {node[align=center]{\esp{Debemos} + inf.\\« Nous devons »}}
    }
    child[concept color=popgreen] {node {Opinion}
        child {node[align=center]{\esp{Creo/Pienso que}\\+ \textbf{Indicatif}}}
        child {node[align=center]{\esp{No creo/pienso que}\\+ \textbf{Subjonctif}}}
    }
    child[concept color=poppurple] {node {Vocabulaire clé}
        child {node[align=center]{\esp{Los demás}\\« Les autres »}}
        child {node[align=center]{\esp{Luchar contra}\\« Lutter contre »}}
        child {node[align=center]{\esp{Propias costumbres}\\« Propres coutumes »}}
    }
    child[concept color=poppink] {node {Connecteurs}
        child {node {\esp{En primer lugar}}}
        child {node {\esp{Además}}}
        child {node {\esp{Sin embargo}}}
        child {node {\esp{Por último}}}
        child {node {\esp{En resumen}}}
    };
\end{tikzpicture}
\legende{Carte mentale : structures grammaticales et lexique clé de la leçon EA12}
\end{popfigure}

\courssub{Méthode pour le thème (français $\rightarrow$ espagnol)}

\begin{methode}[Traduire une phrase d'obligation ou d'opinion]
\textbf{Étape 1 :} Identifier la structure française.
\begin{itemize}
    \item « Il faut + infinitif » $\rightarrow$ \esp{Hay que} + infinitif.
    \item « Nous devons + infinitif » $\rightarrow$ \esp{Debemos} + infinitif.
    \item « Je pense que… » (affirmatif) $\rightarrow$ \esp{Pienso/Creo que} + \textbf{indicatif}.
    \item « Je ne pense pas que… » $\rightarrow$ \esp{No pienso/creo que} + \textbf{subjonctif}.
\end{itemize}
\textbf{Étape 2 :} Traduire le vocabulaire clé en vérifiant les pièges.
\begin{itemize}
    \item « les autres » $\rightarrow$ \esp{los demás} (pas \emph{los otros}).
    \item « lutter contre » $\rightarrow$ \esp{luchar contra}.
    \item « les différences » $\rightarrow$ \esp{las diferencias}.
    \item « propre » $\rightarrow$ \esp{propio/propia/propias/propios} (accord genre/nombre).
    \item « être d'accord » $\rightarrow$ \esp{estar de acuerdo con}.
\end{itemize}
\textbf{Étape 3 :} Construire la phrase espagnole en respectant l'ordre des mots (SVO) et les accords (adjectifs, participes).
\textbf{Étape 4 :} Relire : vérifiez les accents (\esp{acción, nación, convivencia, tolerancia, exclusión, única, propia}), la ponctuation espagnole (¿ ? ¡ !) et les accords.
\end{methode}

\courssub{Exemples traduits variés (thème et version)}

\begin{exemplebox}[Série thème — Français $\rightarrow$ Espagnol]
\begin{enumerate}
    \item « Il faut dialoguer pour vivre en paix. » $\rightarrow$ \esp{Hay que dialogar para vivir en paz.}
    \item « Nous devons accepter la diversité. » $\rightarrow$ \esp{Debemos aceptar la diversidad.}
    \item « Je pense que l'égalité est un droit. » $\rightarrow$ \esp{Pienso que la igualdad es un derecho.} (indicatif \emph{es})
    \item « Je ne crois pas que le racisme disparaisse vite. » $\rightarrow$ \esp{No creo que el racismo \textbf{desaparezca} pronto.} (subjonctif \emph{desaparezca})
    \item « Chaque pays a sa propre histoire. » $\rightarrow$ \esp{Cada país tiene su \textbf{propia} historia.} (\emph{historia} fém. sg.)
    \item « Cependant, les préjugés persistent. » $\rightarrow$ \esp{Sin embargo, los prejuicios persisten.}
\end{enumerate}
\end{exemplebox}

\begin{exemplebox}[Série version — Espagnol $\rightarrow$ Français]
\begin{enumerate}
    \item \esp{La tolerancia es la base de la convivencia.} $\rightarrow$ « La tolérance est la base du vivre-ensemble. »
    \item \esp{Hay que educar a los niños en el respeto.} $\rightarrow$ « Il faut éduquer les enfants au respect. »
    \item \esp{No opino que la discriminación sea normal.} $\rightarrow$ « Je ne pense pas que la discrimination soit normale. » (subjonctif \emph{sea})
    \item \esp{Cada cultura tiene sus propias tradiciones.} $\rightarrow$ « Chaque culture a ses propres traditions. »
    \item \esp{Debemos construir una sociedad más justa.} $\rightarrow$ « Nous devons construire une société plus juste. »
    \item \esp{Los demás no son nuestros enemigos.} $\rightarrow$ « Les autres ne sont pas nos ennemis. »
\end{enumerate}
\end{exemplebox}

\courssub{Exercice résolu}

\begin{exoresolu}[Thème et version mélangés]
\textbf{Énoncé :} Traduisez les phrases suivantes.
\begin{enumerate}
    \item Il faut respecter les opinions de los demás.
    \item Nous devons lutter contre l'intolérance.
    \item Je pense que la diversité enrichit la société.
    \item Je ne crois pas que l'exclusion soit une solution.
    \item \esp{La identidad de cada persona es única.}
    \item \esp{Sin embargo, hay que trabajar por la paz.}
\end{enumerate}

\tcblower
\textbf{Solution détaillée :}

1. \esp{Hay que respetar las opiniones de los demás.}
\begin{itemize}
    \item « Il faut » $\rightarrow$ \esp{Hay que} + infinitif.
    \item « respecter » $\rightarrow$ \esp{respetar}.
    \item « les opinions » $\rightarrow$ \esp{las opiniones} (fém. pl.).
    \item « de les autres » $\rightarrow$ \esp{de los demás} (piège : \emph{los demás}, pas \emph{los otros}).
\end{itemize}

2. \esp{Debemos luchar contra la intolerancia.}
\begin{itemize}
    \item « Nous devons » $\rightarrow$ \esp{Debemos} (1\textsuperscript{re} pl. de \esp{deber}).
    \item « lutter contre » $\rightarrow$ \esp{luchar contra} (pas \emph{luchar *contra de*} ni \emph{luchar *en contra*} seul).
    \item « l'intolérance » $\rightarrow$ \esp{la intolerancia} (fém. sg.).
\end{itemize}

3. \esp{Pienso que la diversidad enriquece la sociedad.}
\begin{itemize}
    \item « Je pense que » affirmatif $\rightarrow$ \esp{Pienso que} + \textbf{indicatif}.
    \item « la diversité » $\rightarrow$ \esp{la diversidad}.
    \item « enrichit » $\rightarrow$ \esp{enriquece} (3\textsuperscript{e} sg. présent indicatif de \esp{enriquecer}, c$\rightarrow$zc).
    \item « la société » $\rightarrow$ \esp{la sociedad}.
\end{itemize}

4. \esp{No creo que la exclusión \textbf{sea} una solución.}
\begin{itemize}
    \item « Je ne crois pas que » $\rightarrow$ \esp{No creo que} + \textbf{subjonctif}.
    \item « l'exclusion » $\rightarrow$ \esp{la exclusión} (sujet de la subordonnée).
    \item « soit » $\rightarrow$ \esp{sea} (subjonctif présent 3\textsuperscript{e} sg. de \esp{ser}).
    \item « une solution » $\rightarrow$ \esp{una solución}.
\end{itemize}

5. « L'identité de chaque personne est unique. »
\begin{itemize}
    \item \esp{La identidad} = l'identité.
    \item \esp{de cada persona} = de chaque personne.
    \item \esp{es única} = est unique (accord fém. sg. avec \emph{identidad}).
\end{itemize}

6. « Cependant, il faut travailler pour la paix. »
\begin{itemize}
    \item \esp{Sin embargo} = cependant.
    \item \esp{hay que trabajar} = il faut travailler.
    \item \esp{por la paz} = pour la paix (\esp{por} = cause, motif, « pour » au sens de « en faveur de »).
\end{itemize}
\end{exoresolu}

\courssub{Méthode du débat oral}

\begin{experience}[Activité orale : « Le tribunal de la tolérance »]
\textbf{Objectif :} Réutiliser \esp{hay que}, \esp{debemos}, \esp{no creo que} + subjonctif, connecteurs (\esp{en primer lugar, además, sin embargo, por último}), \esp{los demás}, \esp{propio}.
\textbf{Déroulement (15--20 min) :}
\begin{enumerate}
    \item Par groupes de 4 : un « juge », un « procureur », un « avocat », un « accusé » (un préjugé : ex. « Les gens du Nord sont paresseux », « Les filles ne savent pas jouer au football », « Les étrangers volent notre travail »).
    \item Le procureur accuse (obligation) : \esp{« Hay que eliminar este prejuicio porque divide a los demás. »}
    \item L'avocat défend (ironiquement, opinion négative) : \esp{« No creo que sea tan grave, cada quien tiene su propia opinión. »}
    \item Le juge tranche avec connecteurs : \esp{« En primer lugar, el prejuicio es dañino. Además, frena el desarrollo. Sin embargo, la educación puede cambiarlo. Por último, declaro que... »}
    \item Rotation des rôles. Le professeur note les structures cibles au tableau et corrige la prononciation (\esp{con-vi-ven-cia, to-le-ran-cia, pre-hui-cio}).
\end{enumerate}
\end{experience}

\courssub{Production écrite guidée : texte argumenté (100 mots)}

\begin{methode}[Rédiger un court texte argumentatif (environ 100 mots)]
\textbf{Consigne type (Bac) :} « \esp{Escribe un texto de unos 100 palabras sobre: « La tolerancia: ¿fuerza o debilidad? » Utiliza conectores, \emph{hay que}, \emph{debemos}, \emph{no creo que} + subjuntivo, \emph{los demás}, \emph{propio}. »}

\textbf{Étape 1 : Analyser le sujet et choisir sa thèse.}
Ex. : Thèse = La tolérance est une force (exige du courage, construit la paix).

\textbf{Étape 2 : Structurer en 3 paragraphes courts (Introduction, Développement, Conclusion).}
\begin{itemize}
    \item \textbf{Intro (1 phrase + thèse) :} \esp{En primer lugar, opino que la tolerancia no es debilidad, sino una gran fuerza.}
    \item \textbf{Développement (2--3 arguments + exemples + connecteurs) :}
    \esp{Además, hay que tener valor para aceptar las diferencias.} (Argument 1 : courage)
    \esp{Sin embargo, no creo que sea fácil, porque los prejuicios son fuertes.} (Nuance / Concession)
    \esp{Por último, debemos educar en el respeto a los demás para construir puentes.} (Argument 2 : éducation / action)
    \item \textbf{Conclusion (synthèse + ouverture) :} \esp{En resumen, la tolerancia es el camino hacia la convivencia pacífica. ¿Están de acuerdo?}
\end{itemize}

\textbf{Étape 3 : Compter les mots (objectif 90--110 mots).}
\textbf{Étape 4 : Relire (Check-list) :}
\begin{itemize}
    \item Accords (genre/nombre) : \esp{propias costumbres, propia identidad} ?
    \item Subjonctif après négation : \esp{no creo que sea} (pas \emph{es}) ?
    \item \esp{Hay que} + infinitif / \esp{Debemos} + infinitif ?
    \item \esp{Los demás} (pas \emph{los otros}) ?
    \item Connecteurs + virgules : \esp{En primer lugar, Además, Sin embargo, Por último, En resumen} ?
    \item Accents : \esp{convivencia, tolerancia, exclusión, única, también, además} ?
    \item Ponctuation : ¿ ? ¡ ! ?
\end{itemize}
\end{methode}

\begin{exemplebox}[Modèle de rédaction (environ 100 mots)]
\textbf{Sujet :} \esp{« La tolerancia: ¿fuerza o debilidad? »}

\esp{En primer lugar, pienso que la tolerancia es una fuerza, no una debilidad. Además, hay que tener mucho valor para aceptar a los demás con sus propias identidades. Sin embargo, no creo que sea un camino fácil, porque los prejuicios y los estereotipos dividen a la sociedad. Por último, debemos educar desde la infancia en el respeto mutuo y luchar contra la discriminación. En resumen, la tolerancia es la base de la convivencia pacífica y la única forma de construir un mundo más justo.}
\end{exemplebox}

\courssub{Exercices d'entraînement}

\begin{exercice}[Exercice 1 — Thème : traduisez en espagnol]
\begin{enumerate}
    \item Il faut accepter les différences culturelles.
    \item Nous devons construire un monde sans discrimination.
    \item Je pense que le dialogue résout les conflits.
    \item Je ne pense pas que l'injustice soit nécessaire.
    \item Chaque élève a sa propre identité.
    \item Cependant, les stéréotypes divisent les gens.
\end{enumerate}
\end{exercice}

\begin{exercice}[Exercice 2 — Version : traduisez en français]
\begin{enumerate}
    \item \esp{Hay que promover la tolerancia en la escuela.}
    \item \esp{Debemos escuchar a los demás con atención.}
    \item \esp{Creo que la convivencia pacífica es posible.}
    \item \esp{No opinamos que la violencia sea justificable.}
    \item \esp{Cada comunidad tiene sus propias costumbres.}
    \item \esp{Los prejuicios impiden el diálogo sincero.}
\end{enumerate}
\end{exercice}

\begin{exercice}[Exercice 3 — Production écrite]
Rédigez en espagnol un texte de \textbf{100 mots environ} sur le sujet suivant :
\esp{« ¿Es posible la convivencia sin respeto? »}
Utilisez obligatoirement : \esp{hay que}, \esp{debemos}, un connecteur d'addition, un connecteur adversatif, \esp{no creo que} + subjonctif, \esp{los demás}, \esp{propio}.
\end{exercice}

\begin{corrige}[Exercice 1 — Thème]
\begin{enumerate}
    \item \esp{Hay que aceptar las diferencias culturales.}
    \item \esp{Debemos construir un mundo sin discriminación.}
    \item \esp{Pienso que el diálogo resuelve los conflictos.} (indicatif \emph{resuelve})
    \item \esp{No pienso que la injusticia \textbf{sea} necesaria.} (subjonctif \emph{sea} de \esp{ser})
    \item \esp{Cada estudiante tiene su \textbf{propia} identidad.} (\emph{identidad} fém. sg. $\rightarrow$ \emph{propia})
    \item \esp{Sin embargo, los estereotipos dividen a la gente.} (ou \esp{dividen a las personas})
\end{enumerate}
\end{corrige}

\begin{corrige}[Exercice 2 — Version]
\begin{enumerate}
    \item « Il faut promouvoir la tolérance à l'école. »
    \item « Nous devons écouter les autres avec attention. »
    \item « Je pense que la cohabitation pacifique est possible. »
    \item « Nous ne pensons pas que la violence soit justifiable. » (subjonctif \emph{sea} $\rightarrow$ « soit »)
    \item « Chaque communauté a ses propres coutumes. »
    \item « Les préjugés empêchent le dialogue sincère. »
\end{enumerate}
\end{corrige}

\begin{corrige}[Exercice 3 — Production écrite (Proposition de correction)]
\esp{En primer lugar, opino que la convivencia sin respeto no es posible. Además, el respeto mutuo es la base de la paz. Sin embargo, no creo que sea fácil lograrlo si hay prejuicios. Por último, hay que educar a los niños para que valoren las propias diferencias de los demás. Debemos construir puentes, no muros. En resumen, sin respeto no hay convivencia real, solo conflicto.}
\emph{(98 mots. Structures cibles présentes : \esp{En primer lugar, Además, Sin embargo, Por último, En resumen, hay que, debemos, no creo que + sea (subj.), los demás, propias diferencias}).}
\end{corrige}

\courssub{Pour aller plus loin : Culture et coopération}

\begin{savaistu}[Cameroun — Guinée équatoriale : voisins hispanophones]
Le Cameroun partage plus de 180 km de frontière avec la \textbf{Guinée équatoriale} (pays hispanophone depuis 1778, seul pays d'Afrique où l'espagnol est langue officielle). Dans les villes frontalières comme \textbf{Kyé-Ossi} ou \textbf{Ebolowa}, le bilinguisme français-espagnol est une réalité quotidienne. De nombreux Camerounais y font du commerce (cacao, bois, produits vivriers) et apprennent l'espagnol « sur le tas ». Depuis 2019, l'espagnol est langue facultative au Bac au Cameroun, et des accords universitaires lient Yaoundé I à l'Université Nationale de Guinée Équatoriale (UNGE) à Malabo et Bata. Apprendre l'espagnol en 1ère A, c'est aussi se préparer à cette coopération régionale concrète.
\end{savaistu}

\coursec{Méthode du débat oral}

Après avoir travaillé le vocabulaire du vivre-ensemble et les connecteurs logiques dans la section précédente (\emph{Traduction : phrases sur l'identité et la tolérance}), nous passons maintenant à la mise en pratique orale. Savoir traduire des phrases isolées est une étape nécessaire ; savoir les enchaîner dans un échange contradictoire, respectueux et structuré en est une autre. Le débat oral est l'activité reine de la compétence \emph{Expresión e interacción orales} : il mobilise la compréhension, la réactivité, la correction grammaticale (subjonctif, indicateurs d'opinion) et la cohérence discursive (connecteurs). Voici une méthode en cinq étapes, testée en classe de Première, pour transformer vos élèves en débateurs efficaces.

\courssub{Étape 1 : Prendre position clairement — L'ouverture du débat}

Tout débat commence par une \cle{thèse} (una \esp{tesis} ou \esp{postura}) formulée sans ambiguïté. En espagnol, on utilise des marqueurs d'opinion personnelle suivis de l'indicatif (certitude) ou du subjonctif (subjectivité, doute, émotion, négation).

\begin{definition}[Marqueurs d'opinion pour ouvrir un débat]
\begin{itemize}
    \item \esp{En mi opinión,} / \esp{A mi juicio,} / \esp{Para mí,} + indicatif (affirmation) ou subjonctif (négation/doute).
    \item \esp{Yo pienso que} / \esp{Yo creo que} + \textbf{indicatif} (certitude subjective).
    \item \esp{Me parece que} / \esp{Considero que} + \textbf{indicatif} (opinion perçue comme fait) ; \esp{No me parece que} / \esp{No considero que} + \textbf{subjonctif}.
    \item \esp{Estoy convencido/a de que} + indicatif.
    \item \esp{Desde mi punto de vista,} + indicatif / subjonctif selon le contexte.
\end{itemize}
\end{definition}

\begin{exemplebox}[Prise de position — Exemples authentiques]
\esp{En mi opinión, la diversidad cultural es una riqueza para nuestra sociedad.} \\
« À mon avis, la diversité culturelle est une richesse pour notre société. »

\esp{Yo pienso que la tolerancia no significa estar de acuerdo con todo.} \\
« Je pense que la tolérance ne signifie pas être d'accord avec tout. »

\esp{Me parece que el respeto mutuo es la base de la convivencia.} \\
« Il me semble que le respect mutuel est la base du vivre-ensemble. »

\esp{Estoy convencida de que debemos escuchar antes de juzgar.} \\
« Je suis convaincue que nous devons écouter avant de juger. »
\end{exemplebox}

\begin{attention}[Piège francophone : \textbf{« Je pense que » $\neq$ \esp{Pienso que} seul}]
En français, « Je pense que » s'omet souvent. En espagnol, le sujet \esp{yo} renforce l'appropriation de l'opinion : \esp{Yo pienso que...} (moi, je pense que...). De plus, \esp{Creer que} et \esp{Pensar que} appellent l'\textbf{indicatif} à l'affirmatif (c'est mon opinion, je la tiens pour vraie), alors que le français utilise souvent le conditionnel pour atténuer (« Je pense qu'il \textbf{faudrait}... »). En espagnol : \esp{Creo que \textbf{hay} que...} (indicatif) ou \esp{Creo que \textbf{deberíamos}...} (conditionnel). À la forme négative : \esp{No creo que \textbf{haya} que...} (subjonctif).
\end{attention}

\courssub{Étape 2 : Argumenter avec des connecteurs — Construire le raisonnement}

Une opinion sans argument est une assertion vide. Les connecteurs logiques (travaillés en grammaire) structurent le discours : ils guident l'auditeur et montrent la rigueur de la pensée.

\begin{propriete}[Les quatre piliers de l'argumentation orale]
\begin{center}
\begin{tabularx}{\linewidth}{|X|X|X|}
\hline
\rowcolor{popblueL}
\textbf{Connecteur} & \textbf{Fonction} & \textbf{Exemple en contexte} \\
\hline
\esp{En primer lugar,} / \esp{Primero,} & Introduire l'argument principal & \esp{En primer lugar, la diversidad nos obliga a salir de nuestra zona de confort.} \\
\hline
\esp{Además,} / \esp{También,} / \esp{Por otra parte,} & Ajouter un argument complémentaire & \esp{Además, el intercambio cultural enriquece nuestra gastronomía y nuestra lengua.} \\
\hline
\esp{Por ejemplo,} / \esp{Como ejemplo,} & Illustrer par un fait concret & \esp{Por ejemplo, en Camerún convivimos más de 250 grupos étnicos.} \\
\hline
\esp{Por tanto,} / \esp{En conclusión,} / \esp{Por eso,} & Conclure l'argumentation & \esp{Por tanto, la diversidad no es un problema, es una oportunidad.} \\
\hline
\end{tabularx}
\end{center}
\end{propriete}

\begin{exoresolu}[Structurer une intervention de 1 minute]
\textbf{Énoncé :} Préparez une intervention courte sur le thème : \esp{« La diversidad cultural enriquece a un país »}. Utilisez au moins trois connecteurs différents.

\tcblower
\textbf{Solution détaillée :}
\begin{enumerate}
    \item \textbf{Ouverture :} \esp{En mi opinión, la diversidad cultural enriquece enormemente a un país.} (Marqueur d'opinion + thèse).
    \item \textbf{Argument 1 :} \esp{En primer lugar, permite conocer otras formas de vida y pensar.} (Premier argument).
    \item \textbf{Argument 2 :} \esp{Además, fomenta la creatividad y la innovación.} (Argument additionnel).
    \item \textbf{Exemple :} \esp{Por ejemplo, en la música, el bendskin camerunés mezcla ritmos tradicionales y modernos.} (Ancrage local).
    \item \textbf{Conclusion :} \esp{Por tanto, debemos valorar y proteger esta diversidad.} (Conséquence logique).
\end{enumerate}
\textit{Durée estimée : 45-60 secondes. Débit naturel, regard vers l'auditoire.}
\end{exoresolu}

\courssub{Étape 3 : Écouter et réagir — L'interaction active}

Le débat n'est pas une succession de monologues. L'écoute active (\esp{escucha activa}) se signale par des reformulations, des demandes de clarification ou des marques d'empathie avant de prendre la parole.

\begin{methode}[Réagir à l'argument de l'autre — 3 formules clés]
\begin{enumerate}
    \item \textbf{Reformuler pour valider la compréhension :} \esp{Si te he entendido bien, dices que...} / \esp{¿Quieres decir que...?}
    \item \textbf{Montrer qu'on a saisi le point de vue :} \esp{Entiendo tu punto de vista.} / \esp{Comprendo lo que dices.} / \esp{Tienes parte de razón.}
    \item \textbf{Demander une précision :} \esp{¿Podrías explicarme qué quieres decir con...?} / \esp{¿A qué te refieres exactamente con...?}
\end{enumerate}
\end{methode}

\begin{exemplebox}[Écoute active — Dialogue simulé]
\esp{— Alumno A: En primer lugar, la inmigración trae problemas de integración.} \\
\esp{— Alumno B: Entiendo tu punto de vista. ¿Te refieres a la barrera del idioma o a las costumbres?} \\
\esp{— Alumno A: A las dos cosas. Por ejemplo, en mi barrio hay tensiones por el ruido.} \\
\esp{— Alumno B: Comprendo. Es un problema real. Sin embargo, creo que la escuela ayuda a resolverlo.}
\end{exemplebox}

\courssub{Étape 4 : Réfuter poliment — L'art de la contre-argumentation}

C'est le cœur du débat démocratique : contredire l'idée sans agresser la personne. La structure « \textbf{Reconnaissance + Pivot + Contre-argument} » est la norme en espagnol standard.

\begin{methode}[Réfuter poliment — La structure en 3 temps]
\begin{enumerate}
    \item \textbf{Reconnaître (Valider l'interlocuteur) :} \esp{Entiendo que...} / \esp{Puede que tengas razón en que...} / \esp{Es verdad que...} / \esp{Reconozco que...}
    \item \textbf{Pivoter (Marquer la transition) :} \esp{Pero...} / \esp{Sin embargo...} / \esp{No obstante...} / \esp{Ahora bien...}
    \item \textbf{Contre-argumenter (Apporter la preuve inverse) :} \esp{Creo que...} / \esp{Mi experiencia me dice que...} / \esp{Los datos muestran que...} + connecteur (\esp{porque}, \esp{ya que}).
\end{enumerate}
\end{methode}

\begin{exemplebox}[Réfutation polie — Exemples traduits]
\esp{Entiendo tu opinión, pero no la comparto.} \\
« Je comprends ton opinion, mais je ne la partage pas. »

\esp{Puede que tengas razón en que hay conflictos, sin embargo, la solución no es cerrar las fronteras.} \\
« Tu as peut-être raison sur le fait qu'il y a des conflits, cependant la solution n'est pas de fermer les frontières. »

\esp{Es cierto que la integración cuesta esfuerzo, no obstante, el enriquecimiento mutuo compensa con creces.} \\
« C'est vrai que l'intégration coûte des efforts, néanmoins l'enrichissement mutuel compense largement. »

\esp{Reconozco que hay retos, pero creo que la educación intercultural es la clave.} \\
« Je reconnais qu'il y a des défis, mais je pense que l'éducation interculturelle est la clé. »
\end{exemplebox}

\begin{attention}[Faux amis et pièges grammaticaux : \textbf{Compartir} vs \textbf{Estar de acuerdo}]
\begin{itemize}
    \item \textbf{Compartir} = \emph{Partager} (un bien, un sentiment, une opinion). \esp{Comparto tu opinión} (Je partage ton avis). \textbf{Pas} : \esp{Comparto con tu opinión} (anglicisme/calque de \emph{share with}).
    \item \textbf{Estar de acuerdo \emph{en que}} + \textbf{Indicatif} (accord sur un fait réel) : \esp{Estamos de acuerdo en que el respeto es fundamental.} (Nous sommes d'accord que le respect est fondamental).
    \item \textbf{Estar de acuerdo \emph{en que}} + \textbf{Subjonctif} (souhait, hypothèse, négation) : \esp{No estoy de acuerdo en que \textbf{se prohíba} la inmigración.} (Je ne suis pas d'accord qu'on interdise l'immigration).
    \item \textbf{Estar de acuerdo \emph{con}} + nom/pronom : \esp{Estoy de acuerdo contigo.} (Je suis d'accord avec toi).
\end{itemize}
\end{attention}

\courssub{Étape 5 : Conclure ou chercher un compromis — La synthèse finale}

Un débat scolaire ne vise pas toujours un vainqueur, mais une \cle{synthèse} ou un \cle{consensus}. La conclusion formelle utilise des connecteurs de clôture et, si possible, identifie un terrain d'entente.

\begin{propriete}[Formules de conclusion et de consensus]
\begin{center}
\begin{tabularx}{\linewidth}{|X|X|}
\hline
\rowcolor{popgreenL}
\textbf{Espagnol} & \textbf{Français / Fonction} \\
\hline
\esp{En conclusión,} / \esp{Para terminar,} & Clôture formelle de l'intervention. \\
\esp{Estamos de acuerdo en que...} & Terrain d'entente identifié (indicatif). \\
\esp{Llegamos a la conclusión de que...} & Synthèse collective. \\
\esp{Quedamos en que...} / \esp{Acordamos que...} & Engagement / décision prise. \\
\esp{Aunque no estemos de acuerdo en todo, coincidimos en lo esencial.} & Désaccord partiel / accord sur l'essentiel. \\
\hline
\end{tabularx}
\end{center}
\end{propriete}

\begin{exemplebox}[Conclusion de débat — Modèles]
\esp{En conclusión, aunque tenemos visiones diferentes sobre la inmigración, estamos de acuerdo en que el respeto a los derechos humanos es innegociable.}
\esp{Para terminar, coincidimos en que la educación es la mejor herramienta para la convivencia. Quedamos en proponer un taller intercultural al director.}
\end{exemplebox}

\courssub{Les règles du débat réglé (Debate reglado)}

Pour que l'échange reste pédagogique et non conflictuel, on instaure des règles explicites, affichées au tableau.

\begin{aretenir}[Règles d'or du débat en classe — \emph{Reglas de oro}]
\begin{enumerate}
    \item \textbf{Temps de parole :} 1 minute maximum par intervention (chronomètre visible). Pas de monopole.
    \item \textbf{Tour de parole :} On lève la main ; le modérateur (professeur ou élève désigné) donne la parole. \textbf{Interdiction d'interrompre} (\esp{No interrumpas}).
    \item \textbf{Attaquer les idées, pas les personnes :} \esp{« Esa idea no me convence porque... »} (Cette idée ne me convainc pas car...) \textbf{et non} \esp{« Tú no entiendes nada »} (Tu ne comprends rien).
    \item \textbf{Langage soigné :} Pas d'injures, pas d'argot excessif, register standard (\esp{usted} ou \esp{tú} selon climat de classe, mais toujours respectueux).
    \item \textbf{Preuve ou exemple :} Toute affirmation générale doit être étayée (\esp{Por ejemplo...}, \esp{Según datos de...}).
\end{enumerate}
\end{aretenir}

\begin{savaistu}[Les « Debates Escolares » en Espagne — Compétition officielle]
Saviez-vous que dans les lycées espagnols (ESO et Bachillerato), le débat est une \textbf{compétition officielle} (\esp{Liga de Debate Escolar}) ? Des équipes de 4 élèves s'affrontent sur une motion imposée (ex. \esp{« ¿Debe prohibirse el uso de móviles en los centros educativos? »}). Les règles sont strictes : \textbf{temps de parole chronométré} (discours d'introduction 3 min, réfutation 2 min, conclusion 1 min), \textbf{pas d'interruption}, notation sur la \emph{argumentación}, la \emph{expresión oral}, la \emph{refutación} et le \emph{trabajo en equipo}. Les finales régionales se tiennent souvent dans les parlements autonomiques (Cortes de Madrid, Parlement de Catalogne...). Au Cameroun, le \emph{Club de Débat} du Lycée Leclerc ou de l'École Normale Supérieure de Yaoundé fonctionne sur ce même modèle !
\end{savaistu}

\courssub{Organigramme du débat : La boucle vertueuse}

Le schéma ci-dessous résume le cycle complet d'un échange argumenté réussi. Il ne s'agit pas d'une ligne droite mais d'une \textbf{boucle} : la conclusion de l'un devient l'opinion de départ du suivant.

\begin{popfigure}
\begin{tikzpicture}[
    node distance=1.8cm and 1.2cm,
    box/.style={draw=popblue, thick, fill=popblueL, rounded corners, minimum width=2.8cm, minimum height=1.2cm, align=center, font=\small\bfseries},
    arrow/.style={-Stealth, thick, popdark},
    looparrow/.style={-Stealth, thick, poporange, dashed}
]
% Nœuds principaux (cercle)
\node[box, align=center] (opinion){Opinión\\ \textit{(Tesis / Postura)}};
\node[box, right=of opinion, align=center] (argumentos){Argumentos\\ \textit{(Conectores lógicos)}};
\node[box, below right=of argumentos, align=center] (escucha){Escucha\\ \textit{(Reformulación / Empatía)}};
\node[box, below left=of escucha, align=center] (replica){Réplica\\ \textit{(Refutación educada)}};
\node[box, left=of replica, align=center] (conclusion){Conclusión\\ \textit{(Síntesis / Consenso)}};

% Flèches principales (cycle)
\draw[arrow] (opinion) -- node[above, sloped, font=\tiny, popmuted] {se apoya en} (argumentos);
\draw[arrow] (argumentos) -- node[right, sloped, font=\tiny, popmuted] {exige} (escucha);
\draw[arrow] (escucha) -- node[below, sloped, font=\tiny, popmuted] {provoca} (replica);
\draw[arrow] (replica) -- node[below, sloped, font=\tiny, popmuted] {conduce a} (conclusion);
\draw[arrow] (conclusion) -- node[left, sloped, font=\tiny, popmuted] {alimenta} (opinion);

% Flèche de retour (amélioration)
\draw[looparrow, bend left=20] (conclusion) to node[left, font=\tiny, poporange] {mejora} (opinion);

% Titre central
\node[align=center, font=\small\itshape, popmuted] at (barycentric cs:opinion=1,argumentos=1,escucha=1,replica=1,conclusion=1) {Ciclo del\\ Debate Constructivo};
\end{tikzpicture}
\legende{Organigramme circulaire du débat argumenté : chaque étape nourrit la suivante. La flèche orange en pointillés symbolise l'amélioration continue de la thèse initiale grâce aux apports contradictoires.}
\end{popfigure}

\courssub{Entraînement guidé : Mini-débats par binômes}

Place à la pratique. L'activité suivante permet d'appliquer la méthode en 15 minutes chrono.

\begin{experience}[Mini-débat : \esp{¿Es la diversidad cultural una riqueza?}]
\textbf{Consigne :} Par binômes (A et B). Tirer au sort le rôle : \textbf{A = POUR} (la diversité est une richesse), \textbf{B = CONTRE / NUANCE} (la diversité pose des problèmes d'unité, de communication, de cohésion nationale).

\textbf{Déroulement (chronométré) :}
\begin{enumerate}
    \item \textbf{Préparation individuelle (3 min) :} Chaque élève note 3 arguments avec connecteurs (\esp{En primer lugar, Además, Por ejemplo}) sur une fiche.
    \item \textbf{Tour 1 — Ouverture (1 min chacun) :} A expose sa thèse + 2 arguments. B écoute, note, ne coupe pas.
    \item \textbf{Tour 2 — Réaction \& Réfutation (1 min chacun) :} B reformule (\esp{Entiendo que...}), pivote (\esp{Sin embargo...}), contre-argumente. A fait de même sur les arguments de B.
    \item \textbf{Tour 3 — Recherche de consensus (1 min ensemble) :} Ils essaient de formuler une phrase commune commençant par \esp{Estamos de acuerdo en que...}
    \item \textbf{Retour classe (5 min) :} 2-3 binômes présentent leur consensus. Professeur note au tableau les meilleures formules de réfutation polie entendues.
\end{enumerate}

\textbf{Grille d'auto-évaluation (à remettre) :}
\begin{center}
\begin{tabular}{|p{5cm}|c|c|c|}
\hline
\rowcolor{poppurpleL}
\textbf{Critère} & \textbf{Non acquis} & \textbf{En cours} & \textbf{Acquis} \\
\hline
J'ai ouvert par \esp{En mi opinión / Yo pienso que} & & & \\
\hline
J'ai utilisé \textbf{3 connecteurs} différents & & & \\
\hline
J'ai reformulé l'argument de l'autre (\esp{Entiendo que...}) & & & \\
\hline
J'ai réfuté poliment (\esp{Puede que..., pero...}) & & & \\
\hline
J'ai conclu par \esp{Estamos de acuerdo en que...} & & & \\
\hline
\end{tabular}
\end{center}
\end{experience}

\courssub{Exercice résolu : De l'écrit à l'oral — Transformer un paragraphe en intervention}

\begin{exoresolu}[Transformer un texte écrit en fiche de débat oral]
\textbf{Énoncé :} Voici un paragraphe rédigé par un élève. Transformez-le en \textbf{fiche de débat} (mots-clés seulement, pas de phrases complètes) utilisable en oral. Identifiez : Thèse, Connecteurs, Exemple, Réfutation anticipée, Conclusion.

\textit{Texte source :}
\esp{« En mi opinión, la tolerancia es fundamental para la convivencia en Yaoundé. En primer lugar, en nuestra ciudad conviven personas de muchas etnias y religiones diferentes. Además, si no hay tolerancia, surgen conflictos como los que vemos a veces en los mercados. Por ejemplo, el otro día vi una discusión entre un vendedor bamileke y un cliente bassa por un malentendido cultural. Sin embargo, algunos dicen que la tolerancia debilita la identidad propia. No estoy de acuerdo porque respetar al otro no significa olvidar quién soy yo. En conclusión, la tolerancia nos hace más fuertes como sociedad camerunesa. »}

\tcblower
\textbf{Solution détaillée — Fiche de débat (format carte mentale) :}

\begin{center}
\begin{tikzpicture}[
    mindmap, grow cyclic, every node/.style=concept, concept color=popblue,
    level 1/.append style={level distance=4cm, sibling angle=72},
    level 2/.append style={level distance=3cm, sibling angle=40},
    font=\small
]
\node[concept, concept color=popblue] {Tesis: \textbf{Tolerancia = Fundamental}}
    child[concept color=popgreen] {node[concept, align=center]{Arg 1: \textbf{Convivencia real} \\ (Yaoundé multiethnique)}}
    child[concept color=poporange] {node[concept, align=center]{Arg 2: \textbf{Evita conflictos} \\ (Ex: Marché Bamileke/Bassa)}}
    child[concept color=poppurple] {node[concept, align=center]{Refutación: \textbf{Identidad} \\ (Respeto $\neq$ Olvido)}}
    child[concept color=poppink] {node[concept, align=center]{Conclusión: \textbf{Fortaleza social} \\ (Camerún unido)}};
\end{tikzpicture}
\end{center}

\textbf{Notes pour l'oral (mots-clés à garder sur carte) :}
\begin{itemize}
    \item \textbf{OPENING:} \esp{En mi opinión, tolerancia = base convivencia Ydé.}
    \item \textbf{CONECTORES:} \esp{En primer lugar / Además / Por ejemplo / Sin embargo / En conclusión.}
    \item \textbf{EJEMPLO CONCRETO:} \esp{Mercado — malentendido cultural Bamileke/Bassa.}
    \item \textbf{REFUTACIÓN ANTICIPADA:} \esp{Objeción: «Debilita identidad». Respuesta: \esp{Respeto $\neq$ Olvido yo.}}
    \item \textbf{CIERRE:} \esp{En conclusión, tolerancia = fuerza Camerún.}
\end{itemize}

\textit{Astuce :} En oral, ne lisez pas le paragraphe. Regardez le jury/la classe. La fiche ne sert que de \textbf{filet de sécurité}.
\end{exoresolu}

\courssub{Synthèse : Check-list du bon débateur (Niveau A2+/B1)}

Avant de passer à la production écrite (Section 6), vérifiez que vous maîtrisez cette check-list. C'est la grille que j'utiliserai pour votre évaluation orale.

\begin{definition}[Check-list compétences débat oral — EA12]
\begin{tabular}{|p{6cm}|c|}
\hline
\rowcolor{popgoldL}
\textbf{Compétence observable} & \textbf{Validé ?} \\
\hline
Prendre la parole avec un marqueur d'opinion correct (\esp{Yo pienso que} + ind. / \esp{Me parece que} + ind. affirmatif) & \\
\hline
Enchaîner 3 arguments avec connecteurs variés (\esp{En primer lugar, Además, Por tanto}) & \\
\hline
Citer un exemple concret, si possible camerounais ou hispanophone & \\
\hline
Écouter sans interrompre ; reformuler l'idée de l'autre (\esp{Entiendo que...}) & \\
\hline
Réfuter poliment : Reconnaissance (\esp{Puede que...}) + Pivot (\esp{Sin embargo}) + Contre-argument & \\
\hline
Utiliser \esp{compartir} (transitif direct) et \esp{estar de acuerdo en que} (+ ind./subj.) correctement & \\
\hline
Conclure par une synthèse ou un consensus (\esp{Estamos de acuerdo en que...}) & \\
\hline
Respecter le temps de parole (1 min) et le registre standard & \\
\hline
\end{tabular}
\end{definition}

\begin{attention}[Dernier piège : \textbf{« Estoy de acuerdo \emph{con} que » — INCORRECT}]
Beaucoup d'élèves disent : \esp{*Estoy de acuerdo con que la diversidad es buena.}
\textbf{Règle :} \esp{Estar de acuerdo} \textbf{+ CON + Nom/Pronom} (\esp{con tu idea, contigo}). \\
\esp{Estar de acuerdo} \textbf{+ EN QUE + Verbe} (\esp{en que es bueno}).
Ne mélangez jamais les deux prépositions. C'est une erreur majeure au DELF/DELE et au BAC.
\end{attention}

Vous êtes maintenant prêts à structurer votre pensée par écrit. La Section 6 (\emph{Production écrite guidée : texte argumenté 100 mots}) reprendra exactement ces connecteurs, cette structure « Thèse - Arguments - Exemple - Réfutation - Conclusion » et ce vocabulaire pour produire un paragraphe cohérent, noté sur 10 points. ¡A escribir!

\coursec{Production écrite guidée : texte argumenté (100 mots)}

Dans la section précédente, tu as appris à \cle{débattre à l'oral} : prendre la parole, donner ton avis, réagir aux arguments des autres avec des connecteurs comme \esp{en primer lugar}, \esp{además} ou \esp{sin embargo}. Nous allons maintenant utiliser exactement les mêmes outils, mais \textbf{à l'écrit}. Le débat oral devient un \cle{texte argumenté} : une petite dissertation espagnole d'environ 100 mots, exercice très fréquent à l'évaluation harmonisée et au baccalauréat. Bonne nouvelle : si tu sais débattre, tu sais déjà rédiger. Il suffit de suivre un plan en cinq étapes, puis de relire avec méthode.

\courssub{Le sujet type et sa consigne}

Voici le sujet sur lequel nous allons travailler pas à pas dans toute cette section :

\begin{exemplebox}[Sujet type d'expression écrite]
\esp{¿Es posible la convivencia entre culturas diferentes? Da tu opinión argumentada (100 palabras).}

« La coexistence entre cultures différentes est-elle possible ? Donne ton opinion argumentée (100 mots). »
\end{exemplebox}

Avant d'écrire la moindre phrase, il faut \textbf{décoder la consigne}. Trois informations essentielles :

\begin{itemize}
\item \esp{¿Es posible...?} : une \cle{question} à laquelle tu dois répondre clairement par oui ou par non (ou « oui, mais... »). Ton texte doit contenir une \textbf{position personnelle}.
\item \esp{Da tu opinión argumentada} : il ne suffit pas de dire \esp{pienso que sí} ; il faut \textbf{justifier} ton avis avec des arguments.
\item \esp{100 palabras} : une \textbf{limite de longueur}. En dessous de 80 mots, le texte est jugé incomplet ; au-dessus de 120, tu risques les fautes inutiles.
\end{itemize}

\begin{attention}[Lire la consigne avant d'écrire]
L'erreur la plus grave au baccalauréat n'est pas la faute de grammaire : c'est le \textbf{hors-sujet}. Si la consigne demande ton opinion (\esp{tu opinión}), ne raconte pas ta vie et ne fais pas un résumé du texte étudié. Réponds à la question posée, avec des arguments. Souligne dans la consigne les mots-clés : \esp{opinión}, \esp{argumentada}, \esp{100 palabras}.
\end{attention}

\courssub{Étape 1 : le remue-méninges (brainstorming)}

Un bon texte de 100 mots se prépare en 5 minutes sur le brouillon. La règle d'or : \textbf{2 arguments + 1 objection}. Pourquoi une objection ? Parce qu'un texte qui reconnaît la difficulté (\esp{sin embargo, los prejuicios existen}) puis la réfute montre un esprit critique : c'est ce qui distingue une bonne copie.

Commence en français si nécessaire, puis traduis en espagnol simple. Exemple de fiche de préparation :

\begin{center}
\begin{tabularx}{\linewidth}{|X|X|}
\hline
\rowcolor{popblueL} \textbf{Français (idée)} & \textbf{Español (phrase simple)} \\
\hline
Argument 1 : la diversité est une richesse (exemple du Cameroun) & \esp{La diversidad es una riqueza: en Camerún conviven más de doscientas lenguas.} \\
\hline
Argument 2 : le dialogue permet de connaître les autres & \esp{El diálogo nos ayuda a conocer a los demás.} \\
\hline
Objection : les préjugés créent des conflits & \esp{Los prejuicios pueden crear conflictos.} \\
\hline
Réfutation : l'éducation au respect est la solution & \esp{La educación en el respeto es la clave para vivir en paz.} \\
\hline
\end{tabularx}
\end{center}

\begin{attention}[Des phrases simples valent mieux que des phrases fausses]
Au niveau Première, n'écris que ce que tu sais écrire correctement. Une phrase simple et juste (\esp{El diálogo es importante}) rapporte plus de points qu'une phrase compliquée pleine d'erreurs. Évite de calquer le français : on dit \esp{estoy de acuerdo} (« je suis d'accord ») et non \esp{*soy acuerdo} ; on dit \esp{me parece que} (« il me semble que ») et non une forme calquée du français « il me semble ».
\end{attention}

\courssub{Étapes 2 à 5 : la rédaction, marche par marche}

Ton texte est un petit escalier de cinq marches. Chaque marche correspond à une ou deux phrases, et chaque marche est signalée par un \cle{connecteur logique} étudié dans la section 3.

\begin{popfigure}
\begin{center}
\begin{tikzpicture}[font=\small]
\draw[fill=popblueL, draw=popblue] (0,0) rectangle (2.4,0.8);
\node[align=center] at (1.2,0.4) {\textbf{1. Introducción}\\ \esp{Pienso que...}};
\draw[fill=popgreenL, draw=popgreen] (1.6,0.8) rectangle (4.4,1.6);
\node[align=center] at (3.0,1.2) {\textbf{2. Argumento 1}\\ \esp{En primer lugar...}};
\draw[fill=poporangeL, draw=poporange] (3.2,1.6) rectangle (6.4,2.4);
\node[align=center] at (4.8,2.0) {\textbf{3. Argumento 2}\\ \esp{Además...}};
\draw[fill=poppinkL, draw=poppink] (4.8,2.4) rectangle (8.4,3.2);
\node[align=center] at (6.6,2.8) {\textbf{4. Objeción refutada}\\ \esp{Sin embargo...}};
\draw[fill=popgoldL, draw=popgold] (6.4,3.2) rectangle (10.4,4.0);
\node[align=center] at (8.4,3.6) {\textbf{5. Conclusión}\\ \esp{Por último...}};
\end{tikzpicture}
\end{center}
\legende{Les cinq marches du texte argumenté : chaque étape est annoncée par un connecteur.}
\end{popfigure}

\textbf{Étape 2 : l'introduction avec ta position.} Une seule phrase suffit. Formules utiles : \esp{Pienso que...} (« je pense que... »), \esp{En mi opinión,...} (« à mon avis... »), \esp{Para mí,...} (« pour moi... »), \esp{Creo que...} (« je crois que... »). Exemple : \esp{En mi opinión, la convivencia entre culturas diferentes es posible y necesaria.} « À mon avis, la coexistence entre cultures différentes est possible et nécessaire. »

\textbf{Étape 3 : les deux arguments.} Le premier est introduit par \esp{en primer lugar} (« en premier lieu »), le second par \esp{además} (« de plus »). Donne si possible un \textbf{exemple concret} (le Cameroun, ton lycée, ta ville) : c'est ce qui rend l'argument vivant. Exemple : \esp{En primer lugar, la diversidad es una riqueza: en Camerún conviven más de doscientas lenguas. Además, el diálogo nos ayuda a conocer a los demás.} « En premier lieu, la diversité est une richesse : au Cameroun cohabitent plus de deux cents langues. De plus, le dialogue nous aide à connaître les autres. »

\textbf{Étape 4 : l'objection réfutée.} \esp{Sin embargo} (« cependant ») introduit la difficulté ; tu la réfutes juste après avec \esp{pero} ou avec ta solution. Exemple : \esp{Sin embargo, los prejuicios pueden crear conflictos, pero la educación puede vencerlos.} « Cependant, les préjugés peuvent créer des conflits, mais l'éducation peut les vaincre. »

\textbf{Étape 5 : la conclusion.} Un connecteur conclusif + une phrase qui reprend ta position de départ.

\begin{propriete}[La conclusion reprend la position de l'introduction]
La conclusion ne doit jamais introduire une idée nouvelle. Elle \textbf{reprend la position annoncée dans l'introduction}, avec un connecteur conclusif : \esp{por último} (« enfin »), \esp{en conclusión} (« en conclusion »), \esp{por eso} (« c'est pourquoi »), \esp{en resumen} (« en résumé »).

\esp{Por último, pienso que la educación en el respeto es la clave para vivir juntos en paz.} « Enfin, je pense que l'éducation au respect est la clé pour vivre ensemble en paix. »
\end{propriete}

\begin{methode}[Compter et structurer : la recette des 100 mots]
\begin{enumerate}
\item \textbf{1 phrase d'introduction} avec ta position (\esp{Pienso que...}) : environ 15 mots.
\item \textbf{2 arguments} (\esp{en primer lugar} / \esp{además}), avec un exemple : environ 40 mots.
\item \textbf{1 objection réfutée} (\esp{sin embargo... pero...}) : environ 20 mots.
\item \textbf{1 conclusion} (\esp{por último} / \esp{en conclusión}) : environ 20 mots.
\item \textbf{Compte les mots} : total environ 95 à 110 mots. Pour compter vite, compte les mots d'une ligne et multiplie par le nombre de lignes.
\end{enumerate}
\end{methode}

\courssub{Étape 6 : la relecture avec grille}

Ne rends jamais ton texte sans le relire. Utilise cette grille de vérification, la même que celle du correcteur :

\begin{center}
\begin{tabularx}{\linewidth}{|X|X|}
\hline
\rowcolor{popgreenL} \textbf{Question de relecture} & \textbf{Quoi vérifier} \\
\hline
Connecteurs présents ? & Au moins 3 : \esp{en primer lugar}, \esp{además}, \esp{sin embargo}, \esp{por último} \\
\hline
Position claire ? & Une phrase avec \esp{pienso que} / \esp{en mi opinión} \\
\hline
Accords ? & \esp{la diversidad es}, \esp{los prejuicios pueden}, \esp{culturas diferentes} \\
\hline
Verbes conjugués ? & Sujet + verbe à la bonne personne : \esp{yo pienso}, \esp{el diálogo ayuda} \\
\hline
Nombre de mots ? & Entre 90 et 110 mots environ \\
\hline
Orthographe ? & Accents (\esp{opinión}, \esp{además}, \esp{educación}), \esp{ñ}, pas de mots français \\
\hline
\end{tabularx}
\end{center}

\begin{attention}[Au Bac, un texte sans connecteurs perd des points de cohérence]
Même sans faute de grammaire, une suite de phrases juxtaposées sans connecteurs est pénalisée : le correcteur évalue la \textbf{cohérence} du texte. Vérifie la présence d'au moins \textbf{trois connecteurs logiques} différents. Attention aussi aux faux amis : \esp{además} signifie « de plus » (et non « à demain ») ; \esp{actualmente} signifie « en ce moment, à présent » (et non « actuellement » au sens de « en réalité ») ; \esp{sin embargo} signifie « cependant » et ne se traduit jamais mot à mot.
\end{attention}

\courssub{Texte modèle commenté}

Voici un modèle complet, rédigé selon les cinq étapes, sur le vivre-ensemble au Cameroun. Compte les mots : environ 100.

\begin{textebox}[Modelo de texto argumentado : la convivencia en Camerún]
En mi opinión, la convivencia entre culturas diferentes es posible y necesaria. En primer lugar, la diversidad es una riqueza: en Camerún conviven más de doscientas lenguas y cada una tiene sus tradiciones, su música y su cocina. Además, el diálogo nos ayuda a conocer a los demás: en mi liceo, los alumnos de diferentes regiones trabajan juntos y aprenden unos de otros. Sin embargo, los prejuicios pueden crear conflictos, pero la escuela y la familia pueden combatirlos. Por último, pienso que la educación en el respeto es la clave para vivir juntos en paz. Un país diverso como Camerún demuestra cada día que la unidad en la diversidad no es un sueño, sino una realidad.
\end{textebox}

\textbf{Glossaire} : \esp{la riqueza} = la richesse ; \esp{conviven} = cohabitent (verbe \esp{convivir}) ; \esp{los demás} = les autres ; \esp{unos de otros} = les uns des autres ; \esp{los prejuicios} = les préjugés ; \esp{combatirlos} = les combattre ; \esp{la clave} = la clé ; \esp{un sueño} = un rêve ; \esp{sino} = mais (après une négation).

\textbf{Commentaire de la structure} : phrase 1 = introduction avec position (\esp{en mi opinión}) ; phrases 2 et 3 = deux arguments avec exemples camerounais (\esp{en primer lugar}, \esp{además}) ; phrase 4 = objection réfutée (\esp{sin embargo... pero...}) ; phrases 5 et 6 = conclusion qui reprend la position (\esp{por último, pienso que...}). Quatre connecteurs, des verbes correctement conjugués, environ 100 mots : la copie idéale.

\begin{savaistu}[La diversité, une réalité camerounaise et hispanophone]
Le Cameroun, avec ses centaines de langues et ses dix régions, est souvent présenté comme « l'Afrique en miniature » : c'est un exemple parfait pour illustrer \esp{la unidad en la diversidad}. Dans le monde hispanophone aussi, la coexistence des cultures est un thème central : en Espagne coexistent le castillan, le catalan, le basque et le galicien, et la Guinée équatoriale, voisine du Cameroun, est le seul pays d'Afrique subsaharienne dont l'espagnol est langue officielle. Citer ces exemples dans ta copie montre ta culture générale.
\end{savaistu}

\begin{exoresolu}[Corriger un texte d'élève]
Énoncé : voici le texte de Marlyse. Trouve les trois problèmes et corrige-les : \esp{La convivencia es posible. La diversidad es una riqueza. El diálogo ayuda. Los prejuicios existen. La educación es la clave.}
\tcblower
Solution détaillée : le texte ne contient \textbf{aucun connecteur logique} : c'est une simple juxtaposition de phrases, pénalisée en cohérence. Il manque aussi une \textbf{position personnelle} explicite et une \textbf{conclusion}. Version corrigée : \esp{Pienso que la convivencia es posible. En primer lugar, la diversidad es una riqueza. Además, el diálogo ayuda a conocer a los demás. Sin embargo, los prejuicios existen, pero se pueden combatir. En conclusión, la educación es la clave.} Quatre connecteurs ajoutés (\esp{en primer lugar}, \esp{además}, \esp{sin embargo}, \esp{en conclusión}), une position (\esp{pienso que}), une vraie conclusion : le texte gagne immédiatement en cohérence.
\end{exoresolu}

\begin{exercice}[À ton tour]
Rédige ton propre texte argumenté de 100 mots environ, au choix :
\begin{enumerate}
\item \esp{¿Es posible la convivencia entre culturas diferentes? Da tu opinión argumentada.}
\item \esp{¿Debe la escuela enseñar la tolerancia? Da tu opinión argumentada.}
\end{enumerate}
Prépare d'abord ta fiche de brainstorming (2 arguments + 1 objection), rédige en suivant les cinq marches, puis relis-toi avec la grille de vérification. Souligne tes connecteurs au crayon : tu dois en avoir au moins trois.
\end{exercice}

\begin{corrige}[Exercice : pistes de correction]
Il n'existe pas une seule bonne réponse, mais toute bonne copie doit contenir : une introduction avec position (\esp{Pienso que...} / \esp{En mi opinión...}), deux arguments introduits par \esp{en primer lugar} et \esp{además}, une objection réfutée avec \esp{sin embargo... pero...}, et une conclusion avec \esp{por último} ou \esp{en conclusión} qui reprend la position initiale. Pour le sujet 2, arguments possibles : \esp{la escuela forma a los futuros ciudadanos} (« l'école forme les futurs citoyens ») ; \esp{los jóvenes aprenden a respetar a los demás} (« les jeunes apprennent à respecter les autres ») ; objection : \esp{algunos piensan que la tolerancia se aprende en casa} (« certains pensent que la tolérance s'apprend à la maison »), réfutation : \esp{pero la escuela y la familia deben trabajar juntas} (« mais l'école et la famille doivent travailler ensemble »). Vérifie enfin le nombre de mots (90 à 110) et les accords.
\end{corrige}

\begin{aretenir}[L'essentiel de la production écrite argumentée]
\begin{itemize}
\item Un texte argumenté de 100 mots = 1 introduction avec position + 2 arguments + 1 objection réfutée + 1 conclusion.
\item Chaque étape est signalée par un connecteur : \esp{en primer lugar}, \esp{además}, \esp{sin embargo}, \esp{por último} / \esp{en conclusión}.
\item La conclusion reprend la position de l'introduction, sans idée nouvelle.
\item Relis toujours avec la grille : connecteurs, position, accords, conjugaison, nombre de mots, accents.
\item Préfère des phrases simples et correctes à des phrases compliquées et fausses.
\end{itemize}
\end{aretenir}

\newpage

\coursec{Activité d'intégration}

\courssub{Objectifs de l'activité}

À l'issue de cette activité d'intégration, l'élève sera capable de :
\begin{itemize}
\item participer à un \cle{débat oral} organisé sur un sujet de société (la diversité culturelle), en respectant des règles de prise de parole ;
\item mobiliser le \cle{lexique du vivre-ensemble} (\esp{tolerancia}, \esp{respeto}, \esp{discriminación}, \esp{convivencia}, \esp{diversidad}...) dans des échanges authentiques ;
\item utiliser correctement les \cle{connecteurs logiques} de l'argumentation : \esp{en primer lugar}, \esp{además}, \esp{sin embargo}, \esp{por último} ;
\item réagir et répliquer poliment avec des formules de concession : \esp{entiendo, pero...}, \esp{estoy de acuerdo hasta cierto punto, sin embargo...} ;
\item traduire des phrases simples sur l'identité et la tolérance, du français vers l'espagnol et de l'espagnol vers le français ;
\item rédiger individuellement un \cle{texte argumenté} d'environ 100 mots présentant et défendant une position personnelle.
\end{itemize}

\courssub{Situation}

Dans le cadre de la \emph{Semana de la Convivencia} organisée par le club d'espagnol du Lycée bilingue de Yaoundé, la classe de Première A a été choisie pour animer le grand débat de clôture. Le thème retenu par les délégués est : \esp{¿La diversidad cultural es una riqueza o un problema?} Le Cameroun, avec ses plus de deux cents groupes ethniques et ses deux langues officielles, est un terrain idéal pour cette réflexion. Le débat se déroulera selon les règles d'un parlement : deux camps, un président de séance, un temps de parole limité, puis un vote final. Voici l'affiche officielle de l'événement, rédigée par le club d'espagnol.

\begin{textebox}[Affiche du club d'espagnol — Lycée bilingue de Yaoundé]
\textbf{GRAN DEBATE PARLAMENTARIO}\par
\textbf{¿La diversidad cultural es una riqueza o un problema?}\par
El club de español invita a todos los alumnos de Primera A a participar en un debate sobre la convivencia. En nuestro país conviven muchas lenguas, tradiciones y religiones. ¿Es esta diversidad una ventaja para la sociedad o una fuente de conflictos?\par
Reglas del debate : cada intervención dura un minuto como máximo. Es obligatorio escuchar con respeto y responder con fórmulas de cortesía : « Entiendo, pero... », « Sin embargo... ». Al final, la clase votará y cada alumno escribirá un texto de cien palabras con su opinión personal.\par
¡Tu voz cuenta! Ven a debatir con tolerancia y respeto.
\end{textebox}

\emph{Glossaire :} \esp{riqueza} = richesse ; \esp{fuente de conflictos} = source de conflits ; \esp{escuchar} = écouter ; \esp{cortesía} = politesse ; \esp{turno de palabra} = tour de parole ; \esp{¡Tu voz cuenta!} = Ta voix compte !

\emph{Traduction de l'affiche :} « Grand débat parlementaire. La diversité culturelle est-elle une richesse ou un problème ? Le club d'espagnol invite tous les élèves de Première A à participer à un débat sur le vivre-ensemble. Dans notre pays cohabitent de nombreuses langues, traditions et religions. Cette diversité est-elle un avantage pour la société ou une source de conflits ? Règles du débat : chaque intervention dure une minute au maximum. Il est obligatoire d'écouter avec respect et de répondre avec des formules de politesse. À la fin, la classe votera et chaque élève écrira un texte de cent mots avec son opinion personnelle. Ta voix compte ! Viens débattre avec tolérance et respect. »

\courssub{Consignes}

\begin{enumerate}
\item \textbf{Préparation du lexique (10 min).} Par groupes de quatre, complétez une carte mentale avec au moins dix mots du vivre-ensemble : \esp{la tolerancia}, \esp{el respeto}, \esp{la discriminación}, \esp{la convivencia}, \esp{la diversidad}, \esp{la identidad}, \esp{el prejuicio}, \esp{la integración}, \esp{el diálogo}, \esp{la paz}. Vérifiez le genre et le pluriel de chaque nom dans le dictionnaire.

\item \textbf{Choix des camps (5 min).} La classe se divise en deux groupes : le camp \esp{a favor} (\esp{la diversidad es una riqueza}) et le camp \esp{en contra} (\esp{la diversidad puede ser un problema}). Un élève est élu \esp{presidente de la sesión} : il distribue la parole et veille au respect du chronomètre.

\item \textbf{Préparation des arguments (15 min).} Individuellement, rédigez sur une fiche \textbf{deux arguments} en faveur de votre camp. Chaque argument doit contenir au moins un connecteur logique. Exemple de trame : \esp{En primer lugar, creo que... porque...} / \esp{Además, en mi opinión...} Pensez à des exemples concrets : les langues du Cameroun, la cuisine, la musique, le football, les fêtes traditionnelles.

\item \textbf{Déroulement du débat (25 min).} Chaque élève intervient pendant \textbf{une minute maximum}. Toute réplique à un camarade doit commencer par une formule de concession polie : \esp{Entiendo tu opinión, pero...}, \esp{Estoy de acuerdo hasta cierto punto, sin embargo...} Le président sanctionne les interruptions et rappelle à l'ordre : \esp{Por favor, respeta el turno de palabra.}

\item \textbf{Vote final (5 min).} À main levée, la classe vote pour la position la plus convaincante. Le président annonce le résultat : \esp{La mayoría piensa que la diversidad cultural es...}

\item \textbf{Production écrite individuelle (20 min).} Rédigez un texte argumenté d'\textbf{environ 100 mots} intitulé \esp{Mi opinión sobre la diversidad cultural}. Votre texte doit contenir : une introduction avec votre position, deux arguments illustrés d'exemples, une concession à l'avis contraire, une conclusion, et \textbf{au moins quatre connecteurs logiques} différents.

\item \textbf{Valorisation.} Les trois meilleurs textes, choisis par la classe, seront affichés au mur de la salle pendant la Semana de la Convivencia.
\end{enumerate}

\courssub{Ressources à mobiliser}

\begin{aretenir}[Les connecteurs logiques de l'argumentation]
\begin{center}
\begin{tabularx}{\linewidth}{|X|X|}\hline
\rowcolor{popblueL} \textbf{Connecteur} & \textbf{Fonction} \\ \hline
\esp{en primer lugar} & introduire le premier argument \\ \hline
\esp{además} / \esp{también} & ajouter un argument \\ \hline
\esp{sin embargo} / \esp{pero} & exprimer une opposition ou une concession \\ \hline
\esp{por ejemplo} & illustrer \\ \hline
\esp{por eso} / \esp{por lo tanto} & exprimer la conséquence \\ \hline
\esp{por último} / \esp{en conclusión} & conclure \\ \hline
\end{tabularx}
\end{center}
\end{aretenir}

\begin{aretenir}[Formules du débat poli]
\begin{itemize}
\item Donner son avis : \esp{creo que...}, \esp{en mi opinión...}, \esp{pienso que...}, \esp{me parece que...}
\item Être d'accord : \esp{estoy de acuerdo contigo}, \esp{es verdad que...}, \esp{tienes razón}
\item Répliquer poliment : \esp{entiendo, pero...}, \esp{estoy de acuerdo hasta cierto punto, sin embargo...}, \esp{respeto tu opinión, pero pienso que...}
\item Demander la parole : \esp{¿puedo decir algo?}, \esp{quiero añadir que...}
\end{itemize}
\end{aretenir}

\begin{attention}[Pièges fréquents des francophones]
\begin{itemize}
\item Ne dites pas \esp{estoy de acuerdo \textbf{avec} tú} : la préposition est \esp{con}, et \esp{con} + \esp{tú} donne la forme contractée \esp{contigo} : \esp{estoy de acuerdo \textbf{contigo}}. De même, \esp{con} + \esp{mí} donne \esp{conmigo}.
\item \esp{Argumentar} ne signifie pas « argumenter » au sens de « discuter longuement » : c'est « donner des arguments ». On dit aussi \esp{dar argumentos}, \esp{defender una opinión}.
\item Attention au faux ami \esp{una riqueza} : c'est « une richesse », pas « une requête ».
\item Ne confondez pas \esp{además} (« de plus », connecteur) et \esp{además de} + nom (« en plus de ») : \esp{Además de la música, me gusta la cocina.} « En plus de la musique, j'aime la cuisine. »
\item Après \esp{creo que} et \esp{pienso que}, on utilise l'indicatif : \esp{Creo que la diversidad \textbf{es} una riqueza.} Mais après la négation, le subjonctif : \esp{No creo que \textbf{sea} un problema.}
\end{itemize}
\end{attention}

\begin{exoresolu}[Traduction — thème et version]
\textbf{Énoncé A (thème).} Traduisez en espagnol : « Je pense que la tolérance est très importante pour vivre ensemble. De plus, le dialogue peut résoudre beaucoup de conflits. »\par
\textbf{Énoncé B (version).} Traduisez en français : \esp{No creo que la discriminación sea una solución ; sin embargo, todavía existe en muchas sociedades.}
\tcblower
\textbf{Corrigé A.} \esp{Pienso que la tolerancia es muy importante para vivir juntos. Además, el diálogo puede resolver muchos conflictos.}\par
Points de méthode : « pour vivre ensemble » se traduit par \esp{para vivir juntos} (\esp{para} + infinitif exprime le but) ; « beaucoup de conflits » = \esp{muchos conflictos} (accord masculin pluriel).\par
\textbf{Corrigé B.} « Je ne crois pas que la discrimination soit une solution ; cependant, elle existe encore dans beaucoup de sociétés. »\par
Points de méthode : après \esp{no creo que}, l'espagnol exige le subjonctif (\esp{sea}), que le français rend naturellement par le subjonctif « soit » ; \esp{sin embargo} = « cependant » ; \esp{todavía} = « encore ».
\end{exoresolu}

\courssub{Proposition de production et corrigé commenté}

\begin{exoresolu}[Texte argumenté modèle — environ 100 mots]
\textbf{Énoncé.} Rédigez un texte argumenté d'environ 100 mots intitulé \esp{Mi opinión sobre la diversidad cultural}, avec au moins quatre connecteurs logiques, une concession et une conclusion.
\tcblower
\textbf{Production modèle :}\par
\esp{\textbf{Mi opinión sobre la diversidad cultural}}\par
\esp{En mi país, Camerún, viven personas de muchas culturas diferentes. En mi opinión, la diversidad cultural es una gran riqueza.}\par
\esp{En primer lugar, la diversidad enriquece nuestra vida diaria : por ejemplo, podemos escuchar músicas distintas y comer platos de todas las regiones. Además, convivir con personas diferentes nos enseña la tolerancia y el respeto.}\par
\esp{Es verdad que a veces las diferencias pueden crear conflictos ; sin embargo, creo que el diálogo siempre es posible.}\par
\esp{En conclusión, la diversidad no es un problema : es una oportunidad para aprender y vivir juntos en paz.}\par
\textbf{Commentaire des choix de langue :}
\begin{itemize}
\item \textbf{Introduction} : la position est annoncée dès la deuxième phrase avec \esp{en mi opinión} + indicatif, conformément à la règle.
\item \textbf{Connecteurs} : quatre connecteurs variés sont utilisés (\esp{en primer lugar}, \esp{además}, \esp{sin embargo}, \esp{en conclusión}), plus \esp{por ejemplo} pour illustrer. Le texte dépasse donc la consigne minimale.
\item \textbf{Concession} : \esp{Es verdad que... ; sin embargo...} reproduit la stratégie de réplique polie du débat oral, transférée à l'écrit.
\item \textbf{Lexique} : les mots clés de la leçon sont bien présents : \esp{diversidad}, \esp{riqueza}, \esp{convivir}, \esp{tolerancia}, \esp{respeto}, \esp{conflictos}, \esp{diálogo}, \esp{vivir juntos en paz}.
\item \textbf{Grammaire} : \esp{vivir juntos} (infinitif après \esp{oportunidad para}), accord correct de \esp{personas diferentes}, ponctuation espagnole respectée. Le texte compte environ 105 mots : la consigne est respectée.
\end{itemize}
\end{exoresolu}

\courssub{Bilan de l'activité}

Cette activité d'intégration a permis de réinvestir, dans une tâche orale et écrite authentique, l'ensemble des acquis de la leçon. À l'oral, les élèves ont appris à défendre une position en un temps limité, à écouter l'autre et à répliquer avec courtoisie : le débat parlementaire transforme la classe en véritable espace de citoyenneté, où l'on s'oppose sans s'insulter, exactement comme dans la vie démocratique. À l'écrit, la rédaction du texte de 100 mots a consolidé l'emploi des connecteurs logiques, qui donnent sa structure à toute argumentation, compétence directement réinvestie dans l'épreuve d'expression écrite du Baccalauréat. Les exercices de traduction ont enfin renforcé deux points sensibles : la préposition \esp{para} + infinitif et le subjonctif après \esp{no creo que}. Le vote final et l'affichage des meilleurs textes valorisent la production de chacun et prolongent la Semana de la Convivencia au-delà de la salle de classe. Pour aller plus loin, la classe pourra organiser un second débat sur un thème voisin (\esp{¿Las redes sociales unen o dividen a los jóvenes?}) ou transformer les textes affichés en articles pour le journal du lycée.

\newpage

\coursec{Méthodes et exercices résolus}

\coursec{EA12 : La identidad, la tolerancia y la convivencia (continuación) : debate y redacción}

\courssub{Texte support : Un débat au lycée de Yaoundé}

\begin{textebox}[Un débat au lycée de Yaoundé]
\esp{El martes pasado, en el liceo de Yaoundé, la profesora de español organizó un debate en la clase de Primera A. El tema era: « La diversidad cultural: ¿riqueza o problema? ».}

\esp{En primer lugar, Marie tomó la palabra: « En mi opinión, la diversidad es una riqueza. En Camerún tenemos más de doscientos grupos étnicos. Cada cultura aporta sus bailes, su música y sus lenguas. Además, el respeto mutuo nos permite vivir en paz. »}

\esp{Sin embargo, Paul no estaba de acuerdo: « Pienso que las diferencias dividen a la gente. A veces hay conflictos entre alumnos de diferentes regiones. Es difícil convivir si no pensamos igual. »}

\esp{La profesora intervino: « Es verdad que surgen conflictos, pero el diálogo es la solución. Por último, quiero recordar que la tolerancia no es aceptar todo sin pensar, es respetar al otro aunque sea diferente. Hay que construir juntos una sociedad unida. »}

\esp{Al final, la clase aplaudió. Los alumnos entendieron que la identidad camerunesa es fuerte porque es diversa.}
\end{textebox}

\begin{center}
\textbf{Traduction / Glossaire}
\end{center}
\begin{itemize}
\item \esp{El martes pasado} : mardi dernier.
\item \esp{organizó un debate} : a organisé un débat.
\item \esp{tomó la palabra} : a pris la parole (pretérito indefinido de \esp{tomar}).
\item \esp{aporta} : apporte (présent de \esp{aportar}).
\item \esp{mutuo} : mutuel.
\item \esp{surgen conflictos} : surgissent des conflits (présent de \esp{surgir}).
\item \esp{la solución} : la solution.
\item \esp{construir juntos} : construire ensemble.
\item \esp{aplaudió} : a applaudi.
\end{itemize}

\courssub{Lexique thématique : Débat et vivre-ensemble}

\begin{definition}[Vocabulaire essentiel du thème]
\begin{center}
\begin{tabularx}{\linewidth}{|X|X|X|}
\hline
\rowcolor{popblueL} \textbf{Français} & \textbf{Español} & \textbf{Remarque} \\
\hline
la tolérance & \esp{la tolerancia} & nom féminin, accent sur \esp{á} \\
le respect & \esp{el respeto} & \\
l'identité & \esp{la identidad} & \textbf{sans accent} (termine en -d) \\
la discrimination & \esp{la discriminación} & accent sur \esp{ó} \\
la diversité & \esp{la diversidad} & \\
le vivre-ensemble & \esp{la convivencia} & de \esp{convivir} \\
le dialogue & \esp{el diálogo} & accent sur \esp{á} \\
le conflit & \esp{el conflicto} & \\
la solution & \esp{la solución} & accent sur \esp{ó} \\
la richesse & \esp{la riqueza} & \\
le groupe ethnique & \esp{el grupo étnico} & \\
les droits & \esp{los derechos} & \\
l'opinion & \esp{la opinión} & accent sur \esp{ó} \\
\hline
\end{tabularx}
\end{center}

\textbf{Verbes utiles :} \esp{respetar} (respecter), \esp{aceptar} (accepter), \esp{compartir} (partager), \esp{convivir} (cohabiter/vivre ensemble), \esp{dialogar} (dialoguer), \esp{discriminar} (discriminer), \esp{enriquecer} (enrichir), \esp{dividir} (diviser), \esp{unir} (unir), \esp{construir} (construire), \esp{aprender a} (apprendre à), \esp{estar de acuerdo} (être d'accord), \esp{estar en contra} (être contre).

\textbf{Expressions de débat :} \esp{En mi opinión...} (À mon avis...), \esp{Yo creo que...} (Je crois que...), \esp{Me parece que...} (Il me semble que...), \esp{Desde mi punto de vista...} (De mon point de vue...), \esp{Estoy de acuerdo} (Je suis d'accord), \esp{No estoy de acuerdo} (Je ne suis pas d'accord), \esp{Es verdad que...} (Il est vrai que...), \esp{Puede ser, pero...} (C'est possible, mais...).
\end{definition}

\courssub{Grammaire : Les connecteurs logiques de l'argumentation}

\begin{propriete}[Tableau des connecteurs logiques (niveau A2+/B1)]
\begin{center}
\begin{tabularx}{\linewidth}{|X|X|X|}
\hline
\rowcolor{popblueL} \textbf{Fonction} & \textbf{Connecteurs espagnols} & \textbf{Français} \\
\hline
Ordonner (début) & \esp{en primer lugar} & d'abord, en premier lieu \\
Ordonner (suite) & \esp{en segundo lugar}, \esp{luego}, \esp{después} & ensuite, puis \\
Ordonner (fin) & \esp{por último}, \esp{en último lugar} & enfin, en dernier lieu \\
Ajouter & \esp{además}, \esp{también}, \esp{es más}, \esp{igualmente} & de plus, aussi, qui plus est, de même \\
Opposer / Nuancer & \esp{sin embargo}, \esp{no obstante}, \esp{pero}, \esp{aunque} + indicatif/subjonctif & cependant, néanmoins, mais, bien que \\
Concession & \esp{es verdad que... sin embargo...}, \esp{puede ser que... pero...} & il est vrai que... cependant..., il se peut que... mais... \\
Conclure & \esp{en conclusión}, \esp{por último}, \esp{por eso}, \esp{por lo tanto}, \esp{en definitiva} & en conclusion, enfin, c'est pourquoi, donc, en définitive \\
\hline
\end{tabularx}
\end{center}

\textbf{Règles de placement :} Le connecteur se place \textbf{en tête de phrase}, suivi \textbf{obligatoirement d'une virgule}. Ex. : \esp{En primer lugar, la tolerancia es clave.} Ne jamais écrire : \esp{La tolerancia, en primer lugar, es clave.} (sauf cas très rares de mise en relief à l'oral).
\textbf{Orthographe :} \esp{además} (accent sur \esp{á}), \esp{sin embargo} (deux mots), \esp{por último} (pas \esp{por último lugar} seul pour « enfin » dans une énumération).
\end{propriete}

\begin{methode}[Choisir le bon connecteur : démarche en 4 étapes]
\begin{enumerate}
\item \textbf{Analyser la relation logique} entre la phrase à écrire et la précédente : est-ce un ajout (\cle{además}), une opposition (\cle{sin embargo}), un ordre (\cle{en primer lugar}) ou une conclusion (\cle{por último}) ?
\item \textbf{Éviter la répétition} : ne pas utiliser deux fois \esp{además} dans le même paragraphe ; varier avec \esp{también}, \esp{es más}.
\item \textbf{Respecter la ponctuation} : connecteur en début de phrase + virgule. \esp{Sin embargo, creo que...}
\item \textbf{Vérifier l'accord} si le connecteur inclut un participe ou un adjectif (ex. \esp{visto que...} vu que...), mais les connecteurs de cette liste sont invariables.
\end{enumerate}
\end{methode}

\begin{exoresolu}[Compléter un texte avec les connecteurs]
\textbf{Énoncé.} Complète le texte suivant avec les connecteurs : \esp{además, sin embargo, en primer lugar, por último, por eso}.

\esp{La tolerancia es esencial en la escuela. \_\_\_\_\_\_, todos los alumnos tienen derecho al respeto. \_\_\_\_\_\_, la convivencia mejora el aprendizaje. \_\_\_\_\_\_, a veces hay conflictos entre alumnos de diferentes regiones; \_\_\_\_\_\_, el diálogo puede resolverlos. \_\_\_\_\_\_ debemos practicar la tolerancia cada día.}
\tcblower
\textbf{Raisonnement.}
\begin{enumerate}
\item Phrase 2 : premier argument principal $\rightarrow$ \esp{En primer lugar}.
\item Phrase 3 : ajout d'un second argument $\rightarrow$ \esp{Además}.
\item Phrase 4 : structure « problème ; solution » $\rightarrow$ opposition sur le second membre $\rightarrow$ \esp{sin embargo}. Le premier membre de la phrase (\esp{a veces hay conflictos...}) ne prend pas de connecteur de la liste (c'est un énoncé factuel), mais \esp{Por último} peut introduire cette dernière idée argumentative avant la conclusion. \textbf{Choix optimal :} \esp{Por último} pour introduire le dernier point du débat (les conflits), \esp{sin embargo} pour la solution.
\item Phrase 5 : conséquence logique $\rightarrow$ \esp{Por eso}.
\end{enumerate}

\textbf{Texte complété :}

\esp{La tolerancia es esencial en la escuela. \textbf{En primer lugar}, todos los alumnos tienen derecho al respeto. \textbf{Además}, la convivencia mejora el aprendizaje. \textbf{Por último}, a veces hay conflictos entre alumnos de diferentes regiones; \textbf{sin embargo}, el diálogo puede resolverlos. \textbf{Por eso} debemos practicar la tolerancia cada día.}

\textbf{Traduction :} « La tolérance est essentielle à l'école. D'abord, tous les élèves ont droit au respect. De plus, le vivre-ensemble améliore l'apprentissage. Enfin, il y a parfois des conflits entre élèves de différentes régions ; cependant, le dialogue peut les résoudre. C'est pourquoi nous devons pratiquer la tolérance chaque jour. »

\textbf{Justifications grammaticales :} \esp{resolverlos} = infinitif \esp{resolver} + pronom \esp{los} (masculin pluriel, antécédent \esp{conflictos}) accolé ; \esp{tienen derecho al respeto} = structure \esp{tener derecho a + nom} ; \esp{por eso} introduit une conséquence, jamais une opposition.
\end{exoresolu}

\begin{popfigure}
\begin{tikzpicture}[
  mindmap,
  grow cyclic,
  every node/.style=concept,
  concept color=popblue,
  level 1/.append style={level distance=3.5cm, sibling angle=90},
  level 2/.append style={level distance=2.5cm, sibling angle=45},
  text=white,
  font=\small
]
\node {Conectores}
  child [concept color=popgreen] {node {Ordenar}
    child {node {En primer lugar}}
    child {node {En segundo lugar}}
    child {node {Por último}}
  }
  child [concept color=poporange] {node {Añadir}
    child {node {Además}}
    child {node {También}}
    child {node {Es más}}
  }
  child [concept color=poppink] {node {Oponer}
    child {node {Sin embargo}}
    child {node {No obstante}}
    child {node {Pero}}
    child {node {Aunque}}
  }
  child [concept color=popgold] {node {Concluir}
    child {node {En conclusión}}
    child {node {Por eso}}
    child {node {Por lo tanto}}
  };
\end{tikzpicture}
\legende{Carte mentale des connecteurs logiques pour l'argumentation}
\end{popfigure}

\courssub{Traduction : Phrases sur l'identité et la tolérance}

\begin{methode}[Traduire du français vers l'espagnol (Thème) : 5 étapes]
\begin{enumerate}
\item \textbf{Identifier le temps} : vérité générale, habitude, description $\rightarrow$ \textbf{présent de l'indicatif}. Action passée datée $\rightarrow$ \textbf{pretérito indefinido}.
\item \textbf{Choisir le verbe exact} : « respecter » $\rightarrow$ \esp{respetar} ; « accepter » $\rightarrow$ \esp{aceptar} ; « partager » $\rightarrow$ \esp{compartir} ; « vivre ensemble » $\rightarrow$ \esp{convivir} ; « être » $\rightarrow$ \esp{ser} (identité, caractéristique permanente) ou \esp{estar} (état, lieu, résultat).
\item \textbf{Éviter les calques français} :
\begin{itemize}
\item « être d'accord » $\rightarrow$ \esp{estar de acuerdo} (jamais \esp{ser de acuerdo}).
\item « il faut / il est nécessaire » $\rightarrow$ \esp{hay que + infinitif} (impersonnel) ou \esp{deber + infinitif} (obligation morale).
\item « apprendre à » $\rightarrow$ \esp{aprender a + infinitif}.
\item « commencer à » $\rightarrow$ \esp{empezar a + infinitif} (changement de radical \esp{e $\rightarrow$ ie}).
\end{itemize}
\item \textbf{Soigner les prépositions} : \esp{para} (but, destination), \esp{por} (cause, moyen), \esp{con} (compagnie), \esp{de} (possession, origine).
\item \textbf{Vérifier les accents et la ponctuation} : \esp{discriminación, opinión, además, también, país, diálogo} ; ¿ ? pour les questions, ¡ ! pour les exclamations.
\end{enumerate}
\end{methode}

\begin{exoresolu}[Thème : Phrases sur la tolérance et l'identité]
\textbf{Énoncé.} Traduis en espagnol :
\begin{enumerate}
\item Nous devons respecter l'identité de chaque personne.
\item La tolérance est très importante pour le vivre-ensemble.
\item Je ne suis pas d'accord avec la discrimination.
\item Au lycée, les élèves apprennent à dialoguer.
\item Il faut accepter les différences pour vivre en paix.
\end{enumerate}
\tcblower
\textbf{Solutions détaillées :}

1. \esp{Debemos respetar la identidad de cada persona.}
\begin{itemize}
\item \esp{Debemos} = \esp{deber} (obligation morale) à la 1ère personne du pluriel, présent.
\item \esp{identidad} : \textbf{pas d'accent} (mot terminé par consonne \esp{d}, tonique sur dernière syllabe = aguda, mais règle : agudas en consonne $\neq$ n,s $\rightarrow$ pas d'accent).
\item \esp{cada persona} = chaque personne (\esp{cada} invariable).
\end{itemize}

2. \esp{La tolerancia es muy importante para la convivencia.}
\begin{itemize}
\item \esp{es} = \esp{ser} (caractère permanent, définition), 3ème pers. sg. \textbf{Pas \esp{está}}.
\item \esp{para} = but / finalité (pour le vivre-ensemble).
\item \esp{convivencia} = vivre-ensemble (nom féminin, pas d'accent).
\end{itemize}

3. \esp{No estoy de acuerdo con la discriminación.}
\begin{itemize}
\item \esp{estoy} = \esp{estar} (expression figée \esp{estar de acuerdo}), 1ère pers. sg.
\item \esp{discriminación} : \textbf{accent sur \esp{ó}} (aguda terminée par \esp{n}).
\item \esp{con} = avec (préposition régie par \esp{acuerdo con}).
\end{itemize}

4. \esp{En el liceo, los alumnos aprenden a dialogar.}
\begin{itemize}
\item \esp{liceo} = lycée (terme courant au Cameroun et en Am. Latine ; \esp{instituto} en Espagne).
\item \esp{aprenden} = \esp{aprender} (présent, 3ème pl.) + \textbf{\esp{a}} + infinitif.
\item \esp{dialogar} = dialoguer.
\end{itemize}

5. \esp{Hay que aceptar las diferencias para vivir en paz.}
\begin{itemize}
\item \esp{Hay que} = il faut (impersonnel, invariables) + infinitif.
\item \esp{aceptar} = accepter.
\item \esp{las diferencias} = les différences (accord fém. pl.).
\item \esp{para vivir} = pour vivre (but).
\item \esp{en paz} = en paix.
\end{itemize}
\end{exoresolu}

\begin{methode}[Comprendre et commenter un texte de débat (Version / Compréhension) : 5 étapes]
\begin{enumerate}
\item \textbf{Lecture globale} : lire deux fois. 1ère fois : sujet général, ton, protagonistes. 2ème fois : souligner mots-clés (\esp{tolerancia, respeto, identidad, diversidad, conflicto, diálogo}).
\item \textbf{Questions de compréhension} : répondre \textbf{en espagnol}, phrases complètes, \textbf{reformuler} (ne pas recopier bêtement), citer un court extrait entre guillemets comme preuve.
\item \textbf{Analyse du point de vue} : qui dit quoi ? Quels connecteurs utilise-t-il ? (\esp{sin embargo, además}).
\item \textbf{Version (traduction vers le français)} : phrase par phrase. Repérer Sujet + Verbe + Compléments. Respecter les temps : \esp{pretérito indefinido} $\rightarrow$ passé composé ou imparfait selon aspect ; \esp{presente} $\rightarrow$ présent ; \esp{hay que} $\rightarrow$ « il faut ».
\item \textbf{Vérification faux amis} : \esp{discutir} = se disputer / débattre vivement (pas « discuter » = bavarder $\rightarrow$ \esp{charlar, conversar}) ; \esp{actualmente} = actuellement (pas « actuellement » au sens de « en ce moment » seulement, c'est bien le sens) ; \esp{actual} = actuel / présent (pas « actuel » = réel) ; \esp{realizar} = réaliser / accomplir (pas « réaliser » = comprendre $\rightarrow$ \esp{darse cuenta}).
\end{enumerate}
\end{methode}

\begin{exoresolu}[Compréhension + Version : Un débat au lycée]
\textbf{Texte.}
\esp{Ayer, en el liceo de Yaoundé, los alumnos de Primera participaron en un debate sobre la tolerancia. Unos alumnos dijeron que la diversidad cultural es una riqueza; otros pensaron que las diferencias pueden dividir a la sociedad. Sin embargo, todos estuvieron de acuerdo en una cosa: hay que respetar a los demás para vivir juntos en paz.}

\textbf{Énoncé.}
\begin{enumerate}
\item ¿Sobre qué tema fue el debate?
\item ¿Cuáles eran las dos opiniones principales?
\item Traduis la dernière phrase en français.
\end{enumerate}
\tcblower
\textbf{1) Réponse :} \esp{El debate fue sobre la tolerancia y la convivencia (o: la diversidad cultural).}
\textbf{Justification :} Le texte dit explicitement : \esp{« un debate sobre la tolerancia »} et le thème annoncé est \esp{« La diversidad cultural: ¿riqueza o problema? »}.

\textbf{2) Réponse :} \esp{Unos alumnos dijeron que la diversidad cultural es una riqueza; otros pensaron que las diferencias pueden dividir a la sociedad.}
\textbf{Justification :} On reformule avec les verbes du texte (\esp{dijeron, pensaron} = pretérito indefinido, actions passées ponctuelles). Les deux opinions sont contrastées par le point-virgule dans le texte original.

\textbf{3) Version :}
\esp{Sin embargo, todos estuvieron de acuerdo en una cosa: hay que respetar a los demás para vivir juntos en paz.}
$\rightarrow$ « Cependant, tous furent d'accord / ont été d'accord sur une chose : il faut respecter les autres pour vivre ensemble en paix. »

\textbf{Justifications grammaticales détaillées :}
\begin{itemize}
\item \esp{Sin embargo} = Cependant (connecteur d'opposition, en tête, virgule).
\item \esp{todos estuvieron de acuerdo} = \esp{estar} au \textbf{pretérito indefinido} (3ème pl. \esp{estuvieron}) : accord conclu à un moment précis (hier). \esp{de acuerdo} = locution adjectivale invariable.
\item \esp{en una cosa} = sur une chose / un point.
\item \esp{hay que respetar} = \textbf{il faut respecter} (structure impersonnelle \esp{hay que} + infinitif, pas de sujet).
\item \esp{a los demás} = les autres (\esp{a} personnel devant complément d'objet direct animé + \esp{los demás} pronom indéfini masculin pluriel).
\item \esp{para vivir} = pour vivre (\esp{para} + infinitif = but).
\item \esp{juntos} = ensemble (adverbe, s'accorde en genre/nombre avec le sujet sous-entendu de \esp{vivir} $\rightarrow$ \esp{nosotros} $\rightarrow$ masculin pluriel).
\item \esp{en paz} = en paix.
\end{itemize}
\end{exoresolu}

\courssub{Méthode : Participer à un débat oral sur le vivre-ensemble}

\begin{methode}[Débattre en espagnol : la démarche en 5 étapes]
\begin{enumerate}
\item \textbf{Comprendre le sujet} : repère le thème (\esp{la tolerancia}, \esp{la discriminación}, \esp{la convivencia}, \esp{la identidad}) et formule ta position : \esp{estoy a favor de...} (je suis pour), \esp{estoy en contra de...} (je suis contre), \esp{estoy de acuerdo con...} (je suis d'accord avec).
\item \textbf{Prendre la parole poliment} : utilise des formules d'ouverture : \esp{En mi opinión...}, \esp{Yo creo que...}, \esp{Me parece que...}, \esp{Desde mi punto de vista...}, \esp{Permítanme decir que...} (permettez-moi de dire que...).
\item \textbf{Structurer ton intervention} avec les connecteurs logiques : \esp{En primer lugar} (d'abord), \esp{Además} (de plus), \esp{Sin embargo} (cependant), \esp{Por último} (enfin).
\item \textbf{Réagir aux autres} (interaction) :
\begin{itemize}
\item Accord : \esp{Estoy de acuerdo contigo}, \esp{Es verdad que...}, \esp{Exacto}, \esp{Comparto tu opinión}.
\item Désaccord : \esp{No estoy de acuerdo}, \esp{Sin embargo, pienso que...}, \esp{Puede ser, pero...}, \esp{No lo veo así}.
\item Demander clarification : \esp{¿Qué quieres decir con...?}, \esp{¿Podrías explicarlo?}
\end{itemize}
\item \textbf{Conclure} : résume ta position : \esp{En conclusión...}, \esp{Para terminar, creo que...}, \esp{En definitiva...}.
\end{enumerate}
\textbf{Réflexes grammaticaux clés :}
\begin{itemize}
\item Parle au \textbf{présent de l'indicatif} (vérités générales, opinions).
\item Utilise \textbf{\esp{hay que + infinitif}} pour proposer des solutions générales : \esp{Hay que respetar a los demás}, \esp{Hay que dialogar más}.
\item Utilise \textbf{\esp{deber + infinitif}} pour l'obligation morale : \esp{Debemos aceptar las diferencias}.
\item Prépare \textbf{2 ou 3 arguments} écrits sur une fiche avant de parler.
\end{itemize}
\end{methode}

\begin{exoresolu}[Débat : ¿Debe haber cursos de tolerancia en el liceo?]
\textbf{Énoncé.} Ton professeur organise un débat : \esp{¿Debe haber cursos de tolerancia y convivencia en el liceo?} Rédige ton intervention orale (6 à 8 phrases) avec au moins \textbf{trois connecteurs logiques} différents et une réaction à un camarade qui a dit : \esp{« Estos cursos son una pérdida de tiempo »}.
\tcblower
\textbf{Raisonnement.} Je choisis ma position : \textbf{POUR} ces cours. Arguments : 1. Prévention des conflits. 2. Formation citoyenne. 3. Concession : il faut des méthodes pratiques. Réaction à l'argument adverse avec \esp{No estoy de acuerdo} + \esp{Sin embargo}.

\textbf{Intervention rédigée :}

\esp{No estoy de acuerdo con mi compañero: estos cursos no son una pérdida de tiempo. \textbf{En primer lugar}, la tolerancia es fundamental para la convivencia entre alumnos de diferentes culturas y orígenes. \textbf{Además}, el liceo no solo enseña materias académicas: también forma ciudadanos responsables y críticos. \textbf{Sin embargo}, es verdad que la metodología debe ser participativa y no solo teórica. \textbf{Por último}, creo que hablar abiertamente de la discriminación ayuda a prevenirla. \textbf{En conclusión}, estoy a favor de estos cursos porque \textbf{hay que aprender a respetar a los demás} para construir una sociedad justa.}

\textbf{Traduction :} « Je ne suis pas d'accord avec mon camarade : ces cours ne sont pas une perte de temps. D'abord, la tolérance est fondamentale pour le vivre-ensemble entre élèves de différentes cultures et origines. De plus, le lycée n'enseigne pas seulement des matières académiques : il forme aussi des citoyens responsables et critiques. Cependant, il est vrai que la méthodologie doit être participative et pas seulement théorique. Enfin, je pense que parler ouvertement de la discrimination aide à la prévenir. En conclusion, je suis favorable à ces cours parce qu'il faut apprendre à respecter les autres pour construire une société juste. »

\textbf{Justifications grammaticales :}
\begin{itemize}
\item \esp{no estoy de acuerdo con} : formule figée de désaccord (\esp{estar} + \esp{de acuerdo} + \esp{con}).
\item \esp{enseña} / \esp{forma} : présent de l'indicatif (vérité générale).
\item \esp{hay que aprender / respetar} : obligation impersonnelle (pas de sujet).
\item Connecteurs en tête de phrase + virgule : \esp{En primer lugar,}, \esp{Además,}, \esp{Sin embargo,}, \esp{Por último,}, \esp{En conclusión,}.
\item \esp{abiertamente} : adverbe en \esp{-mente} (fém. singulier de \esp{abierta} + \esp{mente}).
\end{itemize}
\end{exoresolu}

\courssub{Production écrite guidée : Texte argumenté (100 mots)}

\begin{methode}[Rédiger un court texte argumentatif (90--110 mots) : Structure en 4 paragraphes]
\begin{enumerate}
\item \textbf{Introduction} (1--2 phrases) : Accroche + Sujet + Thèse claire.
Modèle : \esp{Hoy en día, la convivencia entre personas de diferentes culturas es un tema crucial en Camerún. En mi opinión, la diversidad es una riqueza que hay que proteger.}
\item \textbf{Argument 1} (2--3 phrases) : \textbf{\esp{En primer lugar,}} + idée force + \textbf{exemple concret} (ton quartier, ton lycée, le Cameroun, la Guinée équatoriale).
Modèle : \esp{En primer lugar, cada grupo étnico aporta sus tradiciones, lenguas y comidas, lo que enriquece nuestra cultura nacional.}
\item \textbf{Argument 2 + Concession} (2--3 phrases) : \textbf{\esp{Además,}} + 2ème idée + \textbf{\esp{Es verdad que... sin embargo...}} (nuance).
Modèle : \esp{Además, el diálogo intercultural fomenta la paz. Es verdad que surgen malentendidos; sin embargo, la educación y el respeto mutuo los resuelven.}
\item \textbf{Conclusion} (1--2 phrases) : \textbf{\esp{Por último / En conclusión,}} + rappel thèse + \textbf{ouverture / appel à l'action} (\esp{hay que + infinitif}).
Modèle : \esp{En conclusión, la identidad camerunesa es fuerte porque es diversa. Hay que valorar esta riqueza y construir juntos una sociedad unida.}
\end{enumerate}
\textbf{Check-list avant de rendre :}
\begin{itemize}
\item \cle{Nombre de mots} : 90 à 110 (compte les mots espagnols).
\item \cle{Connecteurs} : au moins 4 différents (\esp{en primer lugar, además, sin embargo, por último / en conclusión}).
\item \cle{Temps} : Présent de l'indicatif (sauf si anecdote passée $\rightarrow$ pretérito indefinido).
\item \cle{Vocabulaire thème} : \esp{tolerancia, respeto, identidad, discriminación, diversidad, derechos, convivencia, diálogo, riqueza, conflicto, sociedad}.
\item \cle{Orthographe} : Accents (\esp{además, opinión, discriminación, país, diálogo}), \esp{ñ}, \esp{¿ ¡}.
\end{itemize}
\end{methode}

\begin{exoresolu}[Rédaction : La diversité culturelle au Cameroun]
\textbf{Énoncé.} \esp{Redacta un texto de unas 100 palabras sobre el tema: « La diversidad cultural es una riqueza para el Camerún ». Utiliza al menos cuatro conectores lógicos.}
\tcblower
\textbf{Raisonnement (Plan de rédaction).}
\begin{enumerate}
\item \textbf{Intro :} Cameroun = pays multiethnique (200+ groupes). Thèse : richesse.
\item \textbf{Arg 1 (\esp{En primer lugar}) :} Apports culturels concrets (musique, langues, cuisine, danses \esp{bendskin, makossa, bikutsi}).
\item \textbf{Arg 2 (\esp{Además}) + Concession (\esp{Es verdad que... sin embargo}) :} Vivre-ensemble $\rightarrow$ tolérance. Conflits possibles mais résolus par dialogue/éducation.
\item \textbf{Conclusion (\esp{En conclusión}) :} Identité forte par la diversité. Appel : \esp{hay que valorar / construir}.
\end{enumerate}

\textbf{Texte rédigé (103 mots) :}

\esp{El Camerún es un país con más de doscientos grupos étnicos. \textbf{En mi opinión}, esta diversidad cultural es una gran riqueza para la nación. \textbf{En primer lugar}, cada cultura aporta su música —como el makossa o el bikutsi—, sus lenguas y su gastronomía, y esto hace la vida más interesante y colorida. \textbf{Además}, la convivencia diaria entre diferentes pueblos enseña la tolerancia y el respeto mutuo a los jóvenes. \textbf{Es verdad que} a veces las diferencias pueden crear tensiones o conflictos; \textbf{sin embargo}, el diálogo sincero y la educación cívica permiten superarlos. \textbf{En conclusión}, la identidad camerunesa es fuerte precisamente porque es diversa. \textbf{Hay que valorar} esta riqueza y \textbf{construir juntos} una sociedad más unida y justa.}

\textbf{Traduction :} « Le Cameroun est un pays de plus de deux cents groupes ethniques. À mon avis, cette diversité culturelle est une grande richesse pour la nation. D'abord, chaque culture apporte sa musique — comme le makossa ou le bikutsi —, ses langues et sa gastronomie, et cela rend la vie plus intéressante et colorée. De plus, le vivre-ensemble quotidien entre différents peuples enseigne la tolérance et le respect mutuel aux jeunes. Il est vrai que parfois les différences peuvent créer des tensions ou des conflits ; cependant, le dialogue sincère et l'éducation civique permettent de les surmonter. En conclusion, l'identité camerounaise est forte précisément parce qu'elle est diverse. Il faut valoriser cette richesse et construire ensemble une société plus unie et juste. »

\textbf{Justifications linguistiques :}
\begin{itemize}
\item \textbf{Connecteurs (5) :} \esp{En mi opinión} (marqueur d'opinion), \esp{En primer lugar}, \esp{Además}, \esp{Es verdad que... sin embargo}, \esp{En conclusión}.
\item \textbf{Grammaire :} \esp{aporta, hace, enseña, permiten, superarlos} (présent indicatif) ; \esp{superarlos} = \esp{superar} + \esp{los} (conflits, masc. pl.) ; \esp{hay que valorar / construir} (obligation impersonnelle) ; \esp{precisamente porque} (insistance sur la cause).
\item \textbf{Vocabulaire contexte Cameroun :} \esp{makossa, bikutsi, gastronomía, educación cívica}.
\item \textbf{Orthographe :} \esp{país} (accent), \esp{opinión} (accent), \esp{nación} (accent), \esp{además} (accent), \esp{convivencia} (pas d'accent), \esp{diálogo} (accent).
\end{itemize}
\end{exoresolu}

\begin{attention}[Les 5 erreurs qui coûtent des points au Bac]
\begin{enumerate}
\item \textbf{Oublier les accents obligatoires} : \esp{discriminación, opinión, además, también, país, diálogo, nación, corazón, ratón, canción} prennent l'accent tonique (mots aigus en \esp{-n, -s} ou voyelle). \esp{Identidad, convivencia, tolerancia, riqueza, verdad, ciudad, universidad} \textbf{ne prennent PAS d'accent} (mots llanos terminés en \esp{-d, -n, -s, -z} ou voyelle). \textbf{Astuce :} compte les syllabes à l'envers.
\item \textbf{Le faux ami « estar de acuerdo »} : « Être d'accord » = \esp{estar de acuerdo} (jamais \esp{ser de acuerdo}). « Discuter / débattre » = \esp{debatir} ou \esp{discutir} (mais \esp{discutir} implique souvent une dispute, \esp{debatir} est plus formel pour un débat scolaire). « Discuter » (bavarder) = \esp{charlar, conversar}.
\item \textbf{Le calque « il faut »} : Ne traduis jamais par \esp{es necesario que} + indicatif (déclenche subjonctif complexe) ni \esp{il faut que}. Utilise \textbf{\esp{hay que + infinitif}} (impersonnel, simple, niveau A2/B1) : \esp{Hay que respetar}, \esp{Hay que dialogar}. Ou \esp{deber + infinitif} si sujet précis : \esp{Los alumnos deben respetar}.
\item \textbf{Confondre \esp{ser} et \esp{estar}} : Qualité permanente, définition, identité $\rightarrow$ \esp{ser} : \esp{La tolerancia es importante}, \esp{El Camerún es diverso}. État passager, lieu, résultat d'action $\rightarrow$ \esp{estar} : \esp{Estoy de acuerdo}, \esp{El liceo está en Yaoundé}, \esp{La puerta está abierta}. \esp{La tolerancia está importante} = \textbf{FAUTE GRAVE}.
\item \textbf{Mal placer ou épeler les connecteurs} : \esp{sin embargo} (deux mots) et \esp{además} (un mot, accent) se placent \textbf{en tête de phrase}, suivis d'une \textbf{virgule}. \esp{Sin embargo, creo que...} $\checkmark$ \quad \esp{Creo que, sin embargo...} $\times$ (mal placé) \quad \esp{Sinembargo} $\times$ (orthographe).
\end{enumerate}
\end{attention}

\begin{savaistu}[Le Cameroun et la Guinée équatoriale : frères hispanophones ?]
Le Cameroun est membre observateur de l'Organisation des États ibéro-américains (OEI) et partage une frontière avec la \textbf{Guinée équatoriale}, seul pays africain dont l'espagnol est langue officielle (avec le français et le portugais). À \esp{Malabo} (capitale, sur l'île de Bioko) et \esp{Bata} (continent), l'espagnol est la langue de l'administration, de l'école et des médias. De nombreux Camerounais du Sud (région du Sud, Kribi, Campo) pratiquent l'espagnol pour le commerce frontalier. Apprendre l'espagnol en 1ère A, c'est aussi s'ouvrir à ce voisin immédiat et à la \esp{Comunidad Hispanohablante Africana}. \esp{¡África también habla español!}
\end{savaistu}

\begin{exercice}[Entraînement personnel : Debate y Redacción]
\begin{enumerate}
\item \textbf{Complète} les phrases avec le connecteur logique approprié (\esp{además, sin embargo, en primer lugar, por último, por eso, en conclusión}) :
\begin{enumerate}
\item \esp{El respeto es la base de la convivencia. \_\_\_\_\_\_, sin respeto no hay confianza.}
\item \esp{Los jóvenes tienen ideas nuevas. \_\_\_\_\_\_, los mayores tienen experiencia.}
\item \esp{Hay problemas en el barrio; \_\_\_\_\_\_, el diálogo puede resolverlos.}
\item \esp{\_\_\_\_\_\_, debemos trabajar juntos por la paz.}
\end{enumerate}
\item \textbf{Traduis en espagnol (Thème)} :
\begin{enumerate}
\item Dans mon lycée, il y a des élèves de plusieurs ethnies.
\item Nous devons dialoguer pour éviter les conflits.
\item La discrimination n'est pas acceptable.
\item Est-ce que tu es d'accord avec moi ? (utiliser \esp{¿...?})
\end{enumerate}
\item \textbf{Mini-rédaction (80 mots)} : \esp{« La tolerancia en mi clase de español ».} Utilise \esp{en primer lugar, además, sin embargo, en conclusión}. Parle de l'ambiance, des origines diverses, d'un conflit résolu.
\end{enumerate}
\end{exercice}

\begin{corrige}[Exercice EA12]
\textbf{1. Connecteurs :}
\begin{enumerate}
\item \esp{En primer lugar} (ordre / argument principal).
\item \esp{Además} (ajout) ou \esp{Por otro lado} (opposition douce non dans liste). Avec la liste imposée : \esp{Además} (ajout d'un 2ème sujet) ou \esp{Sin embargo} (contraste). \textbf{Meilleur :} \esp{Además} (ajout d'une seconde idée positive) ou \esp{Sin embargo} (contraste jeunes/âgés). Disons \esp{Sin embargo} pour marquer le contraste générationnel.
\item \esp{sin embargo} (opposition problème/solution).
\item \esp{En conclusión} ou \esp{Por último} (conclusion / appel final).
\end{enumerate}

\textbf{2. Thème :}
\begin{enumerate}
\item \esp{En mi liceo hay alumnos de varias etnias.} (ou \esp{grupos étnicos}).
\item \esp{Debemos dialogar para evitar los conflictos.} (\esp{evitar} + infinitif).
\item \esp{La discriminación no es aceptable.} (\esp{ser} + adjectif).
\item \esp{¿Estás de acuerdo conmigo?} (\esp{estar de acuerdo con} + pronom tonique \esp{migo}).
\end{enumerate}

\textbf{3. Mini-rédaction (exemple ~85 mots) :}
\esp{En mi clase de español hay alumnos de diferentes regiones de Camerún. \textbf{En primer lugar}, esta diversidad enriquece nuestros debates. \textbf{Además}, aprendemos expresiones de otras culturas. \textbf{Sin embargo}, a veces hay malentendidos por los acentos. \textbf{En conclusión}, la tolerancia nos permite aprender juntos en buen ambiente. \textbf{Hay que respetar} todas las voces.}
\end{corrige}

\newpage

\coursec{Exercices}

\courssub{Je vérifie mes connaissances}

\begin{exercice}[QCM : les connecteurs logiques]
Pour chaque phrase, choisis le connecteur qui convient (entoure la bonne réponse).
\begin{enumerate}
\item \esp{\_\_\_\_\_\_\_\_\_, debemos respetar a todas las personas.} \quad a) \esp{Sin embargo} \quad b) \esp{En primer lugar} \quad c) \esp{Por último}
\item \esp{La discriminación es un problema grave; \_\_\_\_\_\_\_\_\_, muchos jóvenes luchan contra ella.} \quad a) \esp{además} \quad b) \esp{sin embargo} \quad c) \esp{en primer lugar}
\item \esp{El diálogo es necesario. \_\_\_\_\_\_\_\_\_, la tolerancia hace la vida más fácil.} \quad a) \esp{Además} \quad b) \esp{Sin embargo} \quad c) \esp{Pero}
\item \esp{\_\_\_\_\_\_\_\_\_, quiero decir que todos somos iguales.} \quad a) \esp{En primer lugar} \quad b) \esp{Además} \quad c) \esp{Por último}
\end{enumerate}
\end{exercice}

\begin{exercice}[Vrai ou faux ? Justifie ta réponse]
Réponds par V (vrai) ou F (faux) et justifie chaque réponse par une phrase en français.
\begin{enumerate}
\item \esp{Sin embargo} exprime une addition d'idées.
\item \esp{En primer lugar} sert à introduire le premier argument.
\item \esp{Hay que} exprime une obligation générale.
\item \esp{Por último} annonce un exemple.
\end{enumerate}
\end{exercice}

\begin{exercice}[Vocabulaire : le débat et le vivre-ensemble]
Complète chaque phrase avec le mot espagnol qui convient, choisi dans la liste : \esp{la tolerancia -- el respeto -- la discriminación -- la diversidad -- el diálogo}.
\begin{enumerate}
\item \esp{El \_\_\_\_\_\_\_\_\_ entre las personas es la base de la convivencia.}
\item \esp{Rechazar a alguien por su color de piel es \_\_\_\_\_\_\_\_\_.}
\item \esp{La \_\_\_\_\_\_\_\_\_ cultural de Camerún es una riqueza.}
\item \esp{Para resolver los conflictos, necesitamos \_\_\_\_\_\_\_\_\_.}
\item \esp{La \_\_\_\_\_\_\_\_\_ significa aceptar las diferencias de los demás.}
\end{enumerate}
\end{exercice}

\begin{exercice}[Conjugaison à compléter]
Conjugue les verbes entre parenthèses au présent de l'indicatif.
\begin{enumerate}
\item \esp{Nosotros (tener) \_\_\_\_\_\_\_\_\_ que respetar las diferencias.}
\item \esp{Los jóvenes (deber) \_\_\_\_\_\_\_\_\_ dialogar con sus padres.}
\item \esp{Yo (poder) \_\_\_\_\_\_\_\_\_ aprender mucho de otras culturas.}
\item \esp{En mi barrio, todos (vivir) \_\_\_\_\_\_\_\_\_ en paz.}
\item \esp{Tú (querer) \_\_\_\_\_\_\_\_\_ participar en el debate, ¿verdad?}
\end{enumerate}
\end{exercice}

\courssub{Je m'entraîne}

\begin{exercice}[Thème : phrases sur l'identité et la tolérance]
Traduis en espagnol les phrases suivantes.
\begin{enumerate}
\item Il faut respecter les différences.
\item D'abord, chaque culture a ses coutumes.
\item De plus, la diversité est une richesse.
\item Cependant, les préjugés existent.
\item Enfin, nous devons dialoguer pour vivre en paix.
\end{enumerate}
\end{exercice}

\begin{exercice}[Texte à trous : les connecteurs]
Complète le texte suivant avec les connecteurs appropriés : \esp{además -- sin embargo -- por último -- en primer lugar -- por eso}.

\begin{textebox}[La convivencia en mi ciudad]
En mi ciudad, Douala, viven personas de muchas culturas. \_\_\_\_\_\_\_\_\_\_\_\_, creo que esta diversidad es muy positiva porque aprendemos unos de otros. \_\_\_\_\_\_\_\_\_\_\_\_, los mercados y las fiestas muestran la riqueza de nuestras tradiciones. \_\_\_\_\_\_\_\_\_\_\_\_, a veces hay tensiones y prejuicios entre los vecinos. \_\_\_\_\_\_\_\_\_\_\_\_, es necesario organizar debates en las escuelas. \_\_\_\_\_\_\_\_\_\_\_\_, pienso que el diálogo es la única solución para vivir juntos en paz.
\end{textebox}
\end{exercice}

\begin{exercice}[Appariements : mots et définitions]
Relie chaque mot espagnol (1--5) à sa définition française (a--e).
\begin{enumerate}
\item \esp{el prejuicio}
\item \esp{la convivencia}
\item \esp{las raíces}
\item \esp{el respeto}
\item \esp{la costumbre}
\end{enumerate}
\begin{itemize}
\item[a)] le fait de vivre ensemble en harmonie
\item[b)] une opinion négative formée sans connaître la réalité
\item[c)] une habitude transmise par une communauté
\item[d)] les origines culturelles d'une personne
\item[e)] la considération que l'on doit aux autres
\end{itemize}
\end{exercice}

\begin{exercice}[Reformulation : du dialogue au texte argumentatif]
Voici un court dialogue-débat entre deux élèves. Transforme-le en un court texte argumentatif de 5 à 6 phrases à la troisième personne, en utilisant au moins trois connecteurs logiques (\esp{en primer lugar, además, sin embargo, por último}).

\begin{textebox}[Un debate en el pasillo]
-- Ana: Yo creo que los jóvenes deben conservar las costumbres de sus raíces.\\
-- Luis: Sí, pero también deben estar abiertos a otras culturas.\\
-- Ana: Claro, la diversidad es una riqueza.\\
-- Luis: Entonces, hay que respetar las tradiciones y dialogar con los demás.
\end{textebox}
\end{exercice}

\courssub{Je me prépare au Bac}

\begin{exercice}[Compréhension et questions de langue]
Lis le texte suivant, puis réponds aux questions.

\begin{textebox}[Un debate en el liceo de Yaoundé]
El martes pasado, los alumnos de Première del liceo de Yaoundé organizaron un debate sobre la tolerancia y la convivencia. La profesora de español propuso una pregunta: «¿Debemos conservar nuestras tradiciones o abrirnos a otras culturas?»

En primer lugar, Amina explicó que las tradiciones son la memoria de un pueblo: sin ellas, perdemos nuestra identidad. Además, contó el ejemplo de su abuela, que le enseñó los cuentos y las canciones de su región. Sin embargo, Etienne defendió otra idea: para él, los jóvenes deben conocer otras culturas porque el mundo cambia rápidamente. Dio el ejemplo de sus amigos de Douala y de Bata, en Guinea Ecuatorial, que comparten con él su música y su comida. Por último, la clase votó y decidió que las dos ideas son compatibles: hay que respetar las raíces y, al mismo tiempo, aceptar las diferencias. El debate terminó con una frase escrita en la pizarra: «La diversidad es una riqueza, el diálogo es el camino».
\end{textebox}

\textbf{A. Questions de compréhension} (réponds en espagnol, avec des phrases complètes) :
\begin{enumerate}
\item ¿Cuál es el tema del debate del liceo?
\item ¿Qué argumento da Amina para defender las tradiciones?
\item ¿Por qué quiere Etienne conocer otras culturas?
\item ¿Qué decidió la clase al final del debate?
\end{enumerate}

\textbf{B. Questions de langue} :
\begin{enumerate}
\item Relève dans le texte les quatre connecteurs logiques et traduis-les en français.
\item Trouve dans le texte un synonyme de \esp{costumbres} et un contraire de \esp{la paz}.
\item Conjugue \esp{perder} et \esp{defender} au présent de l'indicatif (toutes les personnes).
\item Traduis en français la dernière phrase du texte.
\end{enumerate}
\end{exercice}

\begin{exercice}[Thème et version]
\textbf{A. Thème.} Traduis en espagnol le texte suivant (attention aux connecteurs et à \esp{hay que}) :

« Il faut respecter les différences. D'abord, chaque culture a ses coutumes. De plus, la diversité est une richesse. Cependant, les préjugés existent. Enfin, nous devons dialoguer pour vivre en paix. »

\textbf{B. Version.} Traduis en français le texte suivant :

\begin{textebox}[La tolerancia]
La tolerancia es necesaria en nuestra sociedad. En primer lugar, hay que aceptar que todas las personas son diferentes. Además, cada cultura tiene sus tradiciones y su historia. Sin embargo, todavía existen la discriminación y los prejuicios en muchos países. Por último, creo que el diálogo y la educación son las mejores armas para construir la paz.
\end{textebox}
\end{exercice}

\begin{exercice}[Expression écrite et orale type Bac]
\textbf{A. Expression écrite.} Rédige un texte argumentatif en espagnol d'\textbf{environ 100 mots} sur le sujet suivant :

\esp{¿Es importante respetar las culturas de los demás? Da tu opinión con al menos tres argumentos.}

Consignes obligatoires :
\begin{itemize}
\item utilise au moins quatre connecteurs logiques : \esp{en primer lugar, además, sin embargo, por último} ;
\item utilise au moins une fois \esp{hay que} ou \esp{debemos} ;
\item soigne l'orthographe et les accents.
\end{itemize}

\textbf{B. Expression orale.} Par binômes, préparez un débat de \textbf{4 minutes} sur le sujet : \esp{¿Deben los jóvenes conservar las costumbres de sus raíces?} Un élève défend le « oui », l'autre le « non ». Chaque élève doit utiliser au moins trois connecteurs logiques et deux mots du lexique du vivre-ensemble. Le reste de la classe remplit une fiche d'évaluation : connecteurs utilisés, arguments donnés, prononciation, respect du temps de parole.
\end{exercice}

\newpage

\coursec{Corrigés des exercices}

\begin{corrige}[Exercice 1 — QCM : les connecteurs logiques]
\begin{enumerate}
\item Réponse b) \esp{En primer lugar} : \esp{En primer lugar, debemos respetar a todas las personas.} « D'abord, nous devons respecter toutes les personnes. » Le connecteur introduit le premier argument d'une liste.
\item Réponse b) \esp{sin embargo} : \esp{La discriminación es un problema grave; sin embargo, muchos jóvenes luchan contra ella.} « La discrimination est un problème grave ; cependant, beaucoup de jeunes luttent contre elle. » Il y a opposition entre le problème et la lutte.
\item Réponse a) \esp{Además} : \esp{El diálogo es necesario. Además, la tolerancia hace la vida más fácil.} « Le dialogue est nécessaire. De plus, la tolérance rend la vie plus facile. » On ajoute un deuxième argument dans le même sens.
\item Réponse c) \esp{Por último} : \esp{Por último, quiero decir que todos somos iguales.} « Enfin, je veux dire que nous sommes tous égaux. » Le connecteur annonce la conclusion.
\end{enumerate}
\textbf{Point de vigilance :} après \esp{sin embargo}, \esp{además}, \esp{por último}, on met une virgule ; \esp{sin embargo} et \esp{además} s'écrivent chacun en un seul mot (ne pas écrire \esp{sin en bargo} ni \esp{a demás}).
\end{corrige}

\begin{corrige}[Exercice 2 — Vrai ou faux ?]
\begin{enumerate}
\item \textbf{F.} \esp{Sin embargo} exprime l'\textbf{opposition} (« cependant »), pas l'addition ; l'addition s'exprime avec \esp{además}.
\item \textbf{V.} \esp{En primer lugar} (« d'abord, en premier lieu ») introduit le premier argument d'un texte argumentatif.
\item \textbf{V.} \esp{Hay que} + infinitif exprime une obligation générale et impersonnelle : \esp{Hay que respetar a los demás.} « Il faut respecter les autres. »
\item \textbf{F.} \esp{Por último} (« enfin ») annonce le dernier argument ou la conclusion ; pour introduire un exemple, on dit \esp{por ejemplo}.
\end{enumerate}
\textbf{Point de vigilance :} \esp{hay que} est invariable et impersonnel : on ne conjugue jamais \esp{haber} ici (\esp{tenemos que} = « nous devons », autre structure, personnelle).
\end{corrige}

\begin{corrige}[Exercice 3 — Vocabulaire : le débat et le vivre-ensemble]
\begin{enumerate}
\item \esp{El \textbf{respeto} entre las personas es la base de la convivencia.} « Le respect entre les personnes est la base du vivre-ensemble. »
\item \esp{Rechazar a alguien por su color de piel es \textbf{discriminación}.} « Rejeter quelqu'un à cause de sa couleur de peau, c'est de la discrimination. »
\item \esp{La \textbf{diversidad} cultural de Camerún es una riqueza.} « La diversité culturelle du Cameroun est une richesse. »
\item \esp{Para resolver los conflictos, necesitamos \textbf{diálogo}.} « Pour résoudre les conflits, nous avons besoin de dialogue. »
\item \esp{La \textbf{tolerancia} significa aceptar las diferencias de los demás.} « La tolérance signifie accepter les différences des autres. »
\end{enumerate}
\textbf{Point de vigilance :} \esp{respeto} s'écrit sans accent (faux ami orthographique avec le français « respect ») ; \esp{diálogo} prend un accent sur le a ; \esp{los demás} (« les autres ») prend un accent sur le a.
\end{corrige}

\begin{corrige}[Exercice 4 — Conjugaison à compléter]
\begin{enumerate}
\item \esp{Nosotros \textbf{tenemos} que respetar las diferencias.} (tener, 1\iere{} pers. plur. ; \esp{tener que} + infinitif = « devoir » ; pas de diphtongaison à \esp{nosotros})
\item \esp{Los jóvenes \textbf{deben} dialogar con sus padres.} (deber, 3\ieme{} pers. plur., régulier en -er)
\item \esp{Yo \textbf{puedo} aprender mucho de otras culturas.} (poder, verbe à diphtongaison o $\rightarrow$ ue)
\item \esp{En mi barrio, todos \textbf{viven} en paz.} (vivir, 3\ieme{} pers. plur., régulier en -ir)
\item \esp{Tú \textbf{quieres} participar en el debate, ¿verdad?} (querer, diphtongaison e $\rightarrow$ ie)
\end{enumerate}
\textbf{Point de vigilance :} ne pas écrire \esp{podo} ni \esp{quero} : les diphtongues \esp{puedo}, \esp{quieres} sont obligatoires au présent ; mais la diphtongaison disparaît à \esp{nosotros} et \esp{vosotros} : \esp{podemos}, \esp{queremos}.
\end{corrige}

\begin{corrige}[Exercice 5 — Thème : phrases sur l'identité et la tolérance]
\begin{enumerate}
\item \esp{Hay que respetar las diferencias.} (obligation générale impersonnelle : \esp{hay que} + infinitif)
\item \esp{En primer lugar, cada cultura tiene sus costumbres.} (« d'abord » = \esp{en primer lugar} ; \esp{cada} est invariable)
\item \esp{Además, la diversidad es una riqueza.} (attention : \esp{además} en un seul mot, avec accent sur le a final)
\item \esp{Sin embargo, los prejuicios existen.} (« cependant » = \esp{sin embargo} ; \esp{prejuicio} : préjugé)
\item \esp{Por último, debemos dialogar para vivir en paz.} (« pour » + infinitif de but = \esp{para} + infinitif)
\end{enumerate}
\textbf{Points de vigilance :} ne pas traduire « il faut » par \esp{es necesario que} + subjonctif ici (trop lourd à ce niveau) ; \esp{para} et non \esp{por} devant l'infinitif de but ; \esp{vivir en paz} sans article.
\end{corrige}

\begin{corrige}[Exercice 6 — Texte à trous : les connecteurs]
Texte complété :

\esp{En mi ciudad, Douala, viven personas de muchas culturas. \textbf{En primer lugar}, creo que esta diversidad es muy positiva porque aprendemos unos de otros. \textbf{Además}, los mercados y las fiestas muestran la riqueza de nuestras tradiciones. \textbf{Sin embargo}, a veces hay tensiones y prejuicios entre los vecinos. \textbf{Por eso}, es necesario organizar debates en las escuelas. \textbf{Por último}, pienso que el diálogo es la única solución para vivir juntos en paz.}

Traduction : « Dans ma ville, Douala, vivent des personnes de nombreuses cultures. D'abord, je crois que cette diversité est très positive parce que nous apprenons les uns des autres. De plus, les marchés et les fêtes montrent la richesse de nos traditions. Cependant, il y a parfois des tensions et des préjugés entre les voisins. C'est pourquoi il est nécessaire d'organiser des débats dans les écoles. Enfin, je pense que le dialogue est la seule solution pour vivre ensemble en paix. »

\textbf{Justification :} \esp{en primer lugar} ouvre l'argumentation ; \esp{además} ajoute un argument dans le même sens ; \esp{sin embargo} introduit le constat négatif opposé ; \esp{por eso} (« c'est pourquoi ») exprime la conséquence ; \esp{por último} conclut. Remarque : \esp{creo que} et \esp{pienso que} sont suivis de l'indicatif car ils expriment une opinion affirmée.
\end{corrige}

\begin{corrige}[Exercice 7 — Appariements : mots et définitions]
\begin{enumerate}
\item \esp{el prejuicio} $\rightarrow$ b) une opinion négative formée sans connaître la réalité
\item \esp{la convivencia} $\rightarrow$ a) le fait de vivre ensemble en harmonie
\item \esp{las raíces} $\rightarrow$ d) les origines culturelles d'une personne (littéralement « les racines »)
\item \esp{el respeto} $\rightarrow$ e) la considération que l'on doit aux autres
\item \esp{la costumbre} $\rightarrow$ c) une habitude transmise par une communauté
\end{enumerate}
\textbf{Point de vigilance :} \esp{las raíces} prend un accent (\esp{raíz} au singulier, pluriel \esp{raíces}) ; \esp{convivencia} vient du verbe \esp{convivir} (« vivre ensemble ») ; \esp{la costumbre} est féminin.
\end{corrige}

\begin{corrige}[Exercice 8 — Reformulation : du dialogue au texte argumentatif]
\textbf{Proposition de texte modèle :}

\esp{En el debate, Ana y Luis hablan de las tradiciones y de la apertura a otras culturas. En primer lugar, Ana piensa que los jóvenes deben conservar las costumbres de sus raíces. Sin embargo, Luis opina que también deben estar abiertos a otras culturas. Además, los dos están de acuerdo en que la diversidad es una riqueza. Por último, concluyen que hay que respetar las tradiciones y dialogar con los demás.}

Traduction : « Dans le débat, Ana et Luis parlent des traditions et de l'ouverture aux autres cultures. D'abord, Ana pense que les jeunes doivent conserver les coutumes de leurs racines. Cependant, Luis est d'avis qu'ils doivent aussi être ouverts aux autres cultures. De plus, tous les deux sont d'accord pour dire que la diversité est une richesse. Enfin, ils concluent qu'il faut respecter les traditions et dialoguer avec les autres. »

\textbf{Points de vigilance :} passer à la troisième personne (\esp{piensa}, \esp{opina}, \esp{concluyen}) ; utiliser des verbes de parole et d'opinion (\esp{decir que, pensar que, opinar que, estar de acuerdo}) ; conserver au moins trois connecteurs ; \esp{estar abiertos} avec \esp{estar} (état) et accord de l'adjectif au pluriel.
\end{corrige}

\begin{corrige}[Exercice 9 — Compréhension et questions de langue]
\textbf{A. Questions de compréhension}
\begin{enumerate}
\item \esp{El tema del debate es la tolerancia y la convivencia: ¿debemos conservar nuestras tradiciones o abrirnos a otras culturas?} « Le sujet du débat est la tolérance et le vivre-ensemble : devons-nous conserver nos traditions ou nous ouvrir aux autres cultures ? »
\item \esp{Amina dice que las tradiciones son la memoria de un pueblo y que sin ellas perdemos nuestra identidad.} « Amina dit que les traditions sont la mémoire d'un peuple et que sans elles nous perdons notre identité. »
\item \esp{Etienne quiere conocer otras culturas porque el mundo cambia rápidamente.} « Etienne veut connaître d'autres cultures parce que le monde change rapidement. »
\item \esp{Al final, la clase decidió que las dos ideas son compatibles: hay que respetar las raíces y, al mismo tiempo, aceptar las diferencias.} « Finalement, la classe a décidé que les deux idées sont compatibles : il faut respecter les racines et, en même temps, accepter les différences. »
\end{enumerate}

\textbf{B. Questions de langue}
\begin{enumerate}
\item Les quatre connecteurs : \esp{en primer lugar} « d'abord » ; \esp{además} « de plus » ; \esp{sin embargo} « cependant » ; \esp{por último} « enfin ».
\item Synonyme de \esp{costumbres} : \esp{tradiciones}. Contraire de \esp{la paz} : on peut relever \esp{la discriminación} ou proposer \esp{la guerra} / \esp{el conflicto} (le texte oppose la paix aux tensions et aux préjugés).
\item \esp{perder} (e $\rightarrow$ ie) : \esp{pierdo, pierdes, pierde, perdemos, perdéis, pierden}. \esp{defender} (e $\rightarrow$ ie) : \esp{defiendo, defiendes, defiende, defendemos, defendéis, defienden}. Rappel : pas de diphtongaison à \esp{nosotros} ni \esp{vosotros}.
\item \esp{«La diversidad es una riqueza, el diálogo es el camino».} « La diversité est une richesse, le dialogue est le chemin. »
\end{enumerate}
\textbf{Barème indicatif (sur 20) :} compréhension 8 pts (2 pts par question), questions de langue 12 pts (4 pts connecteurs, 2 pts synonyme/contraire, 4 pts conjugaisons, 2 pts traduction).
\end{corrige}

\begin{corrige}[Exercice 10 — Thème et version]
\textbf{A. Thème.} Traduction proposée :

\esp{Hay que respetar las diferencias. En primer lugar, cada cultura tiene sus costumbres. Además, la diversidad es una riqueza. Sin embargo, los prejuicios existen. Por último, debemos dialogar para vivir en paz.}

Points de vigilance : \esp{hay que} pour « il faut » ; virgule après chaque connecteur ; \esp{para} + infinitif pour exprimer le but ; accents sur \esp{además} et \esp{diversidad}.

\textbf{B. Version.} Traduction proposée :

« La tolérance est nécessaire dans notre société. D'abord, il faut accepter que toutes les personnes sont différentes. De plus, chaque culture a ses traditions et son histoire. Cependant, la discrimination et les préjugés existent encore dans de nombreux pays. Enfin, je crois que le dialogue et l'éducation sont les meilleures armes pour construire la paix. »

Points de vigilance : \esp{todavía} = « encore » (et non « déjà », qui se dit \esp{ya}) ; \esp{las armas} est féminin malgré l'article \esp{el} au singulier (\esp{el arma} mais \esp{las armas}) ; \esp{construir la paz} = « construire la paix » ; \esp{la educación} = « l'éducation » (faux ami partiel : « l'éducation » scolaire et familiale, pas seulement la politesse).

\textbf{Barème indicatif (sur 20) :} thème 10 pts (2 pts par phrase : connecteur, grammaire, orthographe), version 10 pts (fidélité 6 pts, qualité du français 4 pts).
\end{corrige}

\begin{corrige}[Exercice 11 — Expression écrite et orale type Bac]
\textbf{A. Expression écrite.} Proposition de rédaction modèle (environ 100 mots) :

\esp{¿Es importante respetar las culturas de los demás? En mi opinión, sí, es fundamental. En primer lugar, todas las personas son iguales y hay que aceptar las diferencias de religión, de lengua y de costumbres. Además, la diversidad cultural es una riqueza: gracias a ella, aprendemos otras músicas, otras comidas y otras ideas. Sin embargo, todavía existen los prejuicios y la discriminación en muchos países, y eso crea conflictos. Por eso, debemos dialogar y conocer mejor a los demás. Por último, creo que el respeto es la base de la convivencia: sin respeto, no hay paz. En conclusión, respetar las culturas de los demás es construir un mundo mejor.}

Traduction : « Est-il important de respecter les cultures des autres ? À mon avis, oui, c'est fondamental. D'abord, toutes les personnes sont égales et il faut accepter les différences de religion, de langue et de coutumes. De plus, la diversité culturelle est une richesse : grâce à elle, nous découvrons d'autres musiques, d'autres plats et d'autres idées. Cependant, les préjugés et la discrimination existent encore dans de nombreux pays, et cela crée des conflits. C'est pourquoi nous devons dialoguer et mieux connaître les autres. Enfin, je crois que le respect est la base du vivre-ensemble : sans respect, pas de paix. En conclusion, respecter les cultures des autres, c'est construire un monde meilleur. »

\textbf{Barème indicatif (sur 20) :} respect du sujet et de la longueur 4 pts ; connecteurs logiques (au moins 4) 4 pts ; emploi de \esp{hay que} / \esp{debemos} 2 pts ; correction grammaticale 5 pts ; orthographe et accents 3 pts ; richesse du vocabulaire 2 pts.

\textbf{Points de vigilance :} \esp{en mi opinión} avec accent (et non \esp{en mi opinion}) ; \esp{los demás} avec accent ; pas de majuscule après la virgule dans \esp{sin respeto, no hay paz} ; éviter le calque \esp{es importante que respeta} : après \esp{es importante que}, il faudrait le subjonctif (\esp{es importante que respetemos}) — à ce niveau, préférer la tournure avec infinitif (\esp{es importante respetar}) ; \esp{gracias a} + nom (« grâce à »), ne pas confondre avec \esp{por eso} (« c'est pourquoi »).

\textbf{B. Expression orale.} Grille d'évaluation proposée pour la classe (sur 10) : connecteurs logiques utilisés (au moins 3) 3 pts ; arguments clairs et pertinents 3 pts ; lexique du vivre-ensemble (au moins 2 mots : \esp{tolerancia, respeto, diversidad, raíces, convivencia...}) 2 pts ; prononciation et fluidité 1 pt ; respect du temps de parole 1 pt.

Exemples de phrases attendues : pour le « oui » : \esp{En primer lugar, las costumbres son nuestra identidad. Además, nuestras raíces nos dan fuerza.} « D'abord, les coutumes sont notre identité. De plus, nos racines nous donnent de la force. » Pour le « non » : \esp{Sin embargo, el mundo cambia y debemos abrirnos a otras culturas. Por último, la diversidad es una riqueza.} « Cependant, le monde change et nous devons nous ouvrir aux autres cultures. Enfin, la diversité est une richesse. »

\textbf{Conseil de prononciation :} \esp{convivencia} se prononce « kon-bi-bén-sia » ; \esp{prejuicio} « pré-houi-sio » ; le h de \esp{hay que} est muet : « aï ké ».
\end{corrige}

\newpage

\coursec{Fiche bilan}

\begin{popfigure}
\begin{tikzpicture}[mindmap, grow cyclic, every node/.style={concept, font=\small},
root concept/.append style={concept color=popblue, font=\small\bfseries},
level 1/.append style={level distance=3.4cm, sibling angle=60},
level 2/.append style={level distance=2.2cm}]
\node[root concept]{Identidad, tolerancia y convivencia}
  child[concept color=poporange]{ node {Léxico del debate}
    child { node {la tolerancia, el respeto} }
    child { node {la discriminación, el prejuicio} }
  }
  child[concept color=popgreen]{ node {Conectores}
    child { node {en primer lugar, además} }
    child { node {sin embargo, por último} }
  }
  child[concept color=poppurple]{ node {Opinar}
    child { node {pienso que, me parece que} }
    child { node {estoy de acuerdo / no estoy de acuerdo} }
  }
  child[concept color=poppink]{ node {Traducción}
    child { node {identidad = identidad} }
    child { node {convivencia = vivre-ensemble} }
  }
  child[concept color=popgold]{ node {Método del debate}
    child { node {escuchar, argumentar, respetar} }
  }
  child[concept color=popblue!60]{ node {Texto argumentado (100 palabras)}
    child { node {introducción, argumentos, conclusión} }
  };
\end{tikzpicture}
\legende{Carte mentale de la leçon : identité, tolérance et vivre-ensemble}
\end{popfigure}

\begin{bilan}
\textbf{Lexique clé du débat et du vivre-ensemble}
\begin{itemize}
  \item \esp{la tolerancia} : la tolérance ; \esp{el respeto} : le respect ; \esp{la convivencia} : le vivre-ensemble.
  \item \esp{la discriminación} : la discrimination ; \esp{el prejuicio} : le préjugé ; \esp{la identidad} : l'identité.
  \item \esp{la diversidad cultural} : la diversité culturelle ; \esp{la igualdad} : l'égalité ; \esp{la libertad de expresión} : la liberté d'expression.
  \item \esp{aceptar a los demás} : accepter les autres ; \esp{respetar las diferencias} : respecter les différences.
  \item \esp{opinar / pensar} : donner son avis / penser ; \esp{estar de acuerdo (con)} : être d'accord (avec) ; \esp{no estar de acuerdo} : ne pas être d'accord.
\end{itemize}

\textbf{Connecteurs logiques à connaître par cœur}
\begin{center}
\begin{tabularx}{\linewidth}{|X|X|X|}\hline
\rowcolor{popblueL} Español & Français & Fonction \\ \hline
en primer lugar & en premier lieu & introduire le 1er argument \\ \hline
además & de plus & ajouter un argument \\ \hline
sin embargo & cependant & exprimer une opposition \\ \hline
por un lado / por otro lado & d'un côté / de l'autre & présenter deux points de vue \\ \hline
por último & enfin & dernier argument \\ \hline
en conclusión & en conclusion & conclure \\ \hline
\end{tabularx}
\end{center}

\textbf{Grammaire et expressions à retenir}
\begin{itemize}
  \item Après \esp{pienso que / creo que / me parece que} : indicatif (\esp{Pienso que la tolerancia es importante.} « Je pense que la tolérance est importante. »).
  \item Après \esp{no pienso que / no creo que} : subjonctif (\esp{No creo que sea justo.} « Je ne crois pas que ce soit juste. »).
  \item \esp{hay que} + infinitif : il faut (\esp{Hay que respetar a los demás.} « Il faut respecter les autres. »).
  \item \esp{es importante / es necesario} + infinitif : il est important / nécessaire de.
\end{itemize}

\textbf{Méthodes express}
\begin{itemize}
  \item \textbf{Débat oral} : écouter l'autre, donner son avis avec \esp{pienso que...}, argumenter avec \esp{porque...}, réagir avec \esp{estoy de acuerdo / sin embargo...}, respecter la parole de chacun.
  \item \textbf{Texte argumenté (100 mots)} : 1) une phrase d'introduction qui pose le sujet ; 2) deux ou trois arguments reliés par les connecteurs ; 3) une conclusion personnelle avec \esp{en conclusión}.
  \item \textbf{Traduction} : repérer le connecteur français et choisir son équivalent exact ; attention aux faux amis.
\end{itemize}
\end{bilan}

\begin{aretenir}[Checklist avant le Bac]
\begin{itemize}
  \item $\square$ Je connais le lexique du débat : \esp{tolerancia, respeto, discriminación, prejuicio, convivencia}.
  \item $\square$ Je sais utiliser \esp{en primer lugar, además, sin embargo, por último, en conclusión}.
  \item $\square$ Je sais donner mon avis : \esp{pienso que..., me parece que..., en mi opinión...}.
  \item $\square$ Je sais exprimer l'accord et le désaccord : \esp{estoy de acuerdo / no estoy de acuerdo}.
  \item $\square$ Je sais que \esp{no creo que} est suivi du subjonctif.
  \item $\square$ Je sais construire \esp{hay que} + infinitif pour exprimer l'obligation générale.
  \item $\square$ Je sais structurer un texte argumenté de 100 mots : introduction, arguments, conclusion.
  \item $\square$ Je sais traduire \esp{la convivencia} par « le vivre-ensemble » et non par un calque.
  \item $\square$ Je n'oublie pas les accents : \esp{además, opinión, discriminación, conclusión}.
  \item $\square$ Je sais écouter et respecter la parole des autres dans un débat.
  \item $\square$ Je relis ma production : connecteurs, accords, ponctuation ¿ ? ¡ !
\end{itemize}
\end{aretenir}

\findecours

\end{document}
